% Production et traduction de textes liés à l'environnement et à la santé — Allemand 1ère A — Cours signé OnBuch+
% Programme officiel MINESEC. Compiler avec : tectonic ALA14-production-et-traduction-de-textes-lies-a-l-environnement-et.tex
\newcommand{\DOCMATIERE}{Allemand}
\newcommand{\DOCNIVEAU}{1ère A}
\newcommand{\DOCMODULE}{Chapitre 3 : Environnement, santé et bien-être}
\newcommand{\DOCLECON}{ALA14}
\newcommand{\DOCTITRE}{Production et traduction de textes liés à l'environnement et à la santé}
% OnBuch+ — préambule « cours signés » ENRICHI (compilé avec tectonic / XeLaTeX)
% Système visuel « Pop » d'OnBuch+ : fond crème (jamais blanc), bordures encre
% épaisses, ombres « dures » décalées, titres Archivo Black en capitales.
% Version MATHÉMATIQUES : graphes (pgfplots/tikz), boîtes pédagogiques
% (définition, propriété, méthode, attention, le savais-tu, exercice résolu,
% exercice, TP, bilan), page de garde et signature OnBuch+ ancrée en bas.
%
% Macros attendues dans main.tex AVANT \input{preamble} :
%   \DOCMATIERE  \DOCNIVEAU  \DOCMODULE  \DOCLECON (numéro)  \DOCTITRE
\documentclass[11pt]{article}

\usepackage{fontspec}
\usepackage[ngerman,french]{babel} % français = langue principale ; l'allemand via \all{..} / textebox
\usepackage[a4paper,margin=2cm,top=3.4cm,bottom=2.8cm,headheight=1.4cm,footskip=1.3cm]{geometry}
\usepackage{amsmath,amssymb,amsfonts}
\usepackage{mathtools} % \xmapsto, \coloneqq... souvent utilisés par les modèles
\usepackage{cancel} % simplifications barrées (bilans de forces)
% \abs{z} (module, valeur absolue) et \norm{u} (norme), utilisés par les modèles
\providecommand{\abs}{}\let\abs\relax\DeclarePairedDelimiter{\abs}{\lvert}{\rvert}
\providecommand{\norm}{}\let\norm\relax\DeclarePairedDelimiter{\norm}{\lVert}{\rVert}
\usepackage[table]{xcolor} % [table] : fournit aussi \rowcolors (alternance de lignes)
\usepackage{pagecolor}
\usepackage{tikz}
\usetikzlibrary{arrows.meta,positioning,calc,shapes.geometric,shapes.misc,shapes.arrows,decorations.pathmorphing,decorations.pathreplacing,decorations.markings,patterns,shadows,mindmap,trees,fit,backgrounds,matrix,intersections,angles,quotes}
\usepackage{pgfplots}
\usepackage[version=4]{mhchem} % équations nucléaires \ce{^{235}_{92}U}
\usepackage[european,siunitx]{circuitikz} % schémas électriques
\pgfplotsset{compat=1.18}
\usepgfplotslibrary{fillbetween,groupplots} % aires entre courbes (calcul intégral)
\usepackage{enumitem}
\usepackage{fancyhdr}
% [most] charge toutes les bibliothèques tcolorbox (listings, minted, raster,
% documentation...) et gonfle inutilement la mémoire TeX sur un document long
% (~300 boîtes) : on ne charge que ce qui est réellement utilisé.
\usepackage[skins,breakable]{tcolorbox}
\usepackage{graphicx}
\usepackage{colortbl}
\usepackage{booktabs}
\usepackage{tabularx}
\usepackage{diagbox} % case d'en-tête diagonale des tableaux à double entrée
% Colonnes extensibles centrée / gauche / droite (C, L, R), souvent utilisées par les modèles
\newcolumntype{C}{>{\centering\arraybackslash}X}
\newcolumntype{L}{>{\raggedright\arraybackslash}X}
\newcolumntype{R}{>{\raggedleft\arraybackslash}X}
\usepackage{array}
\usepackage{multicol}
\usepackage{multirow}
\usepackage{siunitx}
% parse-numbers=false : accepte aussi \SI{2,5 \times 10^{-3}}{...}
\sisetup{locale=FR,output-decimal-marker={,},group-minimum-digits=4,parse-numbers=false}
\DeclareSIUnit\ohm{\ensuremath{\symup{\Omega}}} % U+2126 absent de la police
\DeclareSIUnit\division{div} % réglages d'oscilloscope (ms/div, V/div)
\DeclareSIUnit\tr{tr} % tours (vitesse de rotation, ex. \SI{1500}{\tr\per\minute})
\DeclareSIUnit\molar{M} % concentration molaire (ex. \SI{3}{\molar}, \SI{1}{\micro\molar})
\DeclareSIUnit\an{an} % année (le modèle confond parfois avec \year, un compteur TeX) — ex. \SI{5}{\kilo\meter\per\an}
\DeclareSIUnit\FCFA{FCFA} % franc CFA (ex. \SI{500}{\FCFA\per\kilo\gram})
\DeclareSIUnit\degreeBrix{\textdegree Brix} % teneur en sucre d'un sirop/jus (réfractométrie)
\DeclareSIUnit\month{mois} % le modèle utilise parfois \month, un compteur TeX, comme unité de durée
\usepackage{qrcode}
% \degree : commande fréquemment utilisée par les modèles pour un angle ou une
% température en dehors de \SI{}{} (non fournie par les paquets chargés).
\providecommand{\degree}{^{\circ}}
% \qed (fin de démonstration), utilisé par les modèles sans amsthm
\providecommand{\qed}{\hfill$\square$}
% \barre : double barre des valeurs interdites dans un tableau de variations (tkz-tab)
\providecommand{\barre}{\ensuremath{\|}}
% environnement proof (amsthm non chargé) utilisé par les modèles
\ifdefined\proof\else\newenvironment{proof}[1][Démonstration]{\par\noindent\textit{#1.}\ }{\qed\par}\fi
\usepackage{lastpage}
\usepackage{hyperref}
\hypersetup{hidelinks}

% ============================================================
% Palette « Pop » OnBuch+ — mêmes hex que la classe P (plus_ui.dart)
% ============================================================
\definecolor{poporange}{HTML}{F59321}
\definecolor{popdark}{HTML}{E07A0C}
\definecolor{popcream}{HTML}{FFF6E8}
\definecolor{popink}{HTML}{1C1714}
\definecolor{poppurple}{HTML}{7A5AE0}
\definecolor{popgreen}{HTML}{2E9E6A}
\definecolor{poppink}{HTML}{E0457B}
\definecolor{popblue}{HTML}{2F7DE1}
\definecolor{popgold}{HTML}{C98A12}
\definecolor{popmuted}{HTML}{8A7F76}
\definecolor{popline}{HTML}{EADFD2}
\definecolor{popgray}{HTML}{8A7F76}
\colorlet{popgrey}{popgray}
% Alias tolérants (le modèle nomme parfois les couleurs différemment)
\colorlet{popwhite}{white}
\colorlet{popblack}{popink}
\colorlet{popred}{poppink}
\colorlet{popyellow}{popgold}
% Teintes claires (fonds de tableaux/encadrés) — le modèle nomme parfois
% les couleurs avec un suffixe L (ex. popblueL) sans les avoir définies.
\colorlet{poporangeL}{poporange!20}
\colorlet{popdarkL}{popdark!20}
\colorlet{poppurpleL}{poppurple!20}
\colorlet{popgreenL}{popgreen!20}
\colorlet{poppinkL}{poppink!20}
\colorlet{popblueL}{popblue!20}
\colorlet{popgoldL}{popgold!20}
\colorlet{popredL}{popred!20}
\colorlet{popgrayL}{popgray!20}
% Teintes claires pour fonds de tableaux / zones
\colorlet{popblueL}{popblue!10!white}
\colorlet{poporangeL}{poporange!14!white}
\colorlet{popgreenL}{popgreen!12!white}
\colorlet{poppurpleL}{poppurple!12!white}
\colorlet{poppinkL}{poppink!12!white}
\colorlet{popgoldL}{popgold!14!white}

\pagecolor{popcream}

% ============================================================
% Typographie
% ============================================================
\defaultfontfeatures{Ligatures=TeX}
\setmainfont{PlusJakartaSans-Medium.ttf}[Path=../../../fonts/,BoldFont=PlusJakartaSans-Bold.ttf]
% Maths en sans-serif (Fira Math) avec chiffres et lettres droites de Plus
% Jakarta, pour que les formules se fondent dans le texte.
\usepackage[math-style=ISO,bold-style=ISO]{unicode-math}
\setmathfont{FiraMath-Regular.otf}[Path=../../../fonts/]
\setmathfont{PlusJakartaSans-Medium.ttf}[Path=../../../fonts/,range={up/num,up/{Latin,latin},\mathup}]
\usepackage{icomma} % virgule décimale sans espace en mode math
% Caractères mathématiques Unicode absents des polices de texte (Plus Jakarta,
% Archivo Black) : rendus par leur équivalent mathématique, en texte comme en
% formule — sinon ils s'affichent en carré vide (ex. « Arithmétique dans ℤ »).
\usepackage{newunicodechar}
\newunicodechar{ℝ}{\ensuremath{\mathbb{R}}}
\newunicodechar{ℤ}{\ensuremath{\mathbb{Z}}}
\newunicodechar{ℂ}{\ensuremath{\mathbb{C}}}
\newunicodechar{ℕ}{\ensuremath{\mathbb{N}}}
\newunicodechar{ℚ}{\ensuremath{\mathbb{Q}}}
\newunicodechar{𝔻}{\ensuremath{\mathbb{D}}}
\newunicodechar{α}{\ensuremath{\alpha}}
\newunicodechar{β}{\ensuremath{\beta}}
\newunicodechar{γ}{\ensuremath{\gamma}}
\newunicodechar{δ}{\ensuremath{\delta}}
\newunicodechar{ε}{\ensuremath{\varepsilon}}
\newunicodechar{θ}{\ensuremath{\theta}}
\newunicodechar{λ}{\ensuremath{\lambda}}
\newunicodechar{μ}{\ensuremath{\mu}}
\newunicodechar{σ}{\ensuremath{\sigma}}
\newunicodechar{τ}{\ensuremath{\tau}}
\newunicodechar{φ}{\ensuremath{\varphi}}
\newunicodechar{ψ}{\ensuremath{\psi}}
\newunicodechar{ω}{\ensuremath{\omega}}
\newunicodechar{Σ}{\ensuremath{\Sigma}}
% majuscules produites par \MakeUppercase dans les titres (α-aminés -> Α)
\newunicodechar{Α}{\ensuremath{\alpha}}
\newunicodechar{Β}{\ensuremath{\beta}}
\newunicodechar{Γ}{\ensuremath{\Gamma}}
\newunicodechar{Δ}{\ensuremath{\Delta}}
\newunicodechar{Ε}{\ensuremath{\varepsilon}}
\newunicodechar{Θ}{\ensuremath{\Theta}}
\newunicodechar{Λ}{\ensuremath{\Lambda}}
\newunicodechar{Μ}{\ensuremath{\mu}}
\newunicodechar{Τ}{\ensuremath{\tau}}
\newunicodechar{Φ}{\ensuremath{\Phi}}
\newunicodechar{Ψ}{\ensuremath{\Psi}}
\newunicodechar{Ω}{\ensuremath{\Omega}}
\newunicodechar{↦}{\ensuremath{\mapsto}}
\newunicodechar{∈}{\ensuremath{\in}}
\newunicodechar{∉}{\ensuremath{\notin}}
\newunicodechar{∧}{\ensuremath{\wedge}}
\newunicodechar{⇔}{\ensuremath{\Leftrightarrow}}
\newunicodechar{⟺}{\ensuremath{\Longleftrightarrow}}
\newunicodechar{⇒}{\ensuremath{\Rightarrow}}
\newunicodechar{⟹}{\ensuremath{\Longrightarrow}}
\newunicodechar{∘}{\ensuremath{\circ}}
\newunicodechar{∪}{\ensuremath{\cup}}
\newunicodechar{∩}{\ensuremath{\cap}}
\newunicodechar{≡}{\ensuremath{\equiv}}
\newunicodechar{⊂}{\ensuremath{\subset}}
\newunicodechar{⊥}{\ensuremath{\perp}}
\newunicodechar{∃}{\ensuremath{\exists}}
\newunicodechar{∀}{\ensuremath{\forall}}
\newunicodechar{∛}{\ensuremath{\sqrt[3]{\,}}}
\newunicodechar{∜}{\ensuremath{\sqrt[4]{\,}}}
\newunicodechar{⃗}{\ensuremath{{}^{\to}}}
\newunicodechar{ⁿ}{\ifmmode^{n}\else\textsuperscript{\ensuremath{n}}\fi}
\newunicodechar{ˣ}{\ifmmode^{x}\else\textsuperscript{\ensuremath{x}}\fi}
\newunicodechar{ᵇ}{\ifmmode^{b}\else\textsuperscript{\ensuremath{b}}\fi}
\newunicodechar{ʳ}{\ifmmode^{r}\else\textsuperscript{\ensuremath{r}}\fi}
\newunicodechar{ᵘ}{\ifmmode^{u}\else\textsuperscript{\ensuremath{u}}\fi}
\newunicodechar{ʸ}{\ifmmode^{y}\else\textsuperscript{\ensuremath{y}}\fi}
\newunicodechar{ᵃ}{\ifmmode^{a}\else\textsuperscript{\ensuremath{a}}\fi}
\newunicodechar{ᵉ}{\ifmmode^{e}\else\textsuperscript{\ensuremath{e}}\fi}
\newunicodechar{ᵏ}{\ifmmode^{k}\else\textsuperscript{\ensuremath{k}}\fi}
\newunicodechar{ᵖ}{\ifmmode^{p}\else\textsuperscript{\ensuremath{p}}\fi}
\newunicodechar{ᴺ}{\ifmmode^{N}\else\textsuperscript{\ensuremath{N}}\fi}
\newunicodechar{ᶜ}{\ifmmode^{c}\else\textsuperscript{\ensuremath{c}}\fi}
\newunicodechar{ʰ}{\ifmmode^{h}\else\textsuperscript{\ensuremath{h}}\fi}
\newunicodechar{ᵠ}{\ifmmode^{q}\else\textsuperscript{\ensuremath{q}}\fi}
\newunicodechar{ₙ}{\ifmmode_{n}\else\textsubscript{\ensuremath{n}}\fi}
\newunicodechar{ₐ}{\ifmmode_{a}\else\textsubscript{\ensuremath{a}}\fi}
\newunicodechar{ᵢ}{\ifmmode_{i}\else\textsubscript{\ensuremath{i}}\fi}
\newunicodechar{ᵣ}{\ifmmode_{r}\else\textsubscript{\ensuremath{r}}\fi}
\newunicodechar{ₖ}{\ifmmode_{k}\else\textsubscript{\ensuremath{k}}\fi}
\newunicodechar{ᵦ}{\ifmmode_{\beta}\else\textsubscript{\ensuremath{\beta}}\fi}
\newunicodechar{ₓ}{\ifmmode_{x}\else\textsubscript{\ensuremath{x}}\fi}
\newfontfamily\popdisplay{ArchivoBlack-Regular.ttf}[Path=../../../fonts/]
\newfontfamily\poplogo{SpaceGrotesk-Bold.ttf}[Path=../../../fonts/]
\newfontfamily\popsemi{PlusJakartaSans-SemiBold.ttf}[Path=../../../fonts/]
\newfontfamily\popxbold{PlusJakartaSans-ExtraBold.ttf}[Path=../../../fonts/]
\newfontfamily\popxboldls{PlusJakartaSans-ExtraBold.ttf}[Path=../../../fonts/,LetterSpace=3.0]

\setlist[enumerate]{leftmargin=1.7em,itemsep=5pt,topsep=4pt}
\setlist[itemize]{leftmargin=1.7em,itemsep=4pt,topsep=4pt,label={\color{poporange}\textbullet}}
\setlength{\parindent}{0pt}
\setlength{\parskip}{5pt}
\renewcommand{\arraystretch}{1.25}

% pgfplots : style Pop par défaut
\pgfplotsset{
  every axis/.append style={axis line style={popink,line width=1pt},tick style={popink},
    grid style={popline,line width=0.5pt},label style={font=\small},tick label style={font=\footnotesize}},
  popcurve/.style={poporange,line width=1.8pt,no markers,smooth},
}
% Même style disponible dans un \draw TikZ simple (hors pgfplots)
\tikzset{popcurve/.style={poporange,line width=1.8pt}}
\tikzset{popgrid/.style={popline,very thin}}
\colorlet{popcurve}{poporange} % le nom du style est parfois utilisé comme couleur

% ============================================================
% Pill / squiggle
% ============================================================
\newcommand{\poppill}[3]{%
  \tikz[baseline=-0.55ex]\node[draw=popink,line width=1.2pt,rounded corners=2.6pt,%
    fill=#2,inner xsep=6.5pt,inner ysep=3pt]{%
    \popxboldls\fontsize{7}{7}\selectfont\color{#3}\MakeUppercase{#1}};%
}
\newcommand{\popsquiggle}[1]{%
  \tikz[baseline=-0.4ex]{
    \draw[#1,line width=2.6pt,line cap=round]
      plot [smooth] coordinates {(0,0.09) (0.34,0.01) (0.68,0.21) (1.02,0.01) (1.36,0.10)};
  }%
}
% Mot-clé surligné (marqueur orange sous le texte)
\newcommand{\cle}[1]{{\popxbold\color{popdark}#1}}

% ============================================================
% En-tête / pied
% ============================================================
\pagestyle{fancy}
\fancyhf{}
\renewcommand{\headrulewidth}{0pt}
\renewcommand{\footrulewidth}{0pt}
\lhead{\raisebox{-0.15ex}{\poplogo\fontsize{13}{13}\selectfont\color{popink}OnBuch\color{poporange}+}}
\rhead{\poppill{\DOCMATIERE\ \textbullet\ \DOCNIVEAU\ \textbullet\ Leçon \DOCLECON}{white}{popink}}
\lfoot{\popxbold\fontsize{7.6}{7.6}\selectfont\color{popmuted}\MakeUppercase{Cours signé OnBuch+}\ \textbullet\ {\popsemi\DOCTITRE}}
\rfoot{\tikz[baseline=-0.6ex]\node[rounded corners=8.5pt,fill=popink,minimum height=17pt,minimum width=17pt,inner xsep=5pt,inner sep=0]{\popxbold\fontsize{8}{8}\selectfont\color{popcream}\thepage\,/\,\pageref*{LastPage}};}

% ============================================================
% Titres
% ============================================================
\newcommand{\chaptitre}[2]{%
  \begin{tcolorbox}[enhanced,colback=poporange,colframe=popink,boxrule=2.6pt,arc=11pt,
      drop shadow=popink,left=15pt,right=15pt,top=12pt,bottom=11pt,before skip=3pt,after skip=18pt]
    \poppill{#2}{popink}{popcream}\par\vspace{9pt}
    {\raggedright\hyphenpenalty=10000\exhyphenpenalty=10000\popdisplay\fontsize{22}{24}\selectfont\color{popcream}\MakeUppercase{#1}\par}
  \end{tcolorbox}
}
% Section numérotée : pastille orange avec le numéro + titre Archivo
\newcounter{popsec}
\newcommand{\coursec}[1]{%
  \stepcounter{popsec}\setcounter{popsub}{0}%
  \par\vspace{16pt}\noindent
  \begin{minipage}{\linewidth}
  \tikz[baseline=-0.7ex]\node[draw=popink,line width=1.6pt,fill=poporange,rounded corners=4pt,minimum size=22pt,inner sep=0]{\popdisplay\fontsize{12}{12}\selectfont\color{popcream}\thepopsec};%
  \hspace{8pt}{\popdisplay\fontsize{15}{17}\selectfont\color{popink}\MakeUppercase{#1}}\par
  \vspace{4pt}\noindent{\color{poporange}\rule{36pt}{3.6pt}}
  \end{minipage}\par\nopagebreak\vspace{8pt}
}
\newcounter{popsub}
\newcommand{\courssub}[1]{%
  \stepcounter{popsub}%
  \par\vspace{9pt}\noindent{\popxbold\fontsize{11.5}{13}\selectfont\color{popink}{\color{poporange}\thepopsec.\thepopsub}\hspace{6pt}#1}\par\nopagebreak\vspace{4pt}
}

% ============================================================
% Boîtes pédagogiques — même recette : fond blanc, bordure épaisse colorée,
% ombre dure encre, badge plein. Titre optionnel [#1].
% ============================================================
\tcbset{popbase/.style={breakable,enhanced,colback=white,boxrule=2.3pt,arc=9pt,
  shadow={3pt}{-3pt}{0pt}{popink},left=14pt,right=14pt,top=11pt,bottom=11pt,before skip=11pt,after skip=15pt,
  fontupper=\popsemi\fontsize{10.6}{14.5}\selectfont\color{popink},
  fontlower=\popsemi\fontsize{10.6}{14.5}\selectfont\color{popink}}}
% Coche dessinée (le glyphe ✓ n'existe pas dans les polices du document)
\newcommand{\popcheck}[1]{\tikz[baseline=-0.5ex]\draw[#1,line width=1.6pt,line cap=round,line join=round](-0.13,0.0)--(-0.04,-0.09)--(0.13,0.1);}
\providecommand{\coche}{\popcheck{popgreen}}
\providecommand{\resultat}[1]{\par\smallskip\noindent\fcolorbox{popgreen}{popgreenL}{\parbox{\dimexpr\linewidth-2\fboxsep-2\fboxrule\relax}{\textbf{Résultat.} #1}}\par\smallskip}
\newcommand{\popboxhead}[3]{\poppill{#1}{#2}{white}\ifx\relax#3\relax\else\hspace{6pt}{\popxbold\fontsize{10.6}{12}\selectfont\color{popink}#3}\fi\par\vspace{7pt}}

\newtcolorbox{aretenir}[1][]{popbase,colframe=popgreen,before upper={\popboxhead{À retenir}{popgreen}{#1}}}
% Texte en allemand (document de lecture, dialogue, extrait) : cadre
% d'étude. Un titre optionnel (source, auteur) s'affiche dans l'en-tête.
\newtcolorbox{textebox}[1][]{popbase,colframe=popblue,before upper={\popboxhead{Text}{popblue}{#1}\begin{otherlanguage}{ngerman}},after upper={\end{otherlanguage}}}
% Marques ✓ / ✗ (forme correcte / fautive) souvent utilisées par les modèles
\providecommand{\cross}{\textcolor{poppink}{\ensuremath{\times}}}
\providecommand{\xmark}{\cross}\providecommand{\crossmark}{\cross}\providecommand{\cmark}{\ensuremath{\checkmark}}
% Ligne de réponse / justification à compléter (inventée par les modèles)
\providecommand{\justification}[1]{\par\noindent\textit{Justification~:} \dotfill\par}
% \textipa (package tipa non chargé) : la police ne couvre pas l'API -> simple passage
\providecommand{\textipa}[1]{#1}
% Soulignement (ulem non chargé) : \uline, \ul
\providecommand{\uline}[1]{\underline{#1}}\providecommand{\ul}[1]{\underline{#1}}
% \alert (beamer) inventée par les modèles : texte mis en évidence
\providecommand{\alert}[1]{\textbf{\textcolor{poppink}{#1}}}
% variantes ASCII d'ordinaux écrites en macro
\providecommand{\iere}{\textsuperscript{re}}\providecommand{\ieme}{\textsuperscript{e}}\providecommand{\ere}{\textsuperscript{re}}\providecommand{\eme}{\textsuperscript{e}}\providecommand{\er}{\textsuperscript{er}}\providecommand{\ie}{\textsuperscript{e}}
% Ordinaux français écrits en macro par les modèles : 1\ière, 3\ième
\newcommand{\ière}{\textsuperscript{re}}\newcommand{\ième}{\textsuperscript{e}}
% \exemple (inventée par les modèles) : étiquette « Exemple : »
\providecommand{\exemple}{\textbf{Exemple~:}\ }
% \square utilisable aussi hors mode math (cases à cocher)
\AtBeginDocument{\let\squareMath\square\renewcommand{\square}{\ensuremath{\squareMath}}}
% Mot ou phrase en allemand à l'intérieur d'un paragraphe français
\newcommand{\all}[1]{\ifmmode\text{\textit{#1}}\else\begingroup\selectlanguage{ngerman}\itshape #1\endgroup\fi}
\let\esp\all % alias toléré
\newtcolorbox{exemplebox}[1][]{popbase,colframe=poporange,before upper={\popboxhead{Exemple}{poporange}{#1}}}
\newtcolorbox{definition}[1][]{popbase,colframe=popblue,before upper={\popboxhead{Définition}{popblue}{#1}}}
\newtcolorbox{propriete}[1][]{popbase,colframe=poppurple,before upper={\popboxhead{Propriété}{poppurple}{#1}}}
\newtcolorbox{methode}[1][]{popbase,colframe=popgold,before upper={\popboxhead{Méthode}{popgold}{#1}}}
\newtcolorbox{attention}[1][]{popbase,colframe=poppink,before upper={\popboxhead{Attention}{poppink}{#1}}}
\newtcolorbox{experience}[1][]{popbase,colframe=popblue,colback=popblueL,before upper={\popboxhead{Expérience}{popblue}{#1}}}
% « Le savais-tu ? » : carte violette pleine (contexte, culture, Cameroun)
\newtcolorbox{savaistu}[1][]{popbase,colback=poppurple,colframe=popink,
  fontupper=\popsemi\fontsize{10.4}{14.2}\selectfont\color{white},
  before upper={\poppill{Le savais-tu\,?}{white}{poppurple}\ifx\relax#1\relax\else\hspace{6pt}{\popxbold\color{white}#1}\fi\par\vspace{7pt}}}
% Exercice résolu : énoncé en haut, \tcblower puis la solution sur fond crème
\newcounter{popexo}\newcounter{popexr}
\newtcolorbox{exoresolu}[1][]{popbase,colframe=poporange,colbacklower=popcream,
  segmentation style={popink,line width=1.2pt,dashed},
  before upper={\stepcounter{popexr}\popboxhead{Exercice résolu \thepopexr}{poporange}{#1}},
  before lower={\poppill{Solution}{popink}{popcream}\par\vspace{6pt}}}
% Exercice (non résolu en place) — la correction est dans la partie Corrigés
\newtcolorbox{exercice}[1][]{popbase,colframe=popink,
  before upper={\stepcounter{popexo}\popboxhead{Exercice \thepopexo}{popink}{#1}}}
% Corrigé
\newtcolorbox{corrige}[1][]{popbase,colframe=popgreen,colback=popgreenL,
  before upper={\popboxhead{Corrigé}{popgreen}{#1}}}
% Objectifs / prérequis / bilan
\newtcolorbox{objectifs}{popbase,colframe=popink,colback=poporangeL,
  before upper={\popboxhead{Objectifs de la leçon}{popink}{}}}
\newtcolorbox{prerequis}{popbase,colframe=popink,colback=popblueL,
  before upper={\popboxhead{Prérequis}{popblue}{}}}
\newtcolorbox{bilan}{popbase,colframe=popink,colback=white,boxrule=2.8pt,
  before upper={\popboxhead{Fiche bilan}{poporange}{L'essentiel à connaître pour l'examen}}}
% Figure encadrée avec légende
\newtcolorbox{popfigure}[1][]{enhanced,breakable=false,colback=white,colframe=popline,boxrule=1.4pt,arc=8pt,
  left=10pt,right=10pt,top=10pt,bottom=8pt,before skip=10pt,after skip=12pt,halign=center,
  lower separated=false,halign lower=center,
  fontlower=\popsemi\fontsize{9}{11.5}\selectfont\color{popmuted},
  #1}
\newcommand{\legende}[1]{\par\vspace{6pt}{\popsemi\fontsize{9}{11.5}\selectfont\color{popmuted}#1}}

% ============================================================
% Signature OnBuch+
% ============================================================
\newcommand{\poplogoinline}[1]{{\poplogo\fontsize{#1}{#1}\selectfont\color{popink}OnBuch\color{poporange}+}}
\newcommand{\signatureonbuch}{%
  \begin{tcolorbox}[enhanced,colback=popink,colframe=popink,boxrule=2.2pt,arc=10pt,
      drop shadow=poporange,left=14pt,right=14pt,top=10pt,bottom=10pt,before skip=0pt,after skip=0pt]
    \tikz[baseline=-0.7ex]\node[circle,fill=popgreen,draw=popcream,line width=1.3pt,minimum size=22pt,inner sep=0]{\popcheck{white}};%
    \hspace{9pt}\begin{minipage}[c]{0.62\linewidth}
      {\popxbold\fontsize{12}{13}\selectfont\color{popcream}Signé\ \textbullet\ {\poplogo OnBuch\color{poporange}+}}\par\vspace{2pt}
      {\raggedright\popsemi\fontsize{8}{10.5}\selectfont\color{popline}\MakeUppercase{Équipe pédagogique OnBuch+}\\\MakeUppercase{Conforme au programme officiel MINESEC}\par}
    \end{minipage}\hfill
    {\poplogo\fontsize{22}{22}\selectfont\color{popcream}OnBuch\color{poporange}+}
  \end{tcolorbox}%
}

% ============================================================
% Page de garde
% #1 titre · #2 matière · #3 niveau · #4 module · #5 n° leçon · #6 durée ·
% #7 description / objectifs (texte ou itemize)
% ============================================================
\newcommand{\pagegarde}[7]{%
  \thispagestyle{empty}
  \newgeometry{a4paper,margin=1.7cm,top=1.6cm,bottom=1.4cm}%
  \begin{tikzpicture}[remember picture,overlay]
    \fill[poporange!18] ([xshift=-3.2cm,yshift=-3.6cm]current page.north east) circle (4.6cm);
    \fill[poppurple!14] ([xshift=2.4cm,yshift=7.6cm]current page.south west) circle (3.8cm);
  \end{tikzpicture}%
  {\poplogo\fontsize{30}{30}\selectfont\color{popink}OnBuch\color{poporange}+}\hfill
  \raisebox{0.6ex}{\poppill{Cours signé \textbullet\ Premium}{popink}{popcream}}\par
  \vspace{1pt}\popsquiggle{poppurple}\par
  \vspace{8pt}
  \begin{tcolorbox}[enhanced,colback=poporange,colframe=popink,boxrule=2.8pt,arc=13pt,
      drop shadow=popink,left=18pt,right=18pt,top=12pt,bottom=13pt,after skip=10pt]
    \poppill{#2 \textbullet\ #3}{popink}{popcream}\hspace{5pt}\poppill{#4}{white}{popink}\par\vspace{8pt}
    {\popdisplay\fontsize{14}{14}\selectfont\color{popink}LEÇON #5}\par\vspace{4pt}
    {\raggedright\hyphenpenalty=10000\exhyphenpenalty=10000\popdisplay\fontsize{25}{27}\selectfont\color{popcream}\MakeUppercase{#1}\par}
    \vspace{8pt}
    {\popxbold\fontsize{9.5}{9.5}\selectfont\color{popink}\MakeUppercase{Durée conseillée : #6}}
  \end{tcolorbox}
  \begin{tcolorbox}[enhanced,colback=white,colframe=popink,boxrule=2.2pt,arc=10pt,
      drop shadow=popink,left=16pt,right=16pt,top=10pt,bottom=10pt,after skip=8pt]
    \poppill{Dans cette leçon}{poppurple}{white}\par\vspace{5pt}
    \popsemi\fontsize{9.3}{12.6}\selectfont\color{popink}#7
  \end{tcolorbox}
  \noindent\begin{minipage}[t]{0.487\linewidth}
    \begin{tcolorbox}[enhanced,colback=popgreen,colframe=popink,boxrule=2.2pt,arc=10pt,
        drop shadow=popink,left=11pt,right=11pt,top=10pt,bottom=10pt,height=3.1cm,valign=center]
      \tcbox[colback=white,colframe=popink,boxrule=1.3pt,arc=5pt,boxsep=0pt,nobeforeafter,
        left=3pt,right=3pt,top=3pt,bottom=3pt,valign=center]{\qrcode[height=1.9cm]{https://play.google.com/store/apps/details?id=com.luvvix.onbuch&hl=fr}}%
      \hspace{8pt}\begin{minipage}[c]{0.5\linewidth}
        {\popdisplay\fontsize{11}{12.5}\selectfont\color{white}\MakeUppercase{Télécharge OnBuch}}\par\vspace{4pt}
        {\raggedright\popsemi\fontsize{8.4}{11}\selectfont\color{white}Gratuit \textbullet\ Google Play\\Tuteur IA, annales, concours, résultats\par}
      \end{minipage}
    \end{tcolorbox}
  \end{minipage}\hfill
  \begin{minipage}[t]{0.487\linewidth}
    \begin{tcolorbox}[enhanced,colback=poppurple,colframe=popink,boxrule=2.2pt,arc=10pt,
        drop shadow=popink,left=11pt,right=11pt,top=10pt,bottom=10pt,height=3.1cm,valign=center]
      \tikz[baseline=-0.7ex]\node[circle,fill=white,draw=popink,line width=1.3pt,minimum size=1.6cm,inner sep=0]{\popdisplay\fontsize{24}{24}\selectfont\color{poppurple}+};%
      \hspace{8pt}\begin{minipage}[c]{0.6\linewidth}
        {\popdisplay\fontsize{11}{12.5}\selectfont\color{white}\MakeUppercase{Passe à OnBuch+}}\par\vspace{4pt}
        {\raggedright\popsemi\fontsize{8.4}{11}\selectfont\color{white}Tous les cours signés, Tuteur IA illimité, fiches PDF — onglet OnBuch+ de l'app\par}
      \end{minipage}
    \end{tcolorbox}
  \end{minipage}
  \vspace{8pt}
  \popcontenu
  \vfill
  \signatureonbuch
  \restoregeometry
  \newpage
}

% Rangée « Ce cours contient » (page de garde)
\newcommand{\poptile}[3]{% #1 couleur #2 gros texte #3 légende
  \begin{tcolorbox}[enhanced,colback=#1,colframe=popink,boxrule=2pt,arc=9pt,shadow={2.5pt}{-2.5pt}{0pt}{popink},
      left=6pt,right=6pt,top=9pt,bottom=9pt,halign=center,valign=center,height=1.75cm,width=\linewidth,nobeforeafter]
    {\popdisplay\fontsize{12.5}{13}\selectfont\color{white}\MakeUppercase{#2}}\par\vspace{3pt}
    {\centering\popsemi\fontsize{7.8}{9.5}\selectfont\color{white}#3\par}
  \end{tcolorbox}}
\newcommand{\popcontenu}{%
  \par\noindent{\popxboldls\fontsize{7.5}{7.5}\selectfont\color{popmuted}CE COURS SIGNÉ CONTIENT}\par\vspace{7pt}
  \noindent\begin{minipage}[t]{0.235\linewidth}\poptile{popblue}{Cours}{complet, illustré, pas à pas}\end{minipage}\hfill
  \begin{minipage}[t]{0.235\linewidth}\poptile{poporange}{Méthodes}{et exercices résolus}\end{minipage}\hfill
  \begin{minipage}[t]{0.235\linewidth}\poptile{poppink}{Exercices}{type examen + corrigés}\end{minipage}\hfill
  \begin{minipage}[t]{0.235\linewidth}\poptile{popgreen}{Fiche bilan}{l'essentiel à retenir}\end{minipage}\par}

% Sommaire
\newtcolorbox{sommaire}{enhanced,colback=white,colframe=popink,boxrule=2.2pt,arc=10pt,
  drop shadow=popink,left=18pt,right=18pt,top=13pt,bottom=13pt,before skip=6pt,after skip=20pt,
  before upper={\poppill{Sommaire}{poppurple}{white}\par\vspace{9pt}\popsemi\fontsize{10.6}{14.5}\selectfont\color{popink}}}

% Signature de fin de cours — poussée tout en bas de la dernière page
\newcommand{\findecours}{%
  \par\vfill\nopagebreak
  \begin{center}\popsquiggle{poporange}\end{center}\vspace{4pt}
  \signatureonbuch
}
\AtBeginDocument{\def\llbracket{\mathopen{[\mkern-3.5mu[}}\def\rrbracket{\mathclose{]\mkern-3.5mu]}}} % intervalles d'entiers (glyphe ⟦ absent de la police)
% Glyphes absents de Fira Math : redéfinis à partir de glyphes disponibles
\AtBeginDocument{%
  \def\setminus{\mathbin{\backslash}}\let\smallsetminus\setminus
  \def\vdots{\mathord{\vbox{\baselineskip3.2pt\lineskiplimit0pt\kern4pt\hbox{.}\hbox{.}\hbox{.}}}}%
  \def\ddots{\mathinner{\mkern1mu\raise6.5pt\hbox{.}\mkern2mu\raise3.6pt\hbox{.}\mkern2mu\raise0.7pt\hbox{.}\mkern1mu}}%
  \def\triangle{\mathord{\scalebox{0.9}{$\symup{\Delta}$}}}%
  \def\nabla{\mathord{\raisebox{\depth}{\scalebox{1}[-1]{$\symup{\Delta}$}}}}%
  \def\checkmark{\popcheck{popdark}}%
  \def\star{\mathbin{\ast}}\def\bigstar{\mathord{\ast}}%
  \def\diamond{\mathbin{\scalebox{0.6}{$\Diamond$}}}%
  \def\bigcup{\mathop{\mathchoice{\raisebox{-0.3em}{\scalebox{1.7}{$\cup$}}}{\scalebox{1.25}{$\cup$}}{\cup}{\cup}}\displaylimits}%
  \def\bigcap{\mathop{\mathchoice{\raisebox{-0.3em}{\scalebox{1.7}{$\cap$}}}{\scalebox{1.25}{$\cap$}}{\cap}{\cap}}\displaylimits}%
  \def\lesseqqgtr{\mathrel{\lessgtr}}\def\bigoplus{\mathop{\oplus}\displaylimits}%
}
\providecommand{\ssi}{si et seulement si} % abréviation fréquente
\AtBeginDocument{\providecommand{\wideparen}{\overparen}} % arc de cercle

%% FIGFIX-BEGIN v5 — correctifs globaux des figures (docs/onbuch-generation/tools/figfix)
\usepackage{adjustbox}\usepackage{etoolbox}\tcbuselibrary{raster}
% toute figure TikZ/pgfplots plus large que la ligne est réduite à la largeur disponible
\BeforeBeginEnvironment{tikzpicture}{\begin{adjustbox}{max width=\linewidth}}
\AfterEndEnvironment{tikzpicture}{\end{adjustbox}}
% légendes pgfplots lisibles par-dessus les courbes
\pgfplotsset{every axis/.append style={legend style={fill=white,fill opacity=0.92,text opacity=1,draw=black!15,inner sep=3pt}}}
% étiquettes d'arêtes (syntaxe edge["..."]) avec fond blanc
\tikzset{every edge quotes/.append style={fill=white,inner sep=1.2pt}}
%% FIGFIX-END
\begin{document}

\pagegarde{Production et traduction de textes liés à l'environnement et à la santé}{Allemand}{1ère A}{Chapitre 3 : Environnement, santé et bien-être}{ALA14}{2 h}{Leçon mixte de production et de traduction : à partir d'un tract allemand sur la santé et l'environnement, les élèves enrichissent leur lexique thématique, révisent l'impératif, les verbes pronominaux et les subordonnées, puis s'entraînent au thème et à la version avant de rédiger leur propre tract ou lettre.
\vspace{4pt}
\begin{multicols}{2}\small
\begin{itemize}
\item À la fin de cette leçon, je sais lire et comprendre un tract allemand sur l'environnement et la santé.
\item Je sais utiliser le lexique des champs « Umwelt » et « Gesundheit » dans mes productions.
\item Je sais former et employer l'impératif allemand (du, ihr, Sie) pour donner des conseils.
\item Je sais conjuguer et placer correctement les verbes pronominaux (sich erholen, sich bewegen...).
\item Je sais construire des subordonnées (weil, damit, wenn, dass) avec le verbe en fin de phrase.
\item Je sais rédiger un tract, une liste de conseils et une lettre à un ami en allemand.
\end{itemize}\end{multicols}}

\begin{sommaire}
\begin{multicols}{2}
\begin{enumerate}
\item Texte support : un tract pour la santé et l'environnement
\item Lexique thématique : environnement et santé
\item Grammaire 1 : l'impératif pour conseiller
\item Grammaire 2 : pronominaux et subordonnées
\item Traduction : thème et version guidés
\item Production guidée : tract, conseils et lettre à un ami
\item Activité d'intégration
\item Méthodes et exercices résolus
\item Exercices
\item Corrigés des exercices
\item Fiche bilan
\end{enumerate}
\end{multicols}
\end{sommaire}

\begin{objectifs}
\begin{itemize}
\item À la fin de cette leçon, je sais lire et comprendre un tract allemand sur l'environnement et la santé.
\item Je sais utiliser le lexique des champs « Umwelt » et « Gesundheit » dans mes productions.
\item Je sais former et employer l'impératif allemand (du, ihr, Sie) pour donner des conseils.
\item Je sais conjuguer et placer correctement les verbes pronominaux (sich erholen, sich bewegen...).
\item Je sais construire des subordonnées (weil, damit, wenn, dass) avec le verbe en fin de phrase.
\item Je sais rédiger un tract, une liste de conseils et une lettre à un ami en allemand.
\item Je sais traduire de courtes phrases et un texte simple du français vers l'allemand (thème).
\item Je sais traduire un texte court de l'allemand vers le français (version).
\item Je sais relire et corriger ma production : orthographe, majuscules, ordre des mots, cas.
\end{itemize}
\end{objectifs}

\begin{prerequis}
\begin{itemize}
\item Conjugaison des verbes forts et faibles au présent (Seconde)
\item Vocabulaire de base du corps et de la nature (Seconde)
\item Place du verbe en position 2 dans la phrase principale (Seconde)
\item Notions sur les cas (nominatif, accusatif, datif) vues en début d'année
\item Connecteurs simples : und, aber, denn (début d'année)
\end{itemize}
\end{prerequis}

\newpage

\coursec{Texte support : un tract pour la santé et l'environnement}

\courssub{Situation d'accroche et problématique}

À Yaoundé, les élèves du lycée voient chaque jour des tracts contre le paludisme affichés près de l'infirmerie scolaire et des affiches pour le tri des déchets à l'entrée du marché de Mokolo. En Allemagne, les jeunes écrivent des lettres à leurs amis pour les convaincre de participer à une \all{Müllsammelaktion} (collecte de déchets) ou à une campagne de vaccination. Dans les deux pays, le texte court et percutant — le tract, le conseil, la lettre amicale — est l'outil privilégié pour agir sur la santé et l'environnement.

Problématique : comment rédiger en allemand un tract, donner des conseils de santé ou écrire une lettre à un ami sur l'environnement ? Et comment traduire correctement ces textes courts du français vers l'allemand et inversement ? Cette leçon donne toutes les clés. Nous commençons par lire attentivement un tract : c'est le modèle sur lequel tout le chapitre va s'appuyer.

\courssub{Lecture globale du tract}

Lis le texte suivant deux fois : une première fois sans dictionnaire pour saisir le sens général, une deuxième fois avec le glossaire pour vérifier les détails.

\begin{textebox}[Gesund leben – unsere Umwelt schützen !]
Liebe Schülerinnen und Schüler,

unsere Schule ist unser zweites Zuhause. Aber jeden Tag liegen Plastikflaschen und Papier auf dem Schulhof, und viele von uns trinken nur Limonade und sitzen den ganzen Tag. Das muss sich ändern!

Trennt euren Müll! Werft Plastik in die gelbe Tonne und Papier in die blaue Tonne. Trinkt viel Wasser statt Limonade und bewegt euch jeden Tag mindestens dreißig Minuten: geht zu Fuß zur Schule oder fahrt mit dem Fahrrad. Esst mehr Obst und Gemüse und weniger Süßigkeiten.

Warum ist das wichtig? Erstens: Wenn wir die Umwelt schützen, bleiben wir gesund, weil saubere Luft und sauberes Wasser unser Leben verlängern. Zweitens: Wer sich jeden Tag bewegt, schläft besser und lernt besser. Drittens: Ein sauberer Schulhof macht uns allen Freude und zeigt Respekt vor den anderen.

Die Ärzte sagen auch: Lasst euch impfen! Eine Impfung schützt euch und eure Familie vor gefährlichen Krankheiten.

Kommt am Samstag um 9 Uhr zu unserer Müllsammelaktion auf dem Schulhof! Bringt Handschuhe und gute Laune mit. Gemeinsam können wir viel erreichen. Unsere Schule, unsere Gesundheit, unsere Zukunft!

Euer Umweltclub
\end{textebox}

\begin{definition}[Glossaire du tract]
\begin{itemize}
\item \all{der Abfall, die Abfälle} : le déchet, les ordures
\item \all{der Müll} : les déchets, les ordures (synonyme courant de \all{Abfall})
\item \all{die Mülltrennung} : le tri des déchets
\item \all{trennen} : séparer, trier ; \all{Trennt euren Müll !} : Triez vos déchets !
\item \all{die Tonne, die Tonnen} : la poubelle, le bac ; \all{die gelbe Tonne} : le bac jaune (plastique en Allemagne)
\item \all{die Plastikflasche, die Plastikflaschen} : la bouteille en plastique
\item \all{sich bewegen} : bouger, faire de l'exercice physique
\item \all{die Impfung, die Impfungen} : la vaccination ; \all{sich impfen lassen} : se faire vacciner
\item \all{die Krankheit, die Krankheiten} : la maladie ; \all{gefährlich} : dangereux
\item \all{die Müllsammelaktion, die Müllsammelaktionen} : la collecte de déchets (action organisée)
\item \all{die Gesundheit} : la santé ; \all{gesund} : sain, en bonne santé
\item \all{die Umwelt} : l'environnement ; \all{schützen} : protéger
\item \all{der Schulhof, die Schulhöfe} : la cour de récréation
\item \all{die Süßigkeiten} (pluriel) : les bonbons, les friandises
\item \all{die Handschuhe} (pluriel) : les gants ; \all{die gute Laune} : la bonne humeur
\item \all{erreichen} : atteindre, réussir ; \all{die Zukunft} : l'avenir
\end{itemize}
\end{definition}

\courssub{Compréhension globale : qui écrit ? à qui ? pourquoi ?}

Avant d'analyser les détails, réponds aux trois questions fondamentales de tout commentaire de texte.

\begin{exoresolu}[Questions de compréhension globale]
Énoncé. Réponds en français, en t'appuyant sur le texte : 1) Qui écrit ce tract ? 2) À qui s'adresse-t-il ? 3) Dans quel but ?
\tcblower
Solution détaillée.
\begin{enumerate}
\item \textbf{Qui écrit ?} Le club d'environnement du lycée : la signature \all{Euer Umweltclub} (« Votre club environnement ») et le \all{wir}/\all{unsere Schule} montrent que les rédacteurs sont des élèves engagés.
\item \textbf{À qui ?} Aux autres élèves : l'adresse \all{Liebe Schülerinnen und Schüler} (« Chères élèves, chers élèves ») et le tutoiement collectif des impératifs (\all{trennt, trinkt, kommt}) le prouvent.
\item \textbf{Pourquoi ?} Pour convaincre les élèves d'adopter de bons gestes (trier les déchets, boire de l'eau, bouger, se vacciner) et pour les mobiliser pour la collecte de déchets du samedi (\all{Müllsammelaktion}). C'est un texte \emph{incitatif} : il veut faire agir.
\end{enumerate}
\end{exoresolu}

\begin{methode}[Comprendre un tract en quatre étapes]
\begin{enumerate}
\item \textbf{Lire le titre} : il annonce le thème et le ton. Ici \all{Gesund leben – unsere Umwelt schützen !} combine deux idées : la santé et l'environnement.
\item \textbf{Repérer le destinataire} : cherche la formule d'adresse (\all{Liebe...}) et le mode d'interpellation (tutoiement \all{ihr} ou vouvoiement \all{Sie}).
\item \textbf{Souligner les verbes à l'impératif} : ce sont les conseils et les ordres ; ils forment le cœur du message (\all{Trennt ! Trinkt ! Kommt !}).
\item \textbf{Résumer chaque paragraphe en une phrase} : paragraphe 1 = constat du problème ; paragraphe 2 = les conseils ; paragraphe 3 = les arguments ; paragraphe 4 = la vaccination ; paragraphe 5 = l'appel à l'action.
\end{enumerate}
\end{methode}

\courssub{La structure du tract : un texte pour faire agir}

Un tract n'est pas un texte comme les autres : il doit être lu en dix secondes et donner envie d'agir. Il suit une architecture très codifiée, que tu réutiliseras dans tes propres productions.

\begin{propriete}[Les quatre composantes d'un tract]
\begin{enumerate}
\item \textbf{Le titre accrocheur} : court, souvent avec un point d'exclamation, il nomme l'enjeu : \all{Gesund leben – unsere Umwelt schützen !}
\item \textbf{Les conseils} : phrases à l'impératif, verbe en première position : \all{Trennt euren Müll !}
\item \textbf{Les arguments} : ils justifient les conseils, souvent introduits par \all{Erstens... Zweitens... Drittens...} ou par des subordonnées en \all{wenn, weil, wer}.
\item \textbf{L'appel à l'action} : une invitation concrète avec date, heure et lieu : \all{Kommt am Samstag um 9 Uhr zu unserer Müllsammelaktion !}
\end{enumerate}
\end{propriete}

\begin{popfigure}
\begin{center}
\begin{tikzpicture}[node distance=0.55cm,
box/.style={draw=popblue, fill=popblueL, rounded corners=3pt, align=center, text width=5.2cm, minimum height=1cm, font=\small},
arr/.style={-{Stealth[length=2.5mm]}, thick, poporange}]
\node[box, align=center] (t){\textbf{1. Titre accrocheur}\\ \all{Gesund leben -- unsere Umwelt schützen !}};
\node[box, below=of t, align=center] (c){\textbf{2. Conseils (impératif)}\\ \all{Trennt euren Müll ! Trinkt viel Wasser !}};
\node[box, below=of c, align=center] (a){\textbf{3. Arguments}\\ \all{Erstens... Zweitens... Drittens...}\\ \all{Wenn wir die Umwelt schützen, bleiben wir gesund.}};
\node[box, below=of a, fill=poporangeL, draw=poporange, align=center] (e){\textbf{4. Appel à l'action}\\ \all{Kommt am Samstag um 9 Uhr !}};
\draw[arr] (t) -- (c);
\draw[arr] (c) -- (a);
\draw[arr] (a) -- (e);
\end{tikzpicture}
\end{center}
\legende{La structure d'un tract : du titre à l'appel à l'action}
\end{popfigure}

\courssub{Compréhension détaillée : relever les éléments du texte}

\begin{exercice}[À toi de chercher]
Dans le tract, relève : a) trois conseils ; b) deux arguments ; c) un appel à l'action. Recopie les phrases allemandes exactes.
\end{exercice}

\begin{corrige}[Exercice : compréhension détaillée]
\begin{itemize}
\item \textbf{Trois conseils} : \all{Trennt euren Müll !} ; \all{Trinkt viel Wasser statt Limonade !} ; \all{Lasst euch impfen !} (on pouvait aussi citer \all{Esst mehr Obst und Gemüse !}).
\item \textbf{Deux arguments} : \all{Wenn wir die Umwelt schützen, bleiben wir gesund.} ; \all{Wer sich jeden Tag bewegt, schläft besser und lernt besser.}
\item \textbf{Un appel à l'action} : \all{Kommt am Samstag um 9 Uhr zu unserer Müllsammelaktion auf dem Schulhof !}
\end{itemize}
\end{corrige}

\courssub{Recherche grammaticale dans le texte}

Le tract est une mine d'or pour réviser la grammaire du chapitre. Effectue le triple repérage suivant, comme si tu annotais le texte au crayon de couleur.

\textbf{1. Souligne les impératifs.} Ils sont au début de la phrase, le verbe est conjugué à la 2\up{e} personne du pluriel (\all{ihr}) sans pronom : \all{Trennt, Werft, Trinkt, bewegt, Esst, Lasst, Kommt, Bringt}. Remarque la forme particulière des verbes forts à voyelle \all{e} qui devient \all{i/ie} : \all{werfen} donne \all{Werft !}, \all{essen} donne \all{Esst !}. On reconnaît aussi l'impératif avec \all{lassen} + infinitif : \all{Lasst euch impfen !} (« Faites-vous vacciner ! »).

\textbf{2. Encadre les verbes pronominaux.} Ce sont les verbes accompagnés de \all{sich/euch/uns} : \all{Das muss sich ändern} (cela doit changer), \all{bewegt euch} (bougez-vous), \all{Wer sich bewegt} (celui qui bouge), \all{Lasst euch impfen} (faites-vous vacciner).

\textbf{3. Entoure les subordonnées.} Le verbe conjugué est rejeté à la fin : \all{Wenn wir die Umwelt \textbf{schützen}, bleiben wir gesund} ; \all{weil saubere Luft und sauberes Wasser unser Leben \textbf{verlängern}} ; \all{Wer sich jeden Tag \textbf{bewegt}, schläft besser}.

\begin{exemplebox}[Exemples traduits du tract]
\begin{itemize}
\item \all{Trennt euren Müll !} « Triez vos déchets ! »
\item \all{Bewegt euch jeden Tag !} « Faites de l'exercice chaque jour ! »
\item \all{Lasst euch impfen !} « Faites-vous vacciner ! »
\item \all{Wenn wir die Umwelt schützen, bleiben wir gesund.} « Si nous protégeons l'environnement, nous restons en bonne santé. »
\item \all{Kommt am Samstag um 9 Uhr zu unserer Müllsammelaktion !} « Venez samedi à 9 heures à notre collecte de déchets ! »
\end{itemize}
\end{exemplebox}

\begin{attention}[Piège de traduction : \all{sich bewegen}]
Ne traduis jamais mot à mot \all{sich bewegen} par « se bouger » : en français courant et soigné, on dit « faire du sport », « faire de l'exercice » ou « bouger ». De même, \all{die gute Laune mitbringen} ne se traduit pas par « apporter la bonne humeur » mais par « venir de bonne humeur ». Une bonne traduction cherche toujours l'équivalent \emph{naturel} en français, pas le calque de l'allemand.
\end{attention}

\courssub{Premier lexique thématique : deux familles de mots}

Le tract mêle deux champs lexicaux que nous étudierons en détail dans la section suivante. Voici une première carte mentale pour les mémoriser.

\begin{popfigure}
\begin{center}
\begin{center}\tcbox[colback=popblue,colframe=popblue,coltext=white,arc=9pt,boxsep=4pt,left=6pt,right=6pt]{\bfseries\small Tract \all{Gesund leben}}\end{center}
\vspace{-2pt}
\begin{tcbraster}[raster columns=2,raster equal height=rows,raster column skip=6pt,raster row skip=6pt,enhanced,blankest]
\begin{tcolorbox}[enhanced,colback=white,colframe=popgreen,boxrule=0.9pt,arc=3pt,title={\textbf{Umwelt} environnement},fonttitle=\bfseries\scriptsize,coltitle=white,colbacktitle=popgreen,left=4pt,right=4pt,top=2pt,bottom=2pt,boxsep=1pt,toptitle=1.5pt,bottomtitle=1.5pt,halign=left]
\begin{itemize}[nosep,leftmargin=*,itemsep=1pt,font=\scriptsize]
\item \all{der Müll} déchets
\item \all{die Mülltrennung} tri
\item \all{die Tonne} poubelle
\item \all{schützen} protéger
\end{itemize}
\end{tcolorbox}
\begin{tcolorbox}[enhanced,colback=white,colframe=poporange,boxrule=0.9pt,arc=3pt,title={\textbf{Gesundheit} santé},fonttitle=\bfseries\scriptsize,coltitle=white,colbacktitle=poporange,left=4pt,right=4pt,top=2pt,bottom=2pt,boxsep=1pt,toptitle=1.5pt,bottomtitle=1.5pt,halign=left]
\begin{itemize}[nosep,leftmargin=*,itemsep=1pt,font=\scriptsize]
\item \all{die Impfung} vaccin
\item \all{sich bewegen} bouger
\item \all{die Krankheit} maladie
\item \all{gesund} sain
\end{itemize}
\end{tcolorbox}
\end{tcbraster}

\end{center}
\legende{Carte mentale lexicale : environnement et santé}
\end{popfigure}

\begin{center}
\begin{tabularx}{\linewidth}{|X|X|X|}\hline
\rowcolor{popblueL} \textbf{Français} & \textbf{Deutsch} & \textbf{Remarque} \\ \hline
Trier les déchets ! & \all{Trennt den Müll !} & impératif pluriel, verbe en tête \\ \hline
Faites de l'exercice ! & \all{Bewegt euch !} & verbe pronominal : \all{euch} obligatoire \\ \hline
Faites-vous vacciner ! & \all{Lasst euch impfen !} & \all{lassen} + infinitif = « se faire + infinitif » \\ \hline
Si nous protégeons l'environnement & \all{Wenn wir die Umwelt schützen} & subordonnée : verbe à la fin \\ \hline
Venez samedi à 9 h ! & \all{Kommt am Samstag um 9 Uhr !} & \all{am} + jour, \all{um} + heure \\ \hline
\end{tabularx}
\end{center}

\begin{aretenir}[L'essentiel de cette section]
Un tract comprend quatre parties : un titre accrocheur, des conseils à l'impératif, des arguments (souvent en subordonnées avec \all{wenn} et \all{weil}) et un appel à l'action daté et localisé. Pour le comprendre : lire le titre, repérer le destinataire, souligner les impératifs, résumer chaque paragraphe. Pour le traduire : chercher l'équivalent naturel en français et jamais le mot à mot (\all{sich bewegen} = « faire de l'exercice »). Ces quatre outils — structure, impératif, pronominaux, subordonnées — seront approfondis dans les sections suivantes et te serviront à rédiger ton propre tract.
\end{aretenir}

\begin{savaistu}[Le tri des déchets, une passion allemande]
En Allemagne, le tri des déchets (\all{die Mülltrennung}) est une véritable institution : le bac jaune (\all{die gelbe Tonne}) pour les emballages plastiques, le bac bleu pour le papier, le bac brun pour le compost, et des conteneurs de verre trié par couleur dans les rues. Beaucoup d'écoles allemandes ont un \all{Umweltclub} qui organise des \all{Müllsammelaktionen}, exactement comme dans notre tract. Au Cameroun, des initiatives similaires voient le jour à Yaoundé et à Douala, portées par des élèves et des associations de quartier : une belle occasion de pratiquer l'allemand en parlant d'un sujet qui nous concerne tous !
\end{savaistu}

\coursec{Lexique thématique : environnement et santé}

Dans la section précédente, nous avons lu et commenté le tract \all{„Saubere Stadt, gesunde Menschen!“}, distribué par les élèves du club environnement d'un lycée de Yaoundé. Ce texte regorgeait de mots nouveaux : \all{die Umwelt}, \all{der Müll}, \all{die Gesundheit}, \all{sparen}, \all{wegwerfen}... Pour pouvoir, à la fin de cette leçon, rédiger votre propre tract, donner des conseils à un ami ou traduire un court texte sur ces thèmes, il faut d'abord construire solidement votre vocabulaire. C'est l'objectif de cette section : organiser, mémoriser et employer le lexique de l'\cle{environnement} et de la \cle{santé}.

\courssub{Observer d'abord : un mini-texte pour découvrir le lexique}

Avant toute liste de mots, observons le vocabulaire en situation. Lis le texte suivant et souligne tous les mots qui parlent de la nature, des déchets, du corps ou de la maladie.

\begin{textebox}[Unsere Umwelt, unsere Gesundheit]
In unserer Stadt gibt es viel Müll auf den Straßen. Das ist schlecht für die Umwelt und auch für unsere Gesundheit. Die Verschmutzung macht viele Menschen krank: Sie haben Probleme mit dem Körper, zum Beispiel mit den Augen oder mit dem Hals. Was können wir tun? Zuerst: Wir müssen den Müll trennen. Die Mülltrennung ist in Deutschland normal: Papier, Plastik und Glas kommen in verschiedene Tonnen. Wir können auch Plastikflaschen recyceln. Zweitens: Wir sollen Wasser und Energie sparen. Schaltet das Licht aus, wenn ihr das Zimmer verlasst! Drittens: Für die Gesundheit ist Bewegung wichtig. Wer sich jeden Tag bewegt, bleibt fit. Auch die Ernährung spielt eine große Rolle: Esst viel Obst und Gemüse und trinkt viel Wasser! Wenn man krank ist, geht man zum Arzt oder zur Ärztin. Der Arzt gibt Medikamente, und manchmal ist eine Impfung nötig. Nach der Krankheit muss man sich erholen: schlafen, ruhen, sich nicht stressen. Der Klimawandel ist ein großes Problem für alle Länder, auch für Kamerun. Deshalb müssen wir die Umwelt schützen — für uns und für unsere Kinder. Macht mit!
\end{textebox}

\textbf{Glossaire} (nom = article + pluriel) :

\begin{itemize}
\item \all{die Umwelt} (nur Singular) : l'environnement
\item \all{der Müll} (nur Singular) : les ordures, les déchets
\item \all{die Verschmutzung, -en} : la pollution
\item \all{die Mülltrennung, -en} : le tri des déchets
\item \all{die Tonne, -n} : la poubelle, la benne
\item \all{recyceln} : recycler
\item \all{sparen} : économiser
\item \all{ausschalten} (trennbar) : éteindre
\item \all{die Bewegung, -en} : le mouvement, l'exercice physique
\item \all{die Ernährung} (nur Singular) : l'alimentation
\item \all{der Arzt, Ärzte / die Ärztin, -nen} : le médecin / la femme médecin
\item \all{das Medikament, -e} : le médicament
\item \all{die Impfung, -en} : la vaccination
\item \all{sich erholen} : se remettre, se reposer
\item \all{der Klimawandel} (nur Singular) : le changement climatique
\item \all{schützen} : protéger
\end{itemize}

\textbf{Questions de compréhension} (réponds en allemand, phrases complètes) :

\begin{enumerate}
\item Warum ist der Müll auf den Straßen ein Problem?
\item Was bedeutet \all{Mülltrennung}? Nenne ein Beispiel.
\item Was sollen wir mit Wasser und Energie machen?
\item Was ist wichtig für die Gesundheit? Nenne zwei Dinge.
\item Was macht man, wenn man krank ist?
\end{enumerate}

\courssub{Le champ lexical de l'environnement (\all{die Umwelt})}

\begin{definition}[die Umwelt]
\all{Die Umwelt} (féminin, sans pluriel) désigne l'\cle{environnement}, c'est-à-dire la nature et le milieu dans lequel nous vivons. Attention : on dit \all{die Umwelt schützen} (protéger l'environnement) et \all{umweltfreundlich} (écologique, respectueux de l'environnement). Prononciation : « oum-vèlt ».
\end{definition}

Voici les dix mots essentiels du thème. Apprends-les toujours \textbf{avec leur article et leur pluriel} : c'est la clé pour construire des phrases correctes.

\begin{center}
\begin{tabularx}{\linewidth}{|l|l|X|}\hline
\rowcolor{popblueL} \textbf{Deutsch} & \textbf{Français} & \textbf{Beispiel (exemple)} \\ \hline
\all{die Umwelt} (Sg.) & l'environnement & \all{Wir müssen die Umwelt schützen.} \\ \hline
\all{der Müll} (Sg.) & les ordures & \all{Der Müll liegt auf der Straße.} \\ \hline
\all{die Abfälle} (Pl.) & les déchets & \all{Die Abfälle kommen in die Tonne.} \\ \hline
\all{die Mülltrennung, -en} & le tri des déchets & \all{Die Mülltrennung ist wichtig.} \\ \hline
\all{die Verschmutzung, -en} & la pollution & \all{Die Verschmutzung ist groß.} \\ \hline
\all{die Luft, Lüfte} & l'air & \all{Die Luft in der Stadt ist schlecht.} \\ \hline
\all{die Energie, -n} & l'énergie & \all{Wir sparen Energie.} \\ \hline
\all{das Wasser} (Sg.) & l'eau & \all{Trink viel Wasser!} \\ \hline
\all{der Klimawandel} (Sg.) & le changement climatique & \all{Der Klimawandel ist gefährlich.} \\ \hline
\all{der Baum, Bäume} & l'arbre & \all{Pflanzt Bäume in der Schule!} \\ \hline
\all{die Plastikflasche, -n} & la bouteille en plastique & \all{Wirf die Plastikflasche nicht weg!} \\ \hline
\end{tabularx}
\end{center}

Et les verbes qui vont avec :

\begin{itemize}
\item \all{recyceln} : recycler — \all{Wir recyceln Papier und Glas.} « Nous recyclons le papier et le verre. »
\item \all{sparen} : économiser — \all{Man soll Wasser sparen.} « On devrait économiser l'eau. »
\item \all{schützen} : protéger — \all{Wir schützen die Natur.} « Nous protégeons la nature. »
\item \all{verschmutzen} : polluer — \all{Autos verschmutzen die Luft.} « Les voitures polluent l'air. »
\item \all{trennen} : trier, séparer — \all{In Deutschland trennt man den Müll.} « En Allemagne, on trie les déchets. »
\end{itemize}

\begin{definition}[die Mülltrennung]
\all{Die Mülltrennung} (féminin) signifie littéralement « la séparation des ordures », c'est-à-dire le \cle{tri des déchets}. En Allemagne, cette pratique est \textbf{obligatoire} : chaque matière a sa poubelle de couleur — la poubelle \textbf{jaune} pour le plastique et les emballages (\all{die gelbe Tonne}), la poubelle \textbf{bleue} pour le papier, la poubelle \textbf{brune} ou verte pour les déchets organiques (\all{der Biomüll}), et des conteneurs pour le verre. Mot composé : \all{der Müll} + \all{die Trennung} (la séparation).
\end{definition}

\begin{savaistu}[Le Pfand : la consigne allemande]
En Allemagne, quand tu achètes une bouteille d'eau ou de soda, tu paies une petite somme en plus : c'est le \all{Pfand} (la consigne). Quand tu rapportes la bouteille vide au magasin ou à une machine (\all{der Pfandautomat}), tu récupères ton argent ! Résultat : les rues sont propres et presque toutes les bouteilles sont recyclées. Un vrai geste pour l'environnement — et une bonne idée pour nos villes camerounaises, non ?
\end{savaistu}

\courssub{Le champ lexical de la santé (\all{die Gesundheit})}

\begin{definition}[die Gesundheit]
\all{Die Gesundheit} (féminin, sans pluriel) signifie la \cle{santé}. Elle est construite sur l'adjectif \all{gesund} (en bonne santé) + le suffixe \all{-heit}, qui forme des noms féminins abstraits (comme \all{die Krankheit}, la maladie, à partir de \all{krank}, malade). Contraire : \all{krank} (malade) ; nom : \all{die Krankheit, -en} (la maladie). Prononciation : « gué-zount-haït ».
\end{definition}

\begin{center}
\begin{tabularx}{\linewidth}{|l|l|X|}\hline
\rowcolor{popgreenL} \textbf{Deutsch} & \textbf{Français} & \textbf{Beispiel (exemple)} \\ \hline
\all{die Gesundheit} (Sg.) & la santé & \all{Gesundheit ist wichtig.} \\ \hline
\all{gesund / krank} (Adj.) & en bonne santé / malade & \all{Ich bin heute krank.} \\ \hline
\all{die Krankheit, -en} & la maladie & \all{Malaria ist eine Krankheit.} \\ \hline
\all{der Arzt, Ärzte} & le médecin & \all{Der Arzt hilft den Patienten.} \\ \hline
\all{die Ärztin, -nen} & la femme médecin & \all{Die Ärztin kommt aus Douala.} \\ \hline
\all{der Patient, -en / die Patientin, -nen} & le patient / la patiente & \all{Der Patient wartet vor der Tür.} \\ \hline
\all{das Medikament, -e} & le médicament & \all{Nimm das Medikament dreimal am Tag!} \\ \hline
\all{die Impfung, -en} & la vaccination & \all{Die Impfung schützt uns.} \\ \hline
\all{die Ernährung} (Sg.) & l'alimentation & \all{Gute Ernährung hält gesund.} \\ \hline
\all{der Körper, -} & le corps & \all{Sport ist gut für den Körper.} \\ \hline
\all{der Sport} (Sg.) & le sport & \all{Sport macht fit.} \\ \hline
\end{tabularx}
\end{center}

Et les verbes du domaine :

\begin{itemize}
\item \all{sich bewegen} (réfléchi) : bouger, faire de l'exercice — \all{Ich bewege mich jeden Tag.} « Je fais de l'exercice tous les jours. »
\item \all{sich erholen} (réfléchi) : se remettre, se reposer — \all{Nach der Krankheit erholt er sich.} « Après la maladie, il se remet. »
\item \all{sich impfen lassen} : se faire vacciner — \all{Lass dich impfen!} « Fais-toi vacciner ! »
\item \all{gesund essen / leben} : manger / vivre sainement.
\end{itemize}

Pour visualiser et mémoriser ce vocabulaire, voici une carte mentale autour du mot central \all{Gesundheit} :

\begin{popfigure}
\begin{center}
\begin{tikzpicture}[
  zentrum/.style={circle, draw=popgreen, fill=popgreenL, thick, align=center, font=\bfseries, minimum size=2.2cm},
  ast/.style={rectangle, rounded corners, draw=popblue, fill=popblueL, thick, align=center, font=\small, minimum width=2.6cm},
  blatt/.style={rectangle, rounded corners, draw=popmuted, fill=popcream, align=center, font=\scriptsize}
]
\node[zentrum, align=center] (G){die\\Gesundheit};
\node[ast, above left=1.4cm and 2.6cm of G] (K) {der Körper};
\node[ast, above right=1.4cm and 2.6cm of G] (E) {die Ernährung};
\node[ast, below left=1.4cm and 2.6cm of G] (S) {der Sport};
\node[ast, below right=1.4cm and 2.6cm of G] (A) {der Arzt};
\draw[thick, popline] (G) -- (K);
\draw[thick, popline] (G) -- (E);
\draw[thick, popline] (G) -- (S);
\draw[thick, popline] (G) -- (A);
\node[blatt, above=0.25cm of K, align=center]{gesund, krank,\\sich erholen};
\node[blatt, above=0.25cm of E, align=center]{Obst, Gemüse,\\Wasser trinken};
\node[blatt, below=0.25cm of S, align=center]{sich bewegen,\\fit bleiben};
\node[blatt, below=0.25cm of A, align=center]{die Ärztin, das Medikament,\\die Impfung};
\end{tikzpicture}
\end{center}
\legende{Carte mentale du vocabulaire de la santé : un centre, quatre branches, des mots à compléter par l'élève.}
\end{popfigure}

\begin{attention}[Majuscule à tous les noms !]
En allemand, \textbf{tous les noms} prennent une majuscule, où qu'ils soient dans la phrase : \all{der Müll}, \all{die Gesundheit}, \all{das Medikament}, \all{die Impfung}. C'est l'erreur la plus fréquente des francophones, car en français on écrit « la santé », « le médecin » en minuscules. Mauvais : \emph{ich gehe zum arzt}. Correct : \all{Ich gehe zum Arzt.} À l'écrit, la majuscule t'aide aussi à repérer les noms — et donc à trouver les articles et les cas. Dans tes productions, relis toujours tes majuscules !
\end{attention}

\begin{attention}[Chez le médecin : \all{zum} ou \all{zur} ?]
« Aller chez le médecin » se dit avec la préposition \all{zu} + datif : \all{Ich gehe \textbf{zum} Arzt.} (\all{zu} + \all{dem} = \all{zum}, masculin) ; \all{Ich gehe \textbf{zur} Ärztin.} (\all{zu} + \all{der} = \all{zur}, féminin). Ne traduis jamais « chez » par \emph{bei} ici : \all{zum Arzt gehen} est la formule correcte.
\end{attention}

\courssub{Expressions utiles pour donner des conseils}

Pour rédiger un tract ou une lettre de conseils, trois structures sont indispensables. Observe les exemples avant de lire la règle :

\begin{exemplebox}[Trois façons de conseiller]
\begin{itemize}
\item \all{Es ist wichtig, Wasser zu sparen.} « Il est important d'économiser l'eau. »
\item \all{Man soll den Müll trennen.} « On doit trier les déchets. » / \all{Man sollte mehr Sport machen.} « On devrait faire plus de sport. »
\item \all{Vergiss nicht, dich zu impfen!} « N'oublie pas de te faire vacciner ! »
\end{itemize}
\end{exemplebox}

\begin{propriete}[Forme et emploi des structures de conseil]
\begin{enumerate}
\item \all{Es ist wichtig, ... zu + Infinitiv} : le verbe à l'infinitif se place \textbf{à la fin}, précédé de \all{zu}. Avec un verbe à particule, \all{zu} s'insère dans le verbe : \all{Es ist wichtig, das Licht auszuschalten.} (ausschalten $\rightarrow$ aus\textbf{zu}schalten).
\item \all{Man soll(te) + Infinitiv} : \all{sollen} est un modal ; l'infinitif va \textbf{à la fin} de la phrase. \all{soll} = il faut / on doit ; \all{sollte} (Konjunktiv II) = on devrait, plus poli, plus nuancé — c'est la forme idéale pour \textbf{conseiller}.
\item \all{Vergiss nicht, ... zu + Infinitiv} : impératif de \all{vergessen} + négation + infinitif avec \all{zu}. Avec un verbe réfléchi, le pronom s'accorde avec la personne : \all{Vergiss nicht, \textbf{dich} zu bewegen!} (tu) ; \all{Vergesst nicht, \textbf{euch} zu bewegen!} (vous).
\end{enumerate}
\end{propriete}

\begin{exemplebox}[Exemples traduits]
\begin{itemize}
\item \all{Man soll Wasser sparen.} « On doit économiser l'eau. »
\item \all{Vergiss nicht, dich zu impfen!} « N'oublie pas de te faire vacciner ! »
\item \all{Es ist wichtig, sich gesund zu ernähren.} « Il est important de se nourrir sainement. »
\item \all{Man sollte den Müll nicht auf die Straße werfen.} « On ne devrait pas jeter les ordures dans la rue. »
\item \all{Es ist wichtig, die Umwelt zu schützen.} « Il est important de protéger l'environnement. »
\end{itemize}
\end{exemplebox}

\courssub{Verbes à particule du domaine}

\begin{propriete}[Les verbes à particule séparable]
Les verbes comme \all{wegwerfen}, \all{ausschalten}, \all{aufräumen}, \all{mitmachen} se \textbf{séparent} au présent et à l'impératif : la particule (\all{weg-}, \all{aus-}, \all{auf-}, \all{mit-}) va \textbf{à la fin} de la phrase. Exemple : \all{Ich schalte das Licht aus.} « J'éteins la lumière. » Au Perfekt, le \all{ge-} s'insère entre particule et verbe : \all{Ich habe das Licht ausgeschaltet.} « J'ai éteint la lumière. »
\end{propriete}

\begin{center}
\begin{tabularx}{\linewidth}{|l|X|X|}\hline
\rowcolor{poporangeL} \textbf{Verb} & \textbf{Bedeutung} & \textbf{Beispiel} \\ \hline
\all{wegwerfen} (wirft weg, warf weg, hat weggeworfen) & jeter & \all{Wirf den Müll nicht weg!} « Ne jette pas les déchets ! » \\ \hline
\all{ausschalten} & éteindre & \all{Schalte das Licht aus!} « Éteins la lumière ! » \\ \hline
\all{einschalten} & allumer & \all{Schalte das Radio ein.} « Allume la radio. » \\ \hline
\all{aufräumen} & ranger & \all{Räum dein Zimmer auf!} « Range ta chambre ! » \\ \hline
\all{mitmachen} & participer & \all{Mach bei unserer Aktion mit!} « Participe à notre action ! » \\ \hline
\end{tabularx}
\end{center}

\begin{attention}[Faux amis et pièges lexicaux]
\begin{itemize}
\item \all{die Krankheit} ne veut pas dire « la crainte » ! \all{Die Krankheit} = la \textbf{maladie}. « La crainte, la peur » se dit \all{die Angst} ou \all{die Furcht}.
\item \all{bekommen} ne veut pas dire « devenir » ! \all{Ich bekomme ein Geschenk.} = « Je \textbf{reçois} un cadeau. » « Devenir » se dit \all{werden} : \all{Er wird Arzt.} « Il devient médecin. »
\item \all{der Arzt} : le médecin ; ne confonds pas avec \all{die Art} (la sorte, la manière).
\item « L'eau » est \textbf{neutre} : \all{das Wasser} (et non \emph{die Wasser} comme « la » en français).
\item \all{die Medizin} désigne la \textbf{médecine} (la science, la discipline) ; pour « un médicament » précis, utilise \all{das Medikament}.
\end{itemize}
\end{attention}

\courssub{Complément Bac : pour aller plus loin}

\textbf{(Complément Bac — à connaître pour la classe de Terminale, non exigible en Première.)} Les sujets d'examen aiment le vocabulaire abstrait de l'écologie et de la prévention. Retiens dès maintenant :

\begin{itemize}
\item \all{die erneuerbaren Energien} (Pl.) : les énergies renouvelables — \all{die Solarenergie} (l'énergie solaire), \all{die Windenergie} (l'énergie éolienne).
\item \all{der Umweltschutz} (Sg.) : la protection de l'environnement — \all{Umweltschutz ist eine Aufgabe für alle.} « La protection de l'environnement est une tâche pour tous. »
\item \all{die Gesundheitsvorsorge} (Sg.) : la prévention sanitaire — \all{Vorsorge ist besser als Nachsorge.} « Mieux vaut prévenir que guérir. »
\end{itemize}

Remarque la formation des mots composés, très fréquente en allemand : \all{Umwelt} + \all{Schutz} $\rightarrow$ \all{der Umweltschutz} ; \all{Gesundheit} + \all{Vorsorge} $\rightarrow$ \all{die Gesundheitsvorsorge}. C'est le \textbf{dernier élément} qui donne le genre et le pluriel du mot composé.

\courssub{Exercice résolu et entraînement}

\begin{exoresolu}[Complète avec le mot juste]
Complète les phrases avec : \all{Mülltrennung} — \all{Impfung} — \all{sparen} — \all{wegwerfen} — \all{Gesundheit}.

1. \all{Die \dots\ schützt uns vor Krankheiten.}
2. \all{Man soll Wasser und Energie \dots.}
3. \all{Wirf den Müll nicht auf die Straße — \dots\ ist verboten!} (attention, ici il faut l'infinitif employé comme nom)
4. \all{In Deutschland ist die \dots\ obligatorisch.}
5. \all{Sport und gute Ernährung sind gut für die \dots.}

\tcblower

\textbf{Solution détaillée :}

1. \all{Die \textbf{Impfung} schützt uns vor Krankheiten.} « La vaccination nous protège des maladies. » (féminin, sujet, nominatif — cohérent avec l'article \all{die} ; \all{vor} + datif : \all{vor Krankheiten}.)

2. \all{Man soll Wasser und Energie \textbf{sparen}.} « On doit économiser l'eau et l'énergie. » (infinitif à la fin après le modal \all{soll}.)

3. \all{\dots\ \textbf{Wegwerfen} ist verboten!} « Jeter est interdit ! » (infinitif substantivé : neutre, majuscule obligatoire.)

4. \all{In Deutschland ist die \textbf{Mülltrennung} obligatorisch.} « En Allemagne, le tri des déchets est obligatoire. »

5. \all{\dots\ gut für die \textbf{Gesundheit}.} « Le sport et une bonne alimentation sont bons pour la santé. » (accusatif après \all{für} : \all{die Gesundheit} ne change pas de forme.)
\end{exoresolu}

\begin{exercice}[À toi de jouer !]
1. Traduis en allemand : a) « Il est important de protéger l'environnement. » b) « N'oublie pas d'éteindre la lumière ! » c) « On devrait recycler les bouteilles en plastique. »

2. Construis ta propre carte mentale avec \all{Umwelt} au centre et quatre branches de ton choix (par exemple : \all{Müll}, \all{Energie}, \all{Wasser}, \all{Klima}) avec deux mots par branche.

3. Trouve l'intrus et justifie : \all{der Arzt — die Ärztin — das Medikament — der Klimawandel}.
\end{exercice}

\begin{corrige}[Exercice « À toi de jouer ! »]
1. a) \all{Es ist wichtig, die Umwelt zu schützen.} b) \all{Vergiss nicht, das Licht auszuschalten!} (verbe à particule : \all{zu} s'insère dans le verbe : \all{auszuschalten}.) c) \all{Man sollte Plastikflaschen recyceln.}

2. Exemple de réponse : \all{Müll} $\rightarrow$ \all{der Müll, die Mülltrennung} ; \all{Energie} $\rightarrow$ \all{sparen, ausschalten} ; \all{Wasser} $\rightarrow$ \all{das Wasser, trinken} ; \all{Klima} $\rightarrow$ \all{der Klimawandel, schützen}.

3. L'intrus est \all{der Klimawandel} : c'est le seul mot qui appartient au champ lexical de l'environnement ; les trois autres appartiennent au champ lexical de la santé.
\end{corrige}

\begin{aretenir}[L'essentiel du lexique]
\begin{itemize}
\item Apprends chaque nom avec son \textbf{article} et son \textbf{pluriel} : \all{die Umwelt}, \all{der Müll}, \all{die Verschmutzung, -en}, \all{die Gesundheit}, \all{die Krankheit, -en}, \all{der Arzt, Ärzte}, \all{das Medikament, -e}, \all{die Impfung, -en}.
\item Majuscule à \textbf{tous} les noms allemands, sans exception.
\item Trois structures de conseil : \all{Es ist wichtig, ... zu + Inf.} ; \all{Man soll(te) + Inf.} ; \all{Vergiss nicht, ... zu + Inf.}
\item Verbes à particule : la particule va à la fin (\all{Schalte das Licht aus!}).
\item Faux amis : \all{die Krankheit} = la maladie (pas « la crainte ») ; \all{bekommen} = recevoir (pas « devenir »).
\end{itemize}
\end{aretenir}

Ce vocabulaire est maintenant ta boîte à outils. Dans la prochaine section, nous verrons comment l'\textbf{impératif} permet de transformer ces mots en conseils directs et percutants — exactement ce qu'il faut pour rédiger un tract efficace.

\coursec{Grammaire 1 : l'impératif pour conseiller}

Dans la section précédente, nous avons enrichi notre lexique de l'environnement et de la santé : \all{die Umwelt}, \all{der Müll}, \all{die Gesundheit}, \all{sich bewegen}, \all{sparen}… Mais un tract, une affiche ou un conseil à un ami ne se contente pas de nommer les choses : il faut \cle{inciter à agir}. Pour cela, l'allemand, comme le français, utilise un mode verbal spécial : \cle{l'impératif} (\all{der Imperativ}). C'est le mode du conseil, de la consigne et de l'appel à l'action — exactement ce dont nous avons besoin pour rédiger notre tract de la section 6.

\courssub{Observer avant de comprendre}

Lisons d'abord les slogans d'une campagne de santé publique imaginée pour un lycée de Yaoundé. Soulignez mentalement le verbe de chaque phrase : où est-il placé ? Y a-t-il un sujet ?

\begin{textebox}[Slogans du club santé du lycée]
„Iss mehr Obst und Gemüse!“\\
„Trink jeden Tag genug Wasser!“\\
„Beweg dich! Mach Sport!“\\
„Geht zu Fuß zur Schule!“\\
„Trennt den Müll!“\\
„Schützen Sie die Umwelt! Bleiben Sie gesund!“\\
„Sei ruhig, aber sei aktiv!“
\end{textebox}

\textbf{Glossaire :} \all{das Obst} (les fruits, collectif, pas de pluriel) ; \all{das Gemüse} (les légumes, collectif, pas de pluriel) ; \all{sich bewegen} (bouger, faire de l'exercice) ; \all{der Müll} (les déchets, \textit{pl.} die Müllsorten) ; \all{trennen} (trier, séparer) ; \all{schützen} (protéger) ; \all{bleiben} (rester, demeurer).

\textbf{Observations :} dans toutes ces phrases, le verbe occupe la \cle{première position} et il n'y a \emph{pas de pronom sujet} — sauf dans les deux dernières phrases de la campagne, où \all{Sie} apparaît \emph{après} le verbe : \all{Schützen Sie!}. C'est exactement la logique de l'impératif allemand.

\courssub{La règle fondamentale : position du verbe et disparition du sujet}

\begin{propriete}[Position et sujet à l'impératif]
À l'impératif allemand :
\begin{itemize}
\item le verbe conjugué est en \cle{position 1} (tout au début de la phrase) ;
\item le pronom sujet \cle{disparaît} pour \all{du} et \all{ihr} ;
\item pour la forme de politesse, le verbe est suivi de \all{Sie} : verbe + \all{Sie}.
\end{itemize}
\all{Komm!} (Viens !) — \all{Kommt!} (Venez !) — \all{Kommen Sie!} (Venez ! / politesse).
\end{propriete}

Comparaison avec le français : le français dit « Viens ! », « Venez ! » sans sujet, comme l'allemand. Mais le français n'a que deux formes (tu / vous), alors que l'allemand en a \cle{trois} : \all{du} (une personne familière), \all{ihr} (plusieurs personnes familières) et \all{Sie} (politesse, singulier ou pluriel). Le « vous » français correspond donc tantôt à \all{ihr}, tantôt à \all{Sie} : il faut toujours se demander \emph{à qui l'on parle}.

\begin{exemplebox}[Exemples traduits]
\all{Iss mehr Obst!} = « Mange plus de fruits ! » (à un ami)\\
\all{Geht zu Fuß zur Schule!} = « Allez à l'école à pied ! » (à plusieurs camarades)\\
\all{Bleiben Sie gesund!} = « Restez en bonne santé ! » (formule de politesse, sur une affiche)
\end{exemplebox}

\courssub{Formation de l'impératif \all{du}}

\begin{propriete}[Impératif \all{du} : le radical nu]
On prend le \cle{radical} du présent (l'infinitif sans \all{-en}) :
\begin{itemize}
\item \all{kommen} $\rightarrow$ \all{Komm!}
\item \all{trinken} $\rightarrow$ \all{Trink!}
\item \all{sparen} $\rightarrow$ \all{Spar Wasser!} ou \all{Spare Wasser!} (le \all{-e} final est optionnel pour la plupart des verbes).
\end{itemize}
Le \all{-e} est \cle{obligatoire} quand le radical se termine par \all{-t}, \all{-d}, \all{-ig} ou un groupe de consonnes difficile à prononcer (ex. \all{-chn}, \all{-fn}, \all{-tm}) : \all{Arbeite!} (et non \all{Arbt!}), \all{Warte!}, \all{Öffne!}, \all{Entschuldige!}, \all{Regne!} (de \all{regnen}).
\end{propriete}

\begin{propriete}[Verbes à changement de voyelle \all{e} $\rightarrow$ \all{i}/\all{ie}]
Les verbes qui changent leur voyelle radicale de \all{e} en \all{i} ou \all{ie} à la 2\up{e} et 3\up{e} personne du présent \cle{conservent ce changement} à l'impératif \all{du}, et ne prennent \emph{jamais} de \all{-e} :
\begin{itemize}
\item \all{geben} (du gibst) $\rightarrow$ \all{Gib!}
\item \all{nehmen} (du nimmst) $\rightarrow$ \all{Nimm!}
\item \all{lesen} (du liest) $\rightarrow$ \all{Lies!}
\item \all{sprechen} (du sprichst) $\rightarrow$ \all{Sprich!}
\item \all{essen} (du isst) $\rightarrow$ \all{Iss!}
\item \all{treffen} (du triffst) $\rightarrow$ \all{Triff!}
\end{itemize}
En revanche, les verbes à changement \all{a $\rightarrow$ ä} (comme \all{fahren}, \all{schlafen}, \all{waschen}) ne changent \emph{pas} à l'impératif : \all{Fahr!}, \all{Schlaf!}, \all{Wasch!}.
\end{propriete}

\begin{attention}[« Nimm ! » et non « Nehme ! » — Piège n°1]
Les francophones écrivent souvent \all{Nehme!} ou \all{Gebe!} par analogie avec la 1\up{re} personne (\all{ich nehme}). C'est faux : le changement de voyelle est obligatoire et le \all{-e} disparaît : on écrit \all{Nimm!}, \all{Gib!}, \all{Iss!}. À l'inverse, \all{Arbeite!} garde son \all{-e} pour des raisons de prononciation : ne le supprimez pas. \textbf{Astuce :} si le radical finit par \all{-t} ou \all{-d} (ou \all{-ig}), le \all{-e} reste.
\end{attention}

\courssub{Formation des impératifs \all{ihr} et \all{Sie}}

\begin{propriete}[Impératif \all{ihr} et \all{Sie}]
\begin{itemize}
\item \all{ihr} : on reprend la forme du présent de la 2\up{e} personne du pluriel, \cle{sans le pronom} : \all{ihr kommt} $\rightarrow$ \all{Kommt!} ; \all{ihr trennt} $\rightarrow$ \all{Trennt den Müll!}
\item \all{Sie} : on reprend la forme du présent de politesse, \cle{suivie de Sie} : \all{Sie kommen} $\rightarrow$ \all{Kommen Sie!} ; \all{Sie schützen} $\rightarrow$ \all{Schützen Sie die Umwelt!}
\end{itemize}
Remarque : pour \all{ihr} et \all{Sie}, il n'y a \emph{aucun} changement de voyelle : \all{nehmt!}, \all{nehmen Sie!} (et non \all{nimmt!}).
\end{propriete}

\begin{propriete}[Cas particulier : le verbe \all{sein}]
\all{Sein} est le seul verbe totalement irrégulier à l'impératif :
\begin{itemize}
\item \all{du} : \all{Sei ruhig!} (Sois calme !)
\item \all{ihr} : \all{Seid vorsichtig!} (Soyez prudents !)
\item \all{Sie} : \all{Seien Sie gesund!} (Soyez en bonne santé !)
\end{itemize}
\end{propriete}

\begin{propriete}[Impératif des verbes pronominaux (réfléchis)]
Le pronom réfléchi est conservé et suit le verbe ; il s'accorde avec la personne :
\begin{itemize}
\item \all{du} : \all{Beweg dich!} (Bouge ! / Fais de l'exercice !)
\item \all{ihr} : \all{Bewegt euch!} (Bougez ! entre amis)
\item \all{Sie} : \all{Bewegen Sie sich!} (Bougez ! politesse)
\end{itemize}
Autres exemples : \all{Ruh dich aus!} (Repose-toi !), \all{Erkältet euch nicht!} (Ne vous enrhumez pas !), \all{Waschen Sie sich die Hände!} (Lavez-vous les mains !).
\end{propriete}

\begin{propriete}[Verbes à particule séparable à l'impératif]
La particule se détache et va \cle{en fin de phrase} (comme au présent) :
\begin{itemize}
\item \all{aufstehen} $\rightarrow$ \all{Steh früh auf!} / \all{Steht früh auf!} / \all{Stehen Sie früh auf!}
\item \all{ausruhen} $\rightarrow$ \all{Ruh dich aus!} / \all{Ruht euch aus!} / \all{Ruhen Sie sich aus!}
\item \all{anfangen} $\rightarrow$ \all{Fang an!} / \all{Fangt an!} / \all{Fangen Sie an!}
\end{itemize}
\end{propriete}

\begin{propriete}[Impératif négatif : \all{nicht} et \all{kein}]
La négation se place \textbf{après le verbe} (et après l'objet le cas échéant), la particule en toute fin :
\begin{itemize}
\item \all{Rauch nicht!} (Ne fume pas !)
\item \all{Trink kein Alkohol!} (Ne bois pas d'alcool !)
\item \all{Setzen Sie sich nicht in die Sonne!} (Ne vous asseyez pas au soleil !)
\end{itemize}
\end{propriete}

\begin{center}
\begin{tabular}{|l|l|l|l|}
\hline
\rowcolor{popblueL} Infinitif & du & ihr & Sie \\
\hline
kommen & Komm! & Kommt! & Kommen Sie! \\
\hline
nehmen & Nimm! & Nehmt! & Nehmen Sie! \\
\hline
sich bewegen & Beweg dich! & Bewegt euch! & Bewegen Sie sich! \\
\hline
sein & Sei! & Seid! & Seien Sie! \\
\hline
arbeiten & Arbeite! & Arbeitet! & Arbeiten Sie! \\
\hline
aufstehen & Steh auf! & Steht auf! & Stehen Sie auf! \\
\hline
\end{tabular}
\end{center}

\begin{popfigure}
\begin{tikzpicture}[
  box/.style={draw=popblue, fill=popblueL, rounded corners, align=center, text width=3.2cm, minimum height=1.1cm, font=\small},
  head/.style={draw=popdark, fill=popdark, text=white, rounded corners, align=center, text width=3.2cm, minimum height=0.9cm, font=\small\bfseries},
  node distance=0.5cm]
\node[head] (h1) {du};
\node[box, below=of h1, align=center] (a1){Komm!\\ Nimm!\\ Beweg dich!\\ Sei!\\ Steh auf!};
\node[head, right=1.2cm of h1] (h2) {ihr};
\node[box, below=of h2, align=center] (a2){Kommt!\\ Nehmt!\\ Bewegt euch!\\ Seid!\\ Steht auf!};
\node[head, right=1.2cm of h2] (h3) {Sie};
\node[box, below=of h3, align=center] (a3){Kommen Sie!\\ Nehmen Sie!\\ Bewegen Sie sich!\\ Seien Sie!\\ Stehen Sie auf!};
\end{tikzpicture}
\legende{Les trois formes de l'impératif : \all{kommen}, \all{nehmen}, \all{sich bewegen}, \all{sein}, \all{aufstehen}}
\end{popfigure}

\courssub{À qui je parle ? Choisir la bonne forme}

\begin{methode}[Choisir la bonne forme d'impératif en trois étapes]
\begin{enumerate}
\item \cle{Identifier le destinataire} : une seule personne familière (un ami, un frère) $\rightarrow$ \all{du} ; plusieurs personnes familières (des camarades, des lecteurs jeunes) $\rightarrow$ \all{ihr} ; un adulte, un public inconnu, un lecteur à vouvoyer $\rightarrow$ \all{Sie}.
\item \cle{Former le verbe} : radical nu pour \all{du} (avec changement de voyelle \all{e $\rightarrow$ i/ie} le cas échéant, \all{-e} si radical en \all{-t/-d/-ig}) ; forme du présent sans pronom pour \all{ihr} ; forme du présent + \all{Sie} pour la politesse.
\item \cle{Vérifier} : verbe en position 1, pronom réfléchi placé si nécessaire, particule séparable en fin, négation \all{nicht/kein} après le verbe, point d'exclamation, majuscule en début de phrase.
\end{enumerate}
\end{methode}

\begin{popfigure}
\begin{tikzpicture}[
  q/.style={draw=poporange, fill=poporangeL, rounded corners, align=center, text width=4.5cm, minimum height=1cm, font=\small\bfseries},
  r/.style={draw=popgreen, fill=popgreenL, rounded corners, align=center, text width=3.4cm, minimum height=1.4cm, font=\small},
  node distance=1.6cm]
\node[q] (q) {À qui je parle ?};
\node[r, below left=1.4cm and 2.6cm of q, align=center] (r1){Un ami\\ $\rightarrow$ \textbf{du}\\ Komm!};
\node[r, below=1.4cm of q, align=center] (r2){Plusieurs amis\\ $\rightarrow$ \textbf{ihr}\\ Kommt!};
\node[r, below right=1.4cm and 2.6cm of q, align=center] (r3){Politesse / public\\ $\rightarrow$ \textbf{Sie}\\ Kommen Sie!};
\draw[-{Stealth}, thick, poporange] (q) -- (r1);
\draw[-{Stealth}, thick, poporange] (q) -- (r2);
\draw[-{Stealth}, thick, poporange] (q) -- (r3);
\end{tikzpicture}
\legende{Organigramme de choix : \all{du}, \all{ihr} ou \all{Sie} ?}
\end{popfigure}

\courssub{Emplois de l'impératif}

L'impératif sert à \cle{conseiller}, \cle{donner une consigne}, \cle{encourager}, \cle{avertir} ou \cle{inviter}. On le trouve partout dans la vie quotidienne :

\begin{itemize}
\item les \cle{tracts et affiches} : \all{Rettet die Umwelt!} (Sauvez l'environnement !) ;
\item les \cle{recettes de cuisine} : \all{Schneide die Zwiebeln!} (Coupe les oignons !) ;
\item les \cle{modes d'emploi} : \all{Drücken Sie den Knopf!} (Appuyez sur le bouton !) ;
\item les \cle{conseils de santé} : \all{Trink viel Wasser!} (Bois beaucoup d'eau !) ;
\item les \cle{consignes scolaires} : \all{Öffnet das Buch auf Seite zehn!} (Ouvrez le livre à la page dix !).
\end{itemize}

\begin{center}
\begin{tabularx}{\linewidth}{|X|X|X|}
\hline
\rowcolor{popblueL} Français & Deutsch & Remarque \\
\hline
Mange plus de fruits ! & Iss mehr Obst! & Pas de sujet dans les deux langues \\
\hline
Allez à l'école à pied ! & Geht zu Fuß zur Schule! & « vous » entre amis = \all{ihr} \\
\hline
Restez en bonne santé ! & Bleiben Sie gesund! & « vous » de politesse = \all{Sie} après le verbe \\
\hline
Ne fume pas ! & Rauch nicht! & La négation \all{nicht} suit le verbe \\
\hline
Repose-toi bien ! & Ruh dich gut aus! & Pronom réfléchi + particule en fin \\
\hline
Lavez-vous les mains ! & Waschen Sie sich die Hände! & Pronom réfléchi + objet accusatif \\
\hline
\end{tabularx}
\end{center}

\begin{savaistu}[L'infinitif comme impératif sur les panneaux]
Sur les panneaux publics, les avertissements et les modes d'emploi, on utilise souvent l'\cle{infinitif} à la place de l'impératif \all{Sie}. C'est impersonnel, neutre et valable pour tous.
\begin{itemize}
\item \all{Nicht rauchen!} (Défense de fumer / Ne pas fumer)
\item \all{Bitte nicht stören!} (Prière de ne pas déranger)
\item \all{Rechts fahren!} (Roulez à droite)
\item \all{Vor Gebrauch schütteln!} (Agiter avant emploi)
\end{itemize}
Pour un tract scolaire (section 6), l'impératif \all{ihr} (camarades) ou \all{Sie} (public adulte) est plus engageant que l'infinitif.
\end{savaistu}

\begin{attention}[Pièges fréquents des francophones — Récapitulatif]
\begin{itemize}
\item \textbf{Sujet fantôme :} Ne jamais garder le pronom sujet \all{du} ou \all{ihr} : on ne dit pas \all{Du komm!} mais \all{Komm!} (sauf \all{Sie}, qui reste : \all{Kommen Sie!}).
\item \textbf{Apostrophe inexistante :} Le français « va-t'en ! » n'a pas d'équivalent ; on dit simplement \all{Geh weg!} ou \all{Geh fort!}.
\item \textbf{Point d'exclamation :} Il est la norme à l'impératif, y compris dans les consignes écrites (exercices, examens).
\item \textbf{Particule séparable :} Elle va \cle{en fin de phrase} : \all{Mach das Licht aus!} (Éteins la lumière !), pas \all{Ausmach das Licht!}.
\item \textbf{Verbes à voyelle changeante :} \all{du} garde le \all{i/ie} (\all{Gib!}, \all{Lies!}) mais \textbf{pas} le \all{ä} (\all{Fahr!}, pas \all{Fähr!}).
\item \textbf{\all{-e} de soutien :} Obligatoire après \all{-t, -d, -ig, -chn, -fn, -tm} : \all{Arbeite!}, \all{Öffne!}, \all{Entschuldige!}, \all{Atme!}.
\end{itemize}
\end{attention}

\begin{exoresolu}[Mettre à l'impératif]
Énoncé : mettez les phrases suivantes à l'impératif, à la forme indiquée. (a) \all{du / viel Wasser trinken} ; (b) \all{ihr / den Müll trennen} ; (c) \all{Sie / mehr Sport machen} ; (d) \all{du / dich ausruhen} ; (e) \all{Sie / gesund sein} ; (f) \all{du / nicht rauchen} ; (g) \all{ihr / die Fenster aufmachen}.
\tcblower
Solution détaillée :\\
(a) \all{du} : radical de \all{trinken} = \all{trink-}, pas de changement de voyelle $\rightarrow$ \all{Trink viel Wasser!} (Bois beaucoup d'eau !)\\
(b) \all{ihr} : forme du présent sans pronom $\rightarrow$ \all{Trennt den Müll!} (Triez les déchets !)\\
(c) \all{Sie} : présent + \all{Sie} $\rightarrow$ \all{Machen Sie mehr Sport!} (Faites plus de sport !)\\
(d) verbe pronominal à particule : \all{dich} après le verbe, particule \all{aus} en fin $\rightarrow$ \all{Ruh dich aus!} (Repose-toi !)\\
(e) impératif irrégulier de \all{sein} $\rightarrow$ \all{Seien Sie gesund!} (Soyez en bonne santé !)\\
(f) négation \all{du} : \all{Rauch nicht!} (Ne fume pas !)\\
(g) particule séparable \all{ihr} : \all{Macht die Fenster auf!} (Ouvrez les fenêtres !)
\end{exoresolu}

\begin{exercice}[À vous de conseiller]
Votre camarade Ali est souvent fatigué et mange mal. Donnez-lui cinq conseils à l'impératif \all{du} avec : \all{früh aufstehen}, \all{mehr Obst essen}, \all{genug schlafen}, \all{Wasser trinken}, \all{sich bewegen}. Puis réécrivez ces cinq conseils à la forme \all{Sie}, comme le ferait un médecin sur une affiche.
\end{exercice}

\begin{corrige}[Exercice : À vous de conseiller]
Forme \all{du} : \all{Steh früh auf! Iss mehr Obst! Schlaf genug! Trink Wasser! Beweg dich!}\\
Forme \all{Sie} : \all{Stehen Sie früh auf! Essen Sie mehr Obst! Schlafen Sie genug! Trinken Sie Wasser! Bewegen Sie sich!}\\
Points de vigilance : \all{iss} garde le changement de voyelle (\all{e $\rightarrow$ i}) ; \all{steh} et \all{schlaf} n'ont pas de \all{-e} (pas de \all{-t/-d} final) ; le pronom réfléchi \all{dich/sich} suit immédiatement le verbe ; la particule \all{auf} va en fin de phrase.
\end{corrige}

\begin{aretenir}[L'essentiel de l'impératif]
Verbe en position 1, sujet supprimé (sauf \all{Sie} placé après le verbe).\\
Forme \all{du} : radical nu (\all{Komm!}), avec changement de voyelle \all{e $\rightarrow$ i/ie} conservé et sans \all{-e} (\all{Nimm!}, \all{Gib!}, \all{Lies!}, \all{Iss!}) ; \all{-e} obligatoire après \all{-t/-d/-ig} (\all{Arbeite!}, \all{Öffne!}).\\
Forme \all{ihr} : présent sans pronom (\all{Kommt!}, \all{Nehmt!}).\\
Forme \all{Sie} : présent + \all{Sie} (\all{Kommen Sie!}, \all{Nehmen Sie!}).\\
Irrégulier \all{sein} : \all{Sei! Seid! Seien Sie!}.\\
Pronominaux : \all{Beweg dich! Bewegt euch! Bewegen Sie sich!}.\\
Verbes à particule : particule en fin (\all{Steh auf!}).\\
Négation : \all{nicht/kein} après le verbe (\all{Rauch nicht!}).\\
Panneaux publics : infinitif (\all{Nicht rauchen!}).
\end{aretenir}

\coursec{Grammaire 2 : pronominaux et subordonnées}

Dans la section précédente, nous avons appris à donner des conseils à l'impératif : \all{„Trink viel Wasser!“}, \all{„Beweg dich jeden Tag!“}. Vous avez remarqué ce petit mot \all{dich} collé au verbe. C'est un \cle{pronom réfléchi}. Pour rédiger un tract ou une lettre à un ami sur la santé et l'environnement, il ne suffit pas de donner des ordres : il faut aussi \emph{expliquer pourquoi} (\all{weil}), \emph{dire dans quel but} (\all{damit}, \all{um ... zu}) et \emph{situer dans le temps} (\all{wenn}, \all{als}). C'est l'objet de cette section : maîtriser les verbes pronominaux et construire des phrases complexes correctes.

\courssub{Les verbes pronominaux : révision et observation}

\begin{textebox}[Un dimanche pour se reposer]
Am Sonntag stehe ich spät auf. Zuerst wasche ich mich, und dann frühstücke ich mit meiner Familie. Mein Bruder interessiert sich für Sport, aber heute bewegt er sich nicht: er fühlt sich müde. Ich kümmere mich um den Garten, denn die Pflanzen brauchen Wasser. Am Nachmittag erhole ich mich mit einem guten Buch. Meine Mutter sagt: „Ruht euch aus! Die Schule beginnt morgen wieder!“ Am Abend fühlen wir uns alle glücklich und entspannt.
\end{textebox}

\textbf{Glossaire :} \all{sich ausruhen} (se reposer) ; \all{der Garten, die Gärten} (le jardin) ; \all{die Pflanze, -n} (la plante) ; \all{müde} (fatigué) ; \all{entspannt} (détendu) ; \all{frühstücken} (prendre le petit-déjeuner).

\textbf{Questions de compréhension :}
\begin{enumerate}
\item Que fait le narrateur d'abord le dimanche matin ?
\item Pourquoi le frère ne fait-il pas de sport aujourd'hui ?
\item De quoi le narrateur s'occupe-t-il ?
\end{enumerate}

\begin{definition}[Le verbe pronominal]
Un \cle{verbe pronominal} (allemand : \all{reflexives Verb}) est un verbe accompagné d'un \cle{pronom réfléchi} (\all{mich, dich, sich, uns, euch}) qui renvoie au sujet. En français, ce sont les verbes en « se » : \emph{se laver}, \emph{se reposer}, \emph{s'intéresser à}. Certains verbes sont pronominaux en allemand mais pas en français : \all{sich erholen} = se rétablir, récupérer ; \all{sich interessieren für} = s'intéresser à ; \all{sich kümmern um} = s'occuper de.
\end{definition}

\begin{propriete}[Tableau des pronoms réfléchis (accusatif)]
Le pronom réfléchi s'accorde avec le sujet. À l'accusatif (le cas le plus fréquent à ce niveau) :

\begin{center}
\begin{tabular}{|l|l|l|}
\hline
\rowcolor{popblueL} Personne & Pronom réfléchi & Exemple \\
\hline
ich & mich & Ich wasche \textbf{mich}. \\
\hline
du & dich & Du fühlst \textbf{dich} müde. \\
\hline
er / sie / es & sich & Er erholt \textbf{sich}. \\
\hline
wir & uns & Wir bewegen \textbf{uns}. \\
\hline
ihr & euch & Ihr kümmert \textbf{euch} um den Garten. \\
\hline
sie / Sie & sich & Sie interessieren \textbf{sich} für Sport. \\
\hline
\end{tabular}
\end{center}
\end{propriete}

\begin{popfigure}
\begin{tikzpicture}[
  pers/.style={rectangle, rounded corners, draw=popblue, fill=popblueL, minimum width=1.6cm, minimum height=0.8cm, font=\small},
  refl/.style={rectangle, rounded corners, draw=poporange, fill=poporangeL, minimum width=1.6cm, minimum height=0.8cm, font=\small},
  fleche/.style={-{Stealth[length=2.5mm]}, thick, popdark}]
\node[pers] (p1) at (0,0) {ich};
\node[pers] (p2) at (0,-1.1) {du};
\node[pers] (p3) at (0,-2.2) {er / sie / es};
\node[pers] (p4) at (0,-3.3) {wir};
\node[pers] (p5) at (0,-4.4) {ihr};
\node[pers] (p6) at (0,-5.5) {sie / Sie};
\node[refl] (r1) at (4.5,0) {mich};
\node[refl] (r2) at (4.5,-1.1) {dich};
\node[refl] (r3) at (4.5,-2.2) {sich};
\node[refl] (r4) at (4.5,-3.3) {uns};
\node[refl] (r5) at (4.5,-4.4) {euch};
\node[refl] (r6) at (4.5,-5.5) {sich};
\foreach \i in {1,...,6} {\draw[fleche] (p\i) -- (r\i);}
\end{tikzpicture}
\legende{Correspondance sujet / pronom réfléchi à l'accusatif}
\end{popfigure}

\begin{propriete}[Place du pronom réfléchi]
\begin{itemize}
\item Phrase affirmative : le pronom se place \textbf{juste après le verbe conjugué} : \all{Ich erhole mich am Wochenende.} « Je récupère le week-end. »
\item Impératif (2\textsuperscript{e} personne) : le pronom suit le verbe : \all{Erhole dich!} « Repose-toi ! » ; \all{Wascht euch die Hände!} « Lavez-vous les mains ! »
\item Phrase interrogative sans mot interrogatif : \all{Fühlst du dich gut?} « Te sens-tu bien ? »
\item Cas particulier : avec une \textbf{partie du corps} ou un objet possédé, le pronom réfléchi passe au \textbf{datif} (\all{mir, dir, sich, uns, euch, sich}) et la partie du corps prend l'article défini : \all{Ich wasche \textbf{mir} die Hände.} « Je me lave les mains. » ; \all{Putz \textbf{dir} die Zähne!} « Brosse-toi les dents ! »
\end{itemize}
\end{propriete}

\begin{exemplebox}[Les six verbes du chapitre en contexte]
\begin{itemize}
\item \all{Nach der Krankheit erholt sie sich schnell.} « Après la maladie, elle récupère vite. »
\item \all{Wir bewegen uns jeden Tag dreißig Minuten.} « Nous bougeons trente minutes chaque jour. »
\item \all{Ich fühle mich heute nicht gut.} « Je ne me sens pas bien aujourd'hui. »
\item \all{Er wäscht sich die Hände vor dem Essen.} « Il se lave les mains avant le repas. »
\item \all{Viele Jugendliche interessieren sich für den Umweltschutz.} « Beaucoup de jeunes s'intéressent à la protection de l'environnement. »
\item \all{Sie kümmert sich um ihre kranke Großmutter.} « Elle s'occupe de sa grand-mère malade. »
\end{itemize}
\end{exemplebox}

\begin{attention}[Pièges des francophones]
\begin{itemize}
\item Ne pas oublier le pronom : on ne dit jamais \all{Ich erhole} seul ; il faut \all{Ich erhole \textbf{mich}}.
\item \all{sich interessieren} exige la préposition \all{für} (+ accusatif) : \all{Ich interessiere mich \textbf{für} die Gesundheit.} Ne traduisez pas « s'intéresser \emph{à} » par \all{an}.
\item \all{sich kümmern} exige \all{um} (+ accusatif) : \all{Sie kümmert sich \textbf{um} die Kinder.}
\item Attention au faux ami : \all{sich fühlen} = se sentir (bien, malade), pas « remplir ».
\item Ne confondez pas \all{sich bewegen} (bouger son corps, faire de l'exercice) et \all{umziehen} (déménager) : « je bouge » au sens sportif se dit \all{ich bewege mich}.
\end{itemize}
\end{attention}

\courssub{Les subordonnées : la règle d'or du verbe final}

Observez ces phrases tirées d'un tract sur la santé :

\begin{exemplebox}[Observation]
\all{Wir bleiben zu Hause, weil wir krank sind.} « Nous restons à la maison parce que nous sommes malades. »\\
\all{Wir recyceln den Müll, damit die Umwelt sauber bleibt.} « Nous recyclons les déchets pour que l'environnement reste propre. »\\
\all{Wenn es regnet, bleiben wir zu Hause.} « Quand il pleut, nous restons à la maison. »\\[0.3em]
Dans chaque phrase, où se trouve le verbe conjugué de la subordonnée ? Réponse : \textbf{tout à la fin}, séparé par une virgule.
\end{exemplebox}

\begin{propriete}[La règle d'or de la subordonnée]
Dans toute subordonnée introduite par \all{weil} (parce que), \all{damit} (pour que), \all{wenn} (quand, si), \all{als} (quand, au passé), \all{dass} (que), le \textbf{verbe conjugué se place en dernière position} de la subordonnée, et la subordonnée est séparée de la principale par une \textbf{virgule}.\\
Si la subordonnée \textbf{ouvre la phrase}, elle occupe la position 1 : le verbe de la principale vient \textbf{immédiatement après} (construction verbe-verbe) : \all{\textbf{Wenn es regnet, bleiben wir} zu Hause.}
\end{propriete}

\begin{popfigure}
\begin{tikzpicture}[
  block/.style={rectangle, rounded corners, draw=popdark, fill=popcream, minimum height=1cm, align=center, font=\small},
  vblock/.style={rectangle, rounded corners, draw=poporange, fill=poporangeL, minimum height=1cm, align=center, font=\small},
  fleche/.style={-{Stealth[length=2.5mm]}, thick, popgreen}]
\node[block, align=center] (sub) at (0,0){\textbf{Wenn es regnet,}\\ subordonnée = position 1};
\node[vblock, align=center] (v1) at (5,0){\textbf{bleiben}\\ verbe = position 2};
\node[block, align=center] (rest) at (8.8,0){\textbf{wir zu Hause.}\\ reste de la phrase};
\draw[fleche] (sub) -- (v1);
\draw[fleche] (v1) -- (rest);
\node[font=\small, text=popmuted, align=center] at (0,-1.4) {dans la subordonnée, le verbe\\ conjugué \emph{regnet} est en fin};
\draw[-{Stealth[length=2mm]}, poppurple, thick] (0,-0.55) -- (0,-0.95);
\end{tikzpicture}
\legende{Subordonnée en tête : le verbe de la principale suit immédiatement (verbe-verbe)}
\end{popfigure}

\courssub{La cause avec weil}

\begin{propriete}[weil = parce que]
\all{weil} introduit la cause. Le verbe conjugué va à la fin. La subordonnée peut être placée après ou avant la principale :
\begin{itemize}
\item \all{Ich bleibe zu Hause, weil ich krank bin.} « Je reste à la maison parce que je suis malade. »
\item \all{Weil ich krank bin, bleibe ich zu Hause.} « Parce que je suis malade, je reste à la maison. »
\end{itemize}
\end{propriete}

\begin{exemplebox}[Exemples traduits]
\all{Wir schützen die Umwelt, weil wir gesund bleiben wollen.} « Nous protégeons l'environnement parce que nous voulons rester en bonne santé. »\\
\all{Er trinkt viel Wasser, weil er Sport macht.} « Il boit beaucoup d'eau parce qu'il fait du sport. »\\
\all{Weil die Luft in der Stadt schlecht ist, fahren viele Leute Fahrrad.} « Parce que l'air en ville est mauvais, beaucoup de gens font du vélo. »
\end{exemplebox}

\courssub{Le but avec damit et um ... zu}

\begin{propriete}[Exprimer le but]
Deux constructions pour dire « pour » :
\begin{itemize}
\item \all{um ... zu} + infinitif : \textbf{uniquement quand les deux sujets sont identiques}. L'infinitif se place à la fin, précédé de \all{zu}. \all{Wir recyceln, um die Umwelt zu schützen.} « Nous recyclons pour protéger l'environnement. » (c'est \emph{nous} qui recyclons et \emph{nous} qui protégeons).
\item \all{damit} + subordonnée (verbe à la fin) : quand les sujets sont \textbf{différents}. \all{Wir recyceln, damit die Umwelt sauber bleibt.} « Nous recyclons pour que l'environnement reste propre. »
\end{itemize}
\end{propriete}

\begin{attention}[um ... zu ou damit ?]
\all{um ... zu} s'emploie seulement quand le sujet des deux actions est le même ; sinon on utilise \all{damit}. Comparez :
\begin{itemize}
\item \all{Ich spare Wasser, um die Natur zu schützen.} « J'économise l'eau pour protéger la nature. » (même sujet : \emph{je})
\item \all{Ich spare Wasser, damit meine Kinder genug haben.} « J'économise l'eau pour que mes enfants en aient assez. » (sujets différents : \emph{je} / \emph{mes enfants})
\end{itemize}
Erreur fréquente : \all{Ich spare Wasser, um meine Kinder genug zu haben} est impossible car le sujet change.
\end{attention}

\begin{exemplebox}[Exemples traduits]
\all{Er lernt Deutsch, um in Deutschland zu studieren.} « Il apprend l'allemand pour étudier en Allemagne. »\\
\all{Sie isst viel Obst, um gesund zu bleiben.} « Elle mange beaucoup de fruits pour rester en bonne santé. »\\
\all{Der Arzt erklärt alles langsam, damit die Patienten ihn verstehen.} « Le médecin explique tout lentement pour que les patients le comprennent. »\\
\all{Wir pflanzen Bäume, damit die Luft sauberer wird.} « Nous plantons des arbres pour que l'air devienne plus pur. »
\end{exemplebox}

\courssub{Temps et condition : wenn, als, dass}

\begin{propriete}[wenn, als, dass]
\begin{itemize}
\item \all{wenn} = quand, si (présent, futur, ou habitude au passé) : \all{Wenn es regnet, bleiben wir zu Hause.} « Quand il pleut, nous restons à la maison. » ; \all{Wenn ich Zeit habe, gehe ich zum Arzt.} « Si j'ai le temps, je vais chez le médecin. »
\item \all{als} = quand, pour \textbf{un événement unique au passé} : \all{Als ich krank war, bin ich zum Arzt gegangen.} « Quand j'étais malade, je suis allé chez le médecin. »
\item \all{dass} = que (déclaratif) : \all{Ich denke, dass Sport gesund ist.} « Je pense que le sport est bon pour la santé. »
\end{itemize}
Dans les trois cas, le verbe conjugué va à la fin de la subordonnée.
\end{propriete}

\begin{attention}[wenn ou als ?]
Le français utilise « quand » dans les deux cas, mais l'allemand distingue : \all{als} pour un fait passé unique (\all{Als ich ein Kind war, ...} « Quand j'étais enfant... »), \all{wenn} pour le présent, le futur et les habitudes (\all{Wenn ich krank bin, ...}). Autre piège : ne jamais écrire \all{das} (article) pour \all{dass} (conjonction) : \all{Ich weiß, dass du Recht hast.} « Je sais que tu as raison. »
\end{attention}

\begin{center}
\begin{tabularx}{\linewidth}{|X|X|X|}
\hline
\rowcolor{popblueL} Français & Deutsch & Remarque \\
\hline
Je reste à la maison parce que je suis malade. & Ich bleibe zu Hause, weil ich krank \textbf{bin}. & Verbe à la fin en allemand ; le français garde SVO. \\
\hline
Quand il pleut, nous restons à la maison. & Wenn es regnet, \textbf{bleiben} wir zu Hause. & Subordonnée en tête : verbe juste après la virgule. \\
\hline
Il apprend l'allemand pour étudier en Allemagne. & Er lernt Deutsch, um in Deutschland \textbf{zu studieren}. & Même sujet : \all{um ... zu}. \\
\hline
Je recycle pour que la ville reste propre. & Ich recycle, damit die Stadt sauber \textbf{bleibt}. & Sujets différents : \all{damit}. \\
\hline
\end{tabularx}
\end{center}

\begin{methode}[Vérifier une subordonnée en trois étapes]
\begin{enumerate}
\item \textbf{Souligner le connecteur} (\all{weil}, \all{damit}, \all{wenn}, \all{als}, \all{dass}) : il annonce une subordonnée.
\item \textbf{Vérifier la place du verbe} : le verbe conjugué doit être en dernière position de la subordonnée ; avec \all{um ... zu}, l'infinitif précédé de \all{zu} est à la fin.
\item \textbf{Vérifier la virgule} : elle sépare toujours la subordonnée de la principale. Si la subordonnée ouvre la phrase, contrôler que le verbe de la principale vient juste après (verbe-verbe).
\end{enumerate}
\end{methode}

\begin{exoresolu}[Remets les mots dans l'ordre]
Énoncé : Construis une phrase correcte avec : \all{wir / schützen / die Umwelt / weil / gesund / bleiben / wollen / wir}.
\tcblower
Solution : on identifie d'abord la principale : \all{Wir schützen die Umwelt}. Puis la subordonnée avec \all{weil} : le verbe conjugué \all{wollen} va à la fin, l'infinitif \all{bleiben} juste avant lui. Résultat : \all{Wir schützen die Umwelt, weil wir gesund bleiben wollen.} « Nous protégeons l'environnement parce que nous voulons rester en bonne santé. » Vérification : connecteur \all{weil} souligné, verbe \all{wollen} en dernière position, virgule présente.
\end{exoresolu}

\begin{exoresolu}[Traduis en allemand]
Énoncé : Traduis : « Je fais du sport pour rester en bonne santé. » puis « Je fais du sport pour que mon corps reste fort. »
\tcblower
Solution : première phrase — même sujet (\emph{je}) dans les deux actions, donc \all{um ... zu} : \all{Ich mache Sport, um gesund zu bleiben.} Deuxième phrase — sujets différents (\emph{je} / \emph{mon corps}), donc \all{damit} avec le verbe à la fin : \all{Ich mache Sport, damit mein Körper stark bleibt.}
\end{exoresolu}

\begin{exercice}[À toi de jouer]
\begin{enumerate}
\item Conjugue avec le bon pronom réfléchi : a) \all{Du (sich waschen) die Hände.} b) \all{Wir (sich interessieren) für den Umweltschutz.} c) \all{(Sich erholen) du gut am Wochenende?}
\item Complète avec \all{weil}, \all{damit} ou \all{um ... zu} : a) \all{Ich trinke viel Wasser, \dots\ gesund zu bleiben.} b) \all{Sie lernt viel, \dots\ ihre Eltern stolz sind.} c) \all{Wir gehen zu Fuß, \dots\ wir die Luft nicht verschmutzen wollen.}
\item Traduis en allemand : « Quand je suis malade, je vais chez le médecin. »
\end{enumerate}
\end{exercice}

\begin{corrige}[Exercice « À toi de jouer »]
\begin{enumerate}
\item a) \all{Du wäschst dir die Hände.} (pronom réfléchi au datif \all{dir} devant une partie du corps ; \all{Du wäschst dich.} signifierait « tu te laves » entièrement). b) \all{Wir interessieren uns für den Umweltschutz.} c) \all{Erholst du dich gut am Wochenende?}
\item a) \all{um} (même sujet : \emph{je}) ; b) \all{damit} (sujets différents : \emph{elle} / \emph{ses parents}) ; c) \all{weil} (cause).
\item \all{Wenn ich krank bin, gehe ich zum Arzt.} (subordonnée en tête : le verbe \all{gehe} vient juste après la virgule).
\end{enumerate}
\end{corrige}

\begin{aretenir}[L'essentiel de la section]
\begin{itemize}
\item Les verbes pronominaux s'accompagnent toujours du pronom réfléchi accordé au sujet : \all{mich, dich, sich, uns, euch, sich}, placé après le verbe conjugué ; devant une partie du corps, le pronom passe au datif : \all{Ich wasche mir die Hände.}
\item Dans toute subordonnée (\all{weil}, \all{damit}, \all{wenn}, \all{als}, \all{dass}), le verbe conjugué va \textbf{à la fin}, avec une virgule.
\item But avec le même sujet : \all{um ... zu} ; sujets différents : \all{damit}.
\item Subordonnée en tête de phrase : le verbe de la principale suit immédiatement la virgule.
\end{itemize}
\end{aretenir}

\coursec{Traduction : thème et version guidés}

Dans la section précédente, nous avons révisé les verbes pronominaux (\all{sich erholen}, \all{sich freuen}) et les subordonnées avec \all{weil} et \all{dass}, où le verbe conjugué se place à la fin de la phrase. Ces deux points de grammaire sont précisément ceux qui posent le plus de problèmes quand on passe d'une langue à l'autre. Il est donc temps de les mettre à l'épreuve dans l'activité reine du germaniste : la \cle{traduction}. Nous allons apprendre à traduire de l'allemand vers le français (la \cle{version}) et du français vers l'allemand (le \cle{thème}), en suivant une méthode rigoureuse, étape par étape.

\courssub{Qu'est-ce que traduire ? Observer avant d'agir}

Observons d'abord deux phrases déjà traduites. Que remarquez-vous sur la place du verbe ?

\begin{exemplebox}[Deux phrases, deux directions]
\all{Trennt euren Müll, weil der Planet in Gefahr ist.} « Triez vos déchets, parce que la planète est en danger. »

\all{Er erholt sich, weil er krank ist.} « Il se repose parce qu'il est malade. »
\end{exemplebox}

On constate que le traducteur n'a pas aligné les mots un par un : il a \emph{déplacé} le verbe de la subordonnée (\all{ist} se trouve à la fin en allemand, alors que « est » se trouve au milieu en français), il a adapté le pronom réfléchi (\all{sich} devient « se »), et il a choisi des mots français naturels. Traduire, c'est donc \cle{transférer un sens} d'une langue dans une autre, en respectant les lois de construction de chacune.

\begin{center}
\begin{tabularx}{\linewidth}{|X|X|X|}
\hline
\rowcolor{popblueL} Français & Deutsch & Remarque \\
\hline
Il se repose parce qu'il est malade. & Er erholt sich, weil er krank ist. & Verbe \all{ist} renvoyé à la fin de la subordonnée. \\
\hline
Triez vos déchets, parce que la planète est en danger. & Trennt euren Müll, weil der Planet in Gefahr ist. & Impératif pluriel \all{Trennt} ; verbe \all{ist} en fin. \\
\hline
Je pense que le sport est bon pour la santé. & Ich denke, dass Sport gut für die Gesundheit ist. & \all{dass} + verbe \all{ist} en dernière position. \\
\hline
\end{tabularx}
\end{center}

\courssub{La version : de l'allemand vers le français}

\begin{methode}[Méthode de la version en cinq étapes]
\begin{enumerate}
\item \textbf{Lire} tout le texte une première fois, sans dictionnaire, pour saisir le thème général (ici : l'environnement, la santé).
\item \textbf{Analyser} la structure de chaque phrase : où est le sujet ? où est le verbe conjugué ? y a-t-il une subordonnée (verbe à la fin) ? un verbe à particule séparable ?
\item \textbf{Traduire} d'abord le sens global de la phrase, en commençant par le verbe et son sujet, sans chercher le mot à mot.
\item \textbf{Reformuler} en français correct et naturel : une bonne version « sonne français », elle ne sent pas « l'allemand traduit ».
\item \textbf{Relire} en comparant avec l'original : rien n'a été oublié, rien n'a été ajouté.
\end{enumerate}
\end{methode}

\begin{popfigure}
\begin{tikzpicture}[
  etape/.style={rectangle, rounded corners=4pt, draw=popblue, fill=popblueL, align=center, minimum width=2.1cm, minimum height=0.9cm, font=\small\bfseries},
  fleche/.style={-{Stealth[length=2.5mm]}, thick, popdark}
]
\node[etape] (a) at (0,0) {1. Lire};
\node[etape, fill=popgreenL, draw=popgreen] (b) at (2.7,0) {2. Analyser};
\node[etape, fill=poporangeL, draw=poporange] (c) at (5.4,0) {3. Traduire};
\node[etape, fill=poppinkL, draw=poppink] (d) at (8.1,0) {4. Reformuler};
\node[etape, fill=popgoldL, draw=popgold] (e) at (10.8,0) {5. Relire};
\draw[fleche] (a) -- (b);
\draw[fleche] (b) -- (c);
\draw[fleche] (c) -- (d);
\draw[fleche] (d) -- (e);
\end{tikzpicture}
\legende{Organigramme des cinq étapes de la traduction : on ne passe jamais directement de la lecture à l'écriture.}
\end{popfigure}

\begin{aretenir}[Le secret d'une bonne version]
Une version réussie ne ressemble pas à de l'allemand déguisé en français. Si votre phrase française est bizarre (« Il se réjouit sur les vacances »), c'est qu'elle est calquée sur l'allemand (\all{Er freut sich auf die Ferien}). Reformulez : « Il se réjouit à l'idée des vacances » ou « Il attend les vacances avec impatience ». Traduisez les \emph{idées}, pas les \emph{mots}.
\end{aretenir}

\courssub{Version guidée : traduisons ensemble un extrait du tract}

Reprenons quatre phrases du tract étudié dans la première section et traduisons-les collectivement, en justifiant chaque choix.

\begin{textebox}[Extrait du tract „Gesund leben — sauber leben !“]
Liebe Schülerinnen und Schüler! Unsere Schule startet eine neue Aktion für die Umwelt und für eure Gesundheit. Trennt euren Müll, weil der Planet in Gefahr ist. Trinkt jeden Tag genug Wasser und esst mehr Obst und Gemüse. Wenn ihr krank seid, bleibt zu Hause und erholt euch gut. Wir denken, dass eine saubere Schule und ein gesunder Körper zusammengehören. Macht mit !
\end{textebox}

\textbf{Glossaire :} die Aktion, -en (l'action, la campagne) ; die Umwelt (l'environnement) ; die Gesundheit (la santé) ; der Müll (les déchets) ; die Gefahr, -en (le danger) ; das Obst (les fruits) ; das Gemüse (les légumes) ; sich erholen (se reposer, se rétablir) ; zusammengehören (aller ensemble) ; mitmachen (participer).

\textbf{Phrase 1 :} \all{Trennt euren Müll, weil der Planet in Gefahr ist.}

Analyse : impératif pluriel de \all{trennen} ; subordonnée en \all{weil} avec \all{ist} à la fin ; \all{der Planet} est nominatif. Traduction : « Triez vos déchets, parce que la planète est en danger. »

\textbf{Phrase 2 :} \all{Trinkt jeden Tag genug Wasser und esst mehr Obst und Gemüse.}

Analyse : deux impératifs coordonnés (\all{trinkt}, \all{esst}) ; \all{jeden Tag} est un accusatif de temps. Traduction : « Buvez chaque jour assez d'eau et mangez plus de fruits et de légumes. »

\textbf{Phrase 3 :} \all{Wenn ihr krank seid, bleibt zu Hause und erholt euch gut.}

Analyse : subordonnée en \all{wenn} placée en tête, donc le verbe de la principale (\all{bleibt}) vient juste après la virgule ; \all{erholt euch} est un impératif pronominal. Traduction : « Si vous êtes malades, restez à la maison et reposez-vous bien. »

\textbf{Phrase 4 :} \all{Wir denken, dass eine saubere Schule und ein gesunder Körper zusammengehören.}

Analyse : \all{dass} introduit une subordonnée ; attention, \all{zusammengehören} est un verbe à particule \emph{séparable} : dans la subordonnée, il se recolle en un seul bloc à la fin, mais dans une principale il se détacherait (\all{Sie gehören zusammen.}). Traduction : « Nous pensons qu'une école propre et un corps en bonne santé vont ensemble. »

\courssub{Le thème : du français vers l'allemand}

\begin{methode}[Méthode du thème en cinq étapes]
\begin{enumerate}
\item \textbf{Identifier} le temps et le mode du verbe français (présent ? impératif ? passé composé = Perfekt ?).
\item \textbf{Construire} le squelette : sujet + verbe conjugué en \cle{position 2} dans la principale.
\item \textbf{Placer} la subordonnée : virgule obligatoire, conjonction (\all{weil}, \all{dass}, \all{wenn}), verbe conjugué \emph{à la fin}.
\item \textbf{Vérifier} les cas (nominatif, accusatif, datif) et l'article de chaque nom.
\item \textbf{Relire} : majuscule à tous les noms, Umlaute, particule séparable détachée en principale, recollée en subordonnée.
\end{enumerate}
\end{methode}

\begin{exoresolu}[Thème guidé : cinq phrases sur la santé et l'environnement]
Traduisez en allemand : 1. Repose-toi bien, parce que tu es malade. 2. Triez vos déchets ! 3. Je pense que le sport est bon pour la santé. 4. Quand il fait beau, nous nous promenons au parc. 5. Elle se lève tôt parce qu'elle va à l'école.
\tcblower
\textbf{1.} « Repose-toi » = impératif du pronominal \all{sich erholen} : \all{Erhol dich gut, weil du krank bist.} (verbe \all{bist} à la fin ; réfléchi \all{dich} car 2\ieme{} personne).

\textbf{2.} Impératif pluriel : \all{Trennt euren Müll !} (\all{der Müll} est accusatif : \all{euren Müll}).

\textbf{3.} \all{Ich denke, dass Sport gut für die Gesundheit ist.} (subordonnée en \all{dass}, verbe \all{ist} en dernière position).

\textbf{4.} La subordonnée est en tête : le verbe de la principale suit immédiatement la virgule : \all{Wenn das Wetter schön ist, gehen wir im Park spazieren.} (particule \all{spazieren} de \all{spazieren gehen} rejetée en fin de principale).

\textbf{5.} \all{Sie steht früh auf, weil sie zur Schule geht.} (particule \all{auf} séparée en principale ; en subordonnée, \all{geht} reste simple ; \all{zur} = zu + der).
\end{exoresolu}

\courssub{Les pièges de traduction pour les francophones}

\begin{attention}[La subordonnée : virgule + verbe à la fin]
En thème, deux réflexes obligatoires : la \cle{virgule} avant la subordonnée et le \cle{verbe conjugué à la fin}. On écrit \all{... weil ich krank bin} et jamais \all{... weil ich bin krank}. C'est l'erreur la plus fréquente — et la plus sanctionnée — chez les élèves francophones, car le français garde l'ordre sujet-verbe (« parce que je suis malade »).
\end{attention}

\begin{attention}[Faux amis : bekommen n'est pas « devenir »]
\all{bekommen} signifie « recevoir » : \all{Ich bekomme ein Medikament.} « Je reçois un médicament. » Pour dire « devenir », utilisez \all{werden} : \all{Er wird Arzt.} « Il devient médecin. » Autres pièges : \all{die Mappe} = la chemise, la pochette (et non « la mappe ») ; \all{der Kranke} = le malade (c'est un nom, avec majuscule) ; \all{die Krankheit} = la maladie.
\end{attention}

\begin{attention}[Verbes à particule et majuscules]
En principale, la particule se détache et va à la fin : \all{Er steht um 6 Uhr auf.} En subordonnée, le verbe se recolle à la fin : \all{..., weil er um 6 Uhr aufsteht.} Et n'oubliez jamais la majuscule des noms : \all{die Gesundheit}, \all{der Müll}, \all{das Medikament} — même au milieu de la phrase.
\end{attention}

\begin{center}
\begin{tabularx}{\linewidth}{|X|X|X|}
\hline
\rowcolor{popblueL} Piège & Faux & Correct \\
\hline
Verbe en subordonnée & \all{..., weil ich bin krank} & \all{..., weil ich krank bin} \\
\hline
Faux ami & \all{Ich bekomme Arzt.} & \all{Ich werde Arzt.} \\
\hline
Particule séparable & \all{Er aufsteht früh.} & \all{Er steht früh auf.} \\
\hline
Majuscule des noms & \all{die gesundheit} & \all{die Gesundheit} \\
\hline
Virgule & \all{Ich denke dass...} & \all{Ich denke, dass...} \\
\hline
\end{tabularx}
\end{center}

\courssub{La grille de relecture}

Avant de rendre toute traduction ou production, passez chaque phrase au crible de cette grille. C'est le geste professionnel du traducteur.

\begin{aretenir}[Grille de relecture en cinq points]
\begin{enumerate}
\item \textbf{Orthographe} : Umlaute (ä, ö, ü), \all{ß} ou \all{ss}, pas de lettre oubliée.
\item \textbf{Majuscules} : à tous les noms, sans exception.
\item \textbf{Place du verbe} : position 2 en principale, dernière position en subordonnée, particule séparable détachée ou recollée selon le cas.
\item \textbf{Cas} : nominatif pour le sujet, accusatif après \all{für}, \all{haben}, datif après \all{mit}, \all{zu}, \all{von}.
\item \textbf{Pronom réfléchi} : accordé avec le sujet (\all{ich — mich}, \all{du — dich}, \all{er/sie — sich}, \all{wir — uns}, \all{ihr — euch}, \all{sie/Sie — sich}).
\end{enumerate}
\end{aretenir}

\begin{exercice}[À vous de jouer]
Traduisez en français : \all{Wenn wir den Müll trennen, schützen wir die Umwelt.} Puis traduisez en allemand : « Lève-toi tôt, parce que tu vas à l'école. » Appliquez la grille de relecture à votre propre production.
\end{exercice}

\begin{corrige}[Exercice]
Version : « Si nous trions les déchets, nous protégeons l'environnement. » (subordonnée en tête : en français aussi on commence par « si » ; \all{schützen} = protéger).

Thème : \all{Steh früh auf, weil du zur Schule gehst.} Vérifications : impératif du verbe à particule \all{aufstehen} : \all{Steh ... auf} (particule détachée) ; virgule avant \all{weil} ; \all{gehst} à la fin de la subordonnée ; majuscule à \all{Schule} ; \all{zur} = zu + der.
\end{corrige}

Vous disposez maintenant d'une méthode complète et d'une grille de relecture : la section suivante vous demandera de les réinvestir dans une production guidée — tract, conseils et lettre à un ami.

\coursec{Production guidée : tract, conseils et lettre à un ami}

Après avoir traduit des phrases et des textes courts dans les deux sens, vous êtes maintenant prêts pour l'étape la plus importante du chapitre : \cle{produire} vos propres textes en allemand. Au baccalauréat comme dans la vie réelle, on ne se contente pas de comprendre : on écrit un tract pour sensibiliser le quartier, on donne des conseils à un camarade, on écrit une lettre à un ami allemand. Cette section vous donne les modèles, les plans et les phrases toutes faites pour réussir ces trois types de production.

\courssub{Le tract : caractéristiques et structure}

\begin{definition}[Qu'est-ce qu'un tract ?]
Un \cle{tract} (en allemand \all{das Flugblatt, die Flugblätter} ou \all{der Aufruf, die Aufrufe}) est un texte court et percutant, distribué ou affiché, qui vise à \textbf{informer}, \textbf{convaincre} et \textbf{mobiliser} un groupe de personnes autour d'une cause : la santé, l'environnement, le sport. Il s'adresse à tout le monde, donc il utilise souvent l'impératif \all{ihr} ou la première personne du pluriel \all{wir}.
\end{definition}

\begin{textebox}[Modèle de tract : „Gesund und sauber in Yaoundé !"]
\textbf{GESUND UND SAUBER IN YAOUNDÉ !}\par
Liebe Schülerinnen und Schüler,\par
Unsere Schule ist schmutzig, und viele Schüler sind oft krank. Das muss sich ändern !\par
\textbf{Unsere fünf Tipps :}\par
1. Trennt euren Müll ! Plastik und Papier gehören nicht zusammen.\par
2. Wascht euch die Hände vor dem Essen, weil Seife viele Krankheiten stoppt.\par
3. Trinkt jeden Tag sauberes Wasser, damit ihr gesund bleibt.\par
4. Schlaft immer unter einem Moskitonetz, so schützt ihr euch vor Malaria.\par
5. Macht dreimal pro Woche Sport, denn Bewegung macht den Körper stark.\par
\textbf{Kommt alle am Samstag, den 15. Juni, um 9 Uhr auf den Schulhof !} Wir putzen zusammen den Hof und pflanzen Blumen.\par
\textit{Der Gesundheitsclub der Schule}
\end{textebox}

\begin{definition}[Glossaire du tract]
\all{der Müll} : les déchets ; \all{trennen} : trier, séparer ; \all{sich die Hände waschen} : se laver les mains ; \all{die Seife, -n} : le savon ; \all{das Moskitonetz, -e} : la moustiquaire ; \all{sich schützen vor + Dat.} : se protéger de ; \all{die Bewegung} : le mouvement, l'exercice ; \all{der Schulhof, -höfe} : la cour de l'école ; \all{putzen} : nettoyer ; \all{pflanzen} : planter.
\end{definition}

Observez ce tract : il respecte une architecture fixe que vous devez reproduire.

\begin{propriete}[Les six caractéristiques du tract]
\begin{enumerate}
\item Un \textbf{titre choc}, court, souvent en majuscules, avec un point d'exclamation : \all{Gesund und sauber in Yaoundé !}
\item Une \textbf{accroche} qui pose le problème en une ou deux phrases courtes.
\item \textbf{Trois à cinq conseils} numérotés, à l'\textbf{impératif} (\all{Trennt ! Wascht ! Trinkt !}).
\item Des \textbf{arguments} introduits par \all{weil}, \all{damit} ou \all{denn} pour justifier chaque conseil.
\item Un \textbf{appel à l'action} final avec \textbf{date, heure et lieu} : \all{Kommt alle am Samstag ... um ... Uhr ... !}
\item Une \textbf{signature collective} : \all{Der Gesundheitsclub}, \all{Die Schüler der Klasse 1A}.
\end{enumerate}
\end{propriete}

\begin{popfigure}
\begin{tikzpicture}
\node[draw=popblue, fill=popblueL, rounded corners, minimum width=10cm, minimum height=1cm, align=center, font=\bfseries] (titre) at (0,0) {TITRE CHOC + !};
\node[draw=poporange, fill=poporangeL, rounded corners, minimum width=10cm, minimum height=1cm, align=center] (accroche) at (0,-1.3) {Accroche : le problème en 2 phrases};
\node[draw=popgreen, fill=popgreenL, rounded corners, minimum width=10cm, minimum height=2.2cm, align=center] (conseils) at (0,-3.1) {CONSEILS 1 à 5\\ impératifs + arguments (weil / damit / denn)};
\node[draw=poppurple, fill=poppurpleL, rounded corners, minimum width=10cm, minimum height=1cm, align=center] (appel) at (0,-4.9) {APPEL À L'ACTION : Kommt alle am ... um ... Uhr ... !};
\node[draw=popgold, fill=popgoldL, rounded corners, minimum width=10cm, minimum height=0.9cm, align=center] (sign) at (0,-6.1) {Signature collective};
\draw[-{Stealth}, thick, popmuted] (titre) -- (accroche);
\draw[-{Stealth}, thick, popmuted] (accroche) -- (conseils);
\draw[-{Stealth}, thick, popmuted] (conseils) -- (appel);
\draw[-{Stealth}, thick, popmuted] (appel) -- (sign);
\end{tikzpicture}
\legende{Modèle visuel d'un tract : les cinq zones obligatoires}
\end{popfigure}

\begin{methode}[Rédiger un tract efficace en cinq étapes]
\begin{enumerate}
\item \textbf{Choisir un sujet précis} : le paludisme, le tri des déchets, le sport et la santé.
\item \textbf{Trouver un titre choc} : deux à cinq mots, avec un point d'exclamation : \all{Stoppt die Malaria !}, \all{Unser Quartier, unser Leben !}
\item \textbf{Lister quatre conseils à l'impératif}, en choisissant \textbf{une seule personne} (\all{ihr} pour un groupe, \all{wir} pour s'inclure) et en la gardant : \all{Trennt den Müll!} / \all{Geht zu Fuß zur Schule!}
\item \textbf{Ajouter deux arguments} avec \all{weil} (verbe en fin de phrase) ou \all{damit} : \all{Trinkt kein schmutziges Wasser, weil es krank macht.}
\item \textbf{Conclure par un appel} avec date, heure, lieu, puis signer collectivement.
\end{enumerate}
\end{methode}

\begin{attention}[L'impératif dans le tract : les pièges]
\begin{itemize}
\item Impératif \all{ihr} = radical du présent, \textbf{sans} le pronom \all{ihr} : \all{Trennt den Müll !} (et non \all{Ihr trennt den Müll !}).
\item Impératif \all{du} = radical du présent sans \all{-st} (souvent sans \all{-e} final) : \all{Trink Wasser !} (et non \all{Trinkst Wasser !}). \textbf{Exceptions} : verbes en \all{-t/-d/-n/-m} (\all{arbeite!}, \all{warte!}, \all{öffne!}) et verbes à voyelle modifiée (\all{gib!}, \all{hilf!}, \all{lies!}) gardent le \all{-e} ou changent de voyelle.
\item Les verbes à particule séparable : la particule va \textbf{à la fin} de la phrase : \all{Macht mit !}, \all{Hört zu !}.
\item \all{sein} est irrégulier : \all{Sei vorsichtig !} (\all{du}) / \all{Seid vorsichtig !} (\all{ihr}).
\end{itemize}
\end{attention}

\courssub{La lettre à un ami : structure et formules}

\begin{definition}[La lettre amicale]
La \cle{lettre amicale} (\all{der Brief, -e} ; \all{ein Brief an einen Freund}) est un texte personnel qui suit un plan fixe : lieu et date, salutation, corps en paragraphes, formule finale, signature. On tutoie le destinataire (\all{du}).
\end{definition}

\begin{popfigure}
\begin{tikzpicture}[node distance=0.35cm]
\node[draw=popblue, fill=popblueL, rounded corners, minimum width=9cm, align=center] (a) {Lieu et date : Yaoundé, den 10. Mai};
\node[draw=poporange, fill=poporangeL, rounded corners, minimum width=9cm, align=center, below=of a] (b) {Salutation : Lieber Paul, / Liebe Sarah,};
\node[draw=popgreen, fill=popgreenL, rounded corners, minimum width=9cm, align=center, below=of b] (c) {Corps : nouvelles → conseils → invitation};
\node[draw=poppurple, fill=poppurpleL, rounded corners, minimum width=9cm, align=center, below=of c] (d) {Formule finale : Viele Grüße / Liebe Grüße};
\node[draw=popgold, fill=popgoldL, rounded corners, minimum width=9cm, align=center, below=of d] (e) {Signature : Dein Paul / Deine Amina};
\draw[-{Stealth}, thick, popmuted] (a) -- (b);
\draw[-{Stealth}, thick, popmuted] (b) -- (c);
\draw[-{Stealth}, thick, popmuted] (c) -- (d);
\draw[-{Stealth}, thick, popmuted] (d) -- (e);
\end{tikzpicture}
\legende{Plan de la lettre amicale en cinq blocs}
\end{popfigure}

\begin{propriete}[Les formules obligatoires de la lettre amicale]
\begin{itemize}
\item \textbf{Lieu et date} : \all{Yaoundé, den 10. Mai} (préposition \all{den} + accusatif pour la date).
\item \textbf{Salutation} : \all{Lieber} + prénom masculin, \all{Liebe} + prénom féminin, suivis d'une \textbf{virgule} ; la phrase suivante commence par une \textbf{minuscule}.
\item \textbf{Ouverture} : \all{Wie geht es dir ?} / \all{Vielen Dank für deinen Brief.}
\item \textbf{Corps} : un paragraphe de nouvelles, un paragraphe de conseils, un paragraphe d'invitation.
\item \textbf{Formule finale} : \all{Viele Grüße} ou \all{Liebe Grüße}, \textbf{sans virgule}.
\item \textbf{Signature} : \all{Dein} + prénom masculin, \all{Deine} + prénom féminin.
\end{itemize}
\end{propriete}

\begin{exemplebox}[Extraits de lettres traduits]
\all{Liebe Sarah, wie geht es dir ? Mir geht es gut, weil ich jetzt jeden Tag Sport mache.} = « Chère Sarah, comment vas-tu ? Je vais bien parce que je fais maintenant du sport chaque jour. »\par
\all{Wenn du gesund bleiben willst, musst du mehr Obst essen.} = « Si tu veux rester en bonne santé, tu dois manger plus de fruits. »\par
\all{Kommst du am Samstag zu mir ? Wir spielen Fußball im Quartier.} = « Viens-tu chez moi samedi ? Nous jouons au football dans le quartier. »\par
\all{Viele Grüße, Deine Amina} = « Amitiés, ton amie Amina ».
\end{exemplebox}

\begin{attention}[Majuscule ou minuscule pour « du » ?]
Dans la correspondance traditionnelle allemande, \all{du}, \all{dich}, \all{dein} s'écrivent souvent avec une majuscule (\all{Du}, \all{Dich}, \all{Dein}) par politesse affectueuse. La réforme orthographique de 1996/2006 accepte aussi la minuscule. Au baccalauréat, les deux sont tolérés, mais la \textbf{cohérence est exigée} : si vous écrivez \all{Du} au début, écrivez \all{Dein} à la fin. Ne mélangez jamais les deux systèmes dans la même lettre.
\end{attention}

\begin{center}
\begin{tabularx}{\linewidth}{|X|X|X|}
\hline
\rowcolor{popblueL} \textbf{Français} & \textbf{Deutsch} & \textbf{Remarque} \\
\hline
Cher Paul / Chère Sarah & Lieber Paul / Liebe Sarah & virgule obligatoire \\
\hline
Comment vas-tu ? & Wie geht es dir ? & datif \all{dir} \\
\hline
Merci pour ta lettre & Vielen Dank für deinen Brief & accusatif après \all{für} \\
\hline
Je t'invite chez moi & Ich lade dich zu mir ein & verbe à particule \all{einladen} \\
\hline
Amitiés / Grosses bises & Viele Grüße / Liebe Grüße & sans virgule \\
\hline
Ton ami / Ton amie & Dein Paul / Deine Amina & accord avec l'émetteur \\
\hline
\end{tabularx}
\end{center}

\courssub{Banque de phrases modèles et plan guidé de production}

\begin{aretenir}[Banque de phrases modèles à réutiliser]
\begin{itemize}
\item Appel : \all{Kommt alle zu unserem Sporttag !} (Venez tous à notre journée sportive !)
\item Importance : \all{Es ist wichtig, dass wir die Umwelt schützen.} (Il est important que nous protégions l'environnement.)
\item Conseil conditionnel : \all{Wenn du gesund bleiben willst, musst du regelmäßig schlafen.}
\item Argument cause : \all{Wir müssen den Müll trennen, weil Plastik die Natur zerstört.}
\item Argument but : \all{Wir pflanzen Bäume, damit die Luft sauberer wird.}
\item Invitation : \all{Ich hoffe, dass du kommst !}
\end{itemize}
\end{aretenir}

\begin{methode}[Plan guidé de production en quatre temps]
\begin{enumerate}
\item \textbf{Choisir un sujet} parmi : le paludisme, le tri des déchets, le sport et la santé.
\item \textbf{Lister quatre conseils à l'impératif} sur une feuille de brouillon (choisir \all{ihr} ou \all{wir} et s'y tenir).
\item \textbf{Ajouter deux arguments} avec \all{weil} ou \all{damit} (attention : verbe conjugué à la fin de la subordonnée !).
\item \textbf{Rédiger} le texte choisi (tract ou lettre, 10 à 12 lignes) en insérant titre, appel final ou formules de lettre, puis \textbf{relire} avec la grille.
\end{enumerate}
\end{methode}

\begin{exoresolu}[Rédaction guidée d'un tract]
Énoncé : rédigez un tract de 10 à 12 lignes pour le club santé de votre lycée sur le thème « Stoppons le paludisme ».
\tcblower
\textbf{Solution détaillée.}\par
\textbf{STOPPT DIE MALARIA !}\par
Liebe Schüler, jedes Jahr sind viele von uns krank. Das können wir ändern !\par
1. Schlaft immer unter einem Moskitonetz, weil die Moskitos nachts stechen.\par
2. Reinigt euer Zimmer und den Hof, damit es kein stehendes Wasser gibt.\par
3. Geh sofort zum Arzt, wenn du Fieber hast !\par
4. Es ist wichtig, dass wir alle zusammen kämpfen.\par
Kommt alle am Freitag, den 20. Juni, um 14 Uhr in den großen Saal ! Ein Arzt aus Yaoundé spricht über die Malaria.\par
\textit{Der Gesundheitsclub}\par
Vérification : titre choc (1), accroche (2), quatre conseils à l'impératif varié \all{Schlaft / Reinigt / Geh} (3), deux arguments avec \all{weil} et \all{damit} (4), appel avec date, heure, lieu (5), signature collective (6). Verbes en fin de subordonnée : \all{stechen}, \all{gibt}, \all{hast}, \all{kämpfen}. Correct.
\end{exoresolu}

\courssub{Relecture croisée : la grille de correction}

\begin{methode}[La grille de relecture en cinq points]
Échangez votre texte avec un voisin et cochez chaque point :
\begin{enumerate}
\item \textbf{Orthographe} : Umlaute (\all{ä, ö, ü}), \all{ß}, pas de faute sur les mots du chapitre.
\item \textbf{Majuscules} : tous les noms commencent par une majuscule (\all{der Müll}, \all{die Gesundheit}).
\item \textbf{Ordre des mots} : verbe conjugué en deuxième position dans la phrase principale, en \textbf{fin} de phrase après \all{weil}, \all{damit}, \all{dass}, \all{wenn}.
\item \textbf{Cas} : accusatif après \all{für}, \all{durch}, \all{ohne} ; datif après \all{mit}, \all{zu}, \all{vor}, \all{aus}.
\item \textbf{Impératifs} : formes correctes (\all{Trink !}, \all{Trennt !}, \all{Sei !}, \all{Seid !}) et particules séparables rejetées à la fin.
\end{enumerate}
\end{methode}

\begin{exercice}[À vous de produire]
Rédigez au choix : (a) un tract de 10 à 12 lignes sur le tri des déchets dans votre quartier à Douala ; (b) une lettre de 10 à 12 lignes à votre ami allemand Thomas dans laquelle vous donnez vos nouvelles, trois conseils pour rester en bonne santé et une invitation pour les vacances. Utilisez au moins deux subordonnées avec \all{weil} ou \all{damit} et la banque de phrases modèles.
\end{exercice}

\begin{corrige}[Exercice : éléments attendus]
Pour le tract : titre choc, 3 à 5 impératifs variés (personne unique), au moins deux arguments avec verbe final, appel avec date et lieu, signature collective. Pour la lettre : \all{Lieber Thomas,} + minuscule en début de phrase, trois paragraphes (nouvelles, conseils, invitation), \all{Viele Grüße} + \all{Dein} ou \all{Deine} + prénom. Dans les deux cas : majuscules aux noms, verbe en position 2, verbe final dans les subordonnées. Exemple de phrase attendue : \all{Wenn du fit bleiben willst, musst du jeden Tag laufen, weil Sport das Herz stärkt.}
\end{corrige}

\begin{experience}[Relecture croisée en classe]
Par binômes : échangez vos productions. Le correcteur lit à voix haute, souligne en vert ce qui est réussi, en rouge les erreurs, et remplit la grille en cinq points. L'auteur corrige ensuite son texte et le recopie au propre. L'enseignant ramasse les deux versions pour évaluer le progrès.
\end{experience}

\begin{savaistu}[Comment la production écrite est-elle notée au Bac ?]
Au baccalauréat camerounais, la production écrite d'allemand est évaluée selon trois critères complémentaires : la \textbf{richesse du vocabulaire} (mots variés du chapitre, pas de répétitions), la \textbf{correction grammaticale} (ordre des mots, cas, conjugaisons, majuscules) et le \textbf{respect du type de texte demandé}. Un excellent texte qui oublie le titre du tract ou la formule finale de la lettre perd des points, même sans faute de grammaire ! D'où l'importance des plans fixes appris dans cette leçon.
\end{savaistu}

\begin{aretenir}[L'essentiel de la section]
Un tract = titre choc + accroche + 3 à 5 conseils à l'impératif (personne unique) + arguments (\all{weil/damit}) + appel avec date et lieu + signature collective. Une lettre amicale = lieu/date + \all{Lieber/Liebe} + virgule + minuscule, trois paragraphes, \all{Viele Grüße}, \all{Dein/Deine} + prénom. Dans les deux cas : verbe en position 2, verbe final dans les subordonnées, majuscules aux noms, et relecture systématique avec la grille en cinq points.
\end{aretenir}

\newpage

\coursec{Activité d'intégration}

\coursec{Production et traduction de textes liés à l'environnement et à la santé}

\courssub{Objectifs de l'activité}
Cette activité d'intégration vise à mobiliser l'ensemble des compétences travaillées dans le chapitre « Environnement, santé et bien-être » à travers une tâche finale authentique et motivante. Les élèves devront :
\begin{itemize}
    \item \textbf{Produire} un tract de sensibilisation en allemand (genre \emph{Appelltext}) structuré, cohérent et correct sur le thème « Santé et environnement au lycée ».
    \item \textbf{Rédiger} une lettre personnelle à un ami germanophone pour l'inviter à soutenir l'action (genre \emph{privater Brief}).
    \item \textbf{Traduire} du français vers l'allemand (thème) et de l'allemand vers le français (version) en respectant les structures grammaticales ciblées.
    \item \textbf{Relire} et \textbf{corriger} leurs productions (auto-évaluation et évaluation par les pairs) à l'aide d'une grille critères.
    \item \textbf{Réinvestir} le lexique thématique, l'impératif (conseils, interdictions), les verbes pronominaux (actions réflexives, réciproques), les subordonnées causales (\all{weil}) et finales (\all{damit}, \all{um ... zu}), ainsi que l'ordre des mots (V2, verbe en fin de subordonnée).
\end{itemize}

\courssub{Situation}
\begin{textebox}[Contexte : Le lycée bilingue de Nkolbisson (Yaoundé) lance la campagne « Gesund und sauber »]
Le proviseur du \textbf{Lycée Bilingue de Nkolbisson} à Yaoundé souhaite lancer une grande campagne de sensibilisation bilingue (français-allemand) pour améliorer l'hygiène, la propreté et le bien-être dans l'enceinte de l'établissement. Le projet s'intitule \all{„Gesund und sauber – Gesund leben, sauber lernen“}. Il fait écho aux Objectifs de Développement Durable (ODD 3 : Bonne santé et bien-être ; ODD 6 : Eau propre et assainissement ; ODD 13 : Mesures relatives à la lutte contre les changements climatiques).

Une délégation d'élèves germanistes de Première A est chargée de concevoir les supports de communication en \textbf{allemand} pour les affichages dans les couloirs, la cantine, la bibliothèque et les toilettes. Le proviseur a également demandé qu'une \textbf{lettre type} soit rédigée pour être envoyée au lycée partenaire de \textbf{Stuttgart} (Allemagne) dans le cadre du jumelage scolaire, afin d'inviter les correspondants allemands à mener une action parallèle et à échanger des photos.

Le professeur d'allemand fournit un \textbf{message-cadre} en français qu'il faut traduire en allemand pour l'affichage dans la salle des professeurs.
\end{textebox}

\begin{textebox}[Message du proviseur à traduire (thème FR $\to$ DE)]
« Chers collègues, la campagne „Gesund und sauber“ commence lundi prochain. Merci d'encourager les élèves à jeter les déchets dans les poubelles, à économiser l'eau aux robinets et à aérer les salles de classe régulièrement. Votre engagement est important pour la réussite du projet. Cordialement, Le Proviseur. »
\end{textebox}

\courssub{Consignes}
Travail en \textbf{groupes de 4 élèves} (répartition des rôles suggérée : \emph{Rédacteur principal}, \emph{Correcteur grammaire}, \emph{Traducteur}, \emph{Maquettiste/Présentateur}). Durée : 2 h (une séance de 2 h ou deux séances de 1 h).

\begin{enumerate}
    \item \textbf{Analyse du genre « Tract / Affiche de sensibilisation » (15 min)} \\
    En groupe, lisez le modèle de tract ci-dessous (\emph{Modelltext}). Relevez : la structure (titre, accroche, conseils numérotés, arguments, appel final), les formes verbales utilisées (impératif, infinitif, verbes pronominaux), les connecteurs logiques (\all{weil}, \all{damit}, \all{außerdem}, \all{deshalb}). Complétez le tableau d'analyse fourni par le professeur.

    \item \textbf{Rédaction du tract allemand « Gesund und sauber an unserer Schule » (40 min)} \\
    Rédigez \textbf{un seul tract commun par groupe} sur une feuille A3 (ou format numérique). Le tract doit contenir \textbf{obligatoirement} :
    \begin{itemize}
        \item Un \textbf{titre accrocheur} (nominal, impératif ou question).
        \item Une \textbf{phrase d'accroche} (1--2 phrases) qui présente le problème ou l'objectif.
        \item \textbf{5 conseils concrets} à l'\textbf{impératif} (formes \all{du}, \all{ihr} ou \all{Sie} — choix justifié dans le commentaire) pour des actions quotidiennes au lycée (déchets, eau, électricité, air, bruit, alimentation, hygiène des mains...). Variez les verbes (séparables, pronominaux, modaux).
        \item \textbf{2 arguments développés} avec \textbf{subordonnées} : l'un avec \all{weil} (cause), l'autre avec \all{damit} ou \all{um ... zu} (but). Ex. : \all{Wir trennen den Müll, weil wir die Umwelt schützen wollen.} / \all{Wir lüften regelmäßig, damit die Luft frisch bleibt.}
        \item Un \textbf{appel à l'action final} (impératif ou \all{Lasst uns ...!} / \all{Werden Sie mit!}).
        \item Des \textbf{éléments visuels} (pictogrammes, couleurs, mise en page lisible).
    \end{itemize}

    \item \textbf{Rédaction de la lettre à un ami allemand (30 min)} \\
    Individuellement (chaque élève écrit sa propre lettre), rédigez une \textbf{lettre personnelle} (environ 80--100 mots) à \textbf{Lukas} (ou \textbf{Lena}), votre correspondant(e) du lycée partenaire de Stuttgart. Structure :
    \begin{itemize}
        \item En-tête (lieu, date), formule d'appel (\all{Lieber Lukas,} / \all{Liebe Lena,}).
        \item Présentation de la campagne \all{Gesund und sauber} au lycée (quoi, quand, pourquoi).
        \item Description de \textbf{ton rôle personnel} dans le groupe (verbes pronominaux : \all{ich engagiere mich}, \all{ich kümmere mich um ...}, \all{wir treffen uns}).
        \item Invitation à participer / à organiser une action similaire à Stuttgart.
        \item Formules de politesse de fin (\all{Viele Grüße}, \all{Dein/Deine ...}).
    \end{itemize}
    Utilisez au moins \textbf{3 verbes pronominaux} et \textbf{2 subordonnées} (\all{weil}, \all{dass}, \all{damit}, \all{als}, \all{wenn}).

    \item \textbf{Traduction guidée : Version et Thème (25 min)} \\
    \begin{enumerate}
        \item \textbf{Version (DE $\to$ FR)} : Traduisez \textbf{votre propre tract} (ou celui du groupe) en français naturel, soigné, adapté à l'affichage bilingue. Attention aux faux-amis (\all{Gift} $\neq$ gift, \all{aktuell} $\neq$ actuel), aux majuscules des noms, à l'ordre des mots.
        \item \textbf{Thème (FR $\to$ DE)} : Traduisez le \textbf{message du proviseur} (fourni dans la situation) en allemand standard, forme \all{Sie} (formel, écrit administratif). Respectez l'impératif poli, les verbes à particule séparable (\all{aufräumen}, \all{lüften}, \all{anschalten}), les pronominaux si nécessaire.
    \end{enumerate}

    \item \textbf{Relecture croisée et correction (10 min)} \\
    Échangez vos productions (tract + lettre + traductions) avec un autre groupe. Utilisez la \textbf{grille de correction} (fournie par le professeur) pour annoter : orthographe (majuscules, Umlaute, ß), grammaire (cas, conjugaison, ordre des mots, place du verbe en subordonnée), lexique (précision, variété), cohérence textuelle (connecteurs, progression). Remettez les copies annotées aux auteurs.

    \item \textbf{Finalisation et affichage (10 min)} \\
    Intégrez les corrections. Les 3 meilleurs tracts (vote de la classe ou choix du professeur) sont affichés au tableau / dans le couloir. Les lettres peuvent être numérisées et envoyées au lycée partenaire (projet eTwinning).
\end{enumerate}

\courssub{Ressources à mobiliser}
Voici un rappel synthétique des outils linguistiques nécessaires. Reportez-vous aux fiches de cours des leçons précédentes pour les tableaux complets.

\begin{definition}[Lexique clé : Environnement \& Santé au lycée]
\begin{center}
\begin{tabular}{|l|l|l|}
\hline
\rowcolor{popblueL} \textbf{Deutsch} & \textbf{Français} & \textbf{Remarque} \\
\hline
der Müll (kein Pl.) / der Abfall, die Abfälle & les déchets, les ordures & \all{der Sondermüll} déchets toxiques \\
die Mülltrennung & le tri sélectif & \all{trennen} (séparer) \\
der Mülleimer, die Mülleimer & la poubelle & \all{der Papierkorb} corbeille à papier \\
wegwerfen (wirft weg, warf weg, weggeworfen) & jeter & verbe à particule séparable \\
recyceln & recycler & verbe régulier \\
sparen & économiser & \all{Wasser sparen}, \all{Strom sparen} \\
lüften (lüftet, lüftete, gelüftet) & aérer & \all{regelmäßig lüften} \\
die Luft, die Lüfte & l'air & \all{frische Luft} \\
der Lärm & le bruit, le vacarme & \all{Lärm vermeiden} \\
die Gesundheit & la santé & \all{die Gesundheitserziehung} \\
sich waschen & se laver (les mains) & pronominal, accusatif : \all{sich die Hände waschen} \\
sich bewegen & bouger, faire du sport & pronominal \\
sich ernähren & s'alimenter & pronominal, + datif : \all{sich gesund ernähren} \\
der Wasserhahn, die Wasserhähne & le robinet & \all{den Wasserhahn zudrehen} \\
das Licht, die Lichter & la lumière & \all{das Licht ausschalten} \\
die Pflanze, die Pflanzen & la plante & \all{Pflanzen gießen} \\
\hline
\end{tabular}
\end{center}
\end{definition}

\begin{propriete}[L'impératif (Imperativ) — Révision]
\begin{center}
\begin{tabular}{|l|l|l|l|}
\hline
\rowcolor{popblueL} \textbf{Personne} & \textbf{Verbe régulier} & \textbf{Verbe à particule séparable} & \textbf{Verbe pronominal} \\
\hline
\all{du} (2e sg familier) & \all{Mach die Hausaufgaben!} & \all{Räum dein Zimmer auf!} & \all{Wasch dir die Hände!} \\
\all{ihr} (2e pl familier) & \all{Macht die Hausaufgaben!} & \all{Räumt euer Zimmer auf!} & \all{Wascht euch die Hände!} \\
\all{Sie} (politesse sg/pl) & \all{Machen Sie die Hausaufgaben!} & \all{Räumen Sie Ihr Zimmer auf!} & \all{Waschen Sie sich die Hände!} \\
\all{wir} (1e pl, exhortation) & \all{Machen wir die Hausaufgaben!} & \all{Räumen wir unser Zimmer auf!} & \all{Waschen wir uns die Hände!} \\
\hline
\end{tabular}
\end{center}
\textbf{Règles clés :}
\begin{itemize}
    \item \all{du} : radical + \textbf{-e} (souvent omis à l'oral : \all{Mach!}), pas de \all{-st}. Verbes en \all{-eln/-ern} gardent \all{-e} : \all{handel(e)!}
    \item \all{ihr} : infinitif sans \all{-en} + \textbf{-t} : \all{Macht!}
    \item \all{Sie} / \all{wir} : infinitif + pronom sujet \textbf{après} le verbe : \all{Machen Sie!}, \all{Machen wir!}
    \item Verbes à particule séparable : particule en \textbf{fin de phrase} : \all{Räum ... auf!}
    \item Verbes pronominaux : pronom réfléchi \textbf{après} le verbe (accusatif généralement) : \all{Wasch dich!}, \all{Wascht euch!}, \all{Waschen Sie sich!}
    \item Négation : \all{nicht} en fin de proposition (avant particule séparable) : \all{Werf den Müll nicht auf den Boden!}
\end{itemize}
\end{propriete}

\begin{propriete}[Verbes pronominaux (Reflexivverben) — Emplois principaux]
\begin{enumerate}
    \item \textbf{Vraiment pronominaux (sens réflexe)} : l'action revient au sujet. \all{sich waschen, sich duschen, sich verletzen, sich beeilen}. Accusatif : \all{Ich wasche mich.}
    \item \textbf{Pronominaux à sens réciproque} : action mutuelle. \all{sich treffen, sich kennenlernen, sich unterhalten, sich streiten}. \all{Wir treffen uns vor der Schule.}
    \item \textbf{Pronominaux à sens passif / autocausatif} : \all{sich lohnen} (valoir la peine), \all{sich handeln um} (il s'agit de), \all{sich ändern} (changer).
    \item \textbf{Verbes à préposition figée} : \all{sich kümmern um} (s'occuper de + acc), \all{sich freuen auf} (se réjouir de + acc), \all{sich interessieren für} (s'intéresser à + acc), \all{sich beteiligen an} (participer à + dat).
\end{enumerate}
\textbf{Conjugaison (Präsens) :} \all{ich wasche mich, du wäschst dich, er wäscht sich, wir waschen uns, ihr wascht euch, sie waschen sich}. \cle{Le pronom réfléchi suit le verbe conjugué (position 2 en phrase principale).}
\end{propriete}

\begin{propriete}[Subordonnées : Causale (weil) et Finale (damit / um ... zu)]
\begin{center}
\begin{tabularx}{\linewidth}{|X|X|X|}
\hline
\rowcolor{popblueL} \textbf{Type} & \textbf{Structure} & \textbf{Exemple} \\
\hline
\all{weil} (cause) & \all{weil} + sujet + ... + \textbf{verbe à la fin} & \all{Wir trennen Müll, weil wir die Umwelt schützen.} \\
\hline
\all{damit} (but, sujet différent possible) & \all{damit} + sujet + ... + \textbf{verbe à la fin} & \all{Wir stellen Mülleimer auf, damit alle Müll trennen können.} \\
\hline
\all{um ... zu} + infinitif (but, \textbf{même sujet}) & \all{um} + ... + \textbf{zu} + \textbf{infinitif à la fin} & \all{Wir stellen Mülleimer auf, um die Schule sauber zu halten.} \\
\hline
\end{tabularx}
\end{center}
\cle{Attention :} \all{um ... zu} ne s'utilise que si le sujet de la principale et de la subordonnée est \textbf{identique}. Sinon $\to$ \all{damit}. Le verbe à la fin en subordonnée est \textbf{la règle d'or} (verbe conjugué pour \all{weil/damit}, infinitif pour \all{zu}).
\end{propriete}

\begin{attention}[Pièges fréquents pour francophones]
\begin{itemize}
    \item \textbf{Majuscules :} \emph{Tous} les noms en allemand (\all{der Müll, die Gesundheit, das Wasser}). Pas de majuscule aux adjectifs de nationalité/langue (\all{deutsch, französisch}) ni au pronom \all{ich}.
    \item \textbf{Ordre des mots (V2) :} En phrase principale, le verbe conjugué est \textbf{toujours en 2e position}. \all{Heute lüften wir die Klasse.} (pas \emph{\all{Wir lüften heute die Klasse.}} si \all{Heute} est en tête).
    \item \textbf{Verbe en fin de subordonnée :} \all{..., weil wir die Klasse lüften.} (pas \emph{\all{... weil wir lüften die Klasse.}}).
    \item \textbf{Particule séparable :} \all{aufräumen $\to$ Räum auf! / Wir räumen auf.} Particule à la fin.
    \item \textbf{Pronom réfléchi :} \all{sich waschen $\to$ Ich wasche mich.} Ne pas oublier le pronom, ni le mettre au mauvais cas (généralement accusatif pour le corps, datif si partie du corps + possesseur : \all{Ich wasche mir die Hände}).
    \item \textbf{Faux amis :} \all{Gift} = poison (pas \emph{gift/cadeau} $\to$ \all{das Geschenk}), \all{aktuell} = actuel/en cours (pas \emph{actuel} au sens de \emph{réel} $\to$ \all{tatsächlich}), \all{konsequent} = cohérent/conséquent (pas \emph{conséquent} $\to$ \all{folglich}), \all{sensibel} = sensible (pas \emph{sensible} $\to$ \all{vernünftig}).
    \item \textbf{Thème (FR $\to$ DE) :} \emph{« Merci d'encourager... »} $\not\to$ *\all{Danke ermutigen...} mais \all{Bitten Sie die Schüler, ... zu ermutigen} ou \all{Ermutigen Sie die Schüler, ...}. \emph{« Économiser l'eau »} $\to$ \all{Wasser sparen} (infinitif nominalisé) ou \all{sparen Sie Wasser} (impératif).
\end{itemize}
\end{attention}

\begin{exemplebox}[Modelltext : Tract type (niveau A2) — Analyse en classe]
\textbf{„Saubere Schule – Gesunde Zukunft!“} \\
\textbf{Unsere Schule soll sauber und gesund bleiben. Jeder kann helfen!} \\
\textbf{1. Wirf deinen Müll in den Mülleimer!} \\
\textbf{2. Trenne Papier, Plastik und Bioabfall!} \\
\textbf{3. Dreh den Wasserhahn zu, wenn du dich wäschst!} \\
\textbf{4. Lüfte das Klassenzimmer in der Pause!} \\
\textbf{5. Schalte das Licht aus, wenn niemand drin ist!} \\
\textbf{Wir machen das, \cle{weil} wir unsere Umwelt schützen wollen.} \\
\textbf{Wir trennen den Müll, \cle{damit} er recycelt werden kann.} \\
\textbf{\cle{Lasst uns gemeinsam} für eine saubere Schule sorgen!}
\end{exemplebox}

\courssub{Proposition de production et corrigé commenté}
\begin{exoresolu}[Production modèle du groupe + commentaire didactique]
\textbf{A. Tract finalisé (version affichage)}

„\textbf{Gesund und sauber an unserer Schule – Mach mit!}“ \\
\textbf{Eine saubere Schule macht gesund und glücklich. Jeder zählt!} \\
\textbf{1. Wirf deinen Müll in den richtigen Eimer!} \\
\textbf{2. Trenne deinen Müll: Papier, Plastik, Biomüll!} \\
\textbf{3. Wasch dir die Hände gründlich mit Seife!} \\
\textbf{4. Dreh den Wasserhahn zu, während du dich einseifst!} \\
\textbf{5. Lüftet die Klassenzimmer in jeder großen Pause!} \\
\textbf{Wir trennen den Müll, \cle{weil} wir die Natur schützen und Rohstoffe sparen wollen.} \\
\textbf{Wir lüften regelmäßig, \cle{damit} die Luft frisch bleibt und Viren keine Chance haben.} \\
\textbf{\cle{Werden Sie mit!} Gemeinsam schaffen wir eine gesunde Schule!}

\textbf{B. Lettre type (élève « Aminatou » à « Lukas »)}

Yaoundé, den 12. Oktober 2025 \\
Lieber Lukas, \\
wie geht es dir? Hier in Yaoundé starten wir nächste Woche die Aktion „Gesund und sauber“. \cle{Weil} unser Schulhof oft schmutzig ist, wollen wir etwas ändern. \cle{Ich engagiere mich} in der Gruppe für die Plakate und \cle{kümmere mich um} die Übersetzung ins Französische. \cle{Wir treffen uns} jeden Mittag, um die Mülleimer zu kontrollieren. \cle{Es wäre toll, wenn} ihr in Stuttgart auch eine Aktion macht. Schreib mir bitte, was ihr plant! \\
Viele Grüße aus Kamerun \\
Deine Aminatou

\textbf{C. Version (Tract $\to$ Français) — Proposition corrigée}
„\textbf{École saine et propre – Participez !}“ \\
\textbf{Une école propre rend en bonne santé et heureux. Chacun compte !} \\
\textbf{1. Jetez vos déchets dans la bonne poubelle !} \\
\textbf{2. Triez vos déchets : papier, plastique, biodéchets !} \\
\textbf{3. Lavez-vous bien les mains avec du savon !} \\
\textbf{4. Fermez le robinet pendant que vous vous savonnez !} \\
\textbf{5. Aérez les salles de classe à chaque grande récréation !} \\
\textbf{Nous trions les déchets, \cle{car} nous voulons protéger la nature et économiser les ressources.} \\
\textbf{Nous aérons régulièrement, \cle{afin que} l'air reste frais et que les virus n'aient aucune chance.} \\
\textbf{\cle{Rejoignez-nous !} Ensemble, nous créons une école saine !}

\textbf{D. Thème (Message proviseur FR $\to$ DE) — Proposition corrigée}
„\textbf{Liebe Kolleginnen und Kollegen,} \\
\textbf{die Aktion „Gesund und sauber“ beginnt nächsten Montag. \cle{Bitten Sie} die Schülerinnen und Schüler, den Müll in die Mülleimer zu werfen, an den Wasserhähnen Wasser zu sparen und die Klassenzimmer regelmäßig zu lüften. \cle{Ihr Engagement} ist wichtig für den Erfolg des Projekts.} \\
\textbf{Mit freundlichen Grüßen} \\
\textbf{Der Schulleiter}“

\tcblower
\textbf{Commentaire didactique des choix de langue}

\begin{enumerate}
    \item \textbf{Tract — Impératif :} Choix de la forme \all{du} (familier singulier) pour les conseils 1--4 (\all{Wirf, Trenne, Wasch, Dreh}) car le tract s'adresse \emph{directement} à chaque élève individuellement (« \emph{Toi, élève, fais ceci} »). Conseil 5 (\all{Lüftet}) à la forme \all{ihr} (pluriel familier) car l'aération est une action \emph{collective} de la classe. Appel final \all{Werden Sie mit!} (politesse \all{Sie}) pour inclure le personnel et les visiteurs adultes. \cle{Alternance justifiée selon la cible de l'injonction.}
    \item \textbf{Verbes pronominaux :} \all{Wasch dir die Hände} (accusatif \all{dir} = pour toi / à toi, corps = \all{die Hände} accusatif objet) — structure possessive typique allemande (pas *\all{tes mains} mais \all{dir die Hände}). \all{während du dich einseifst} (pronominal \all{sich einseifen}, subordonnée temporelle \all{während} + verbe à la fin).
    \item \textbf{Subordonnées :} \all{weil} (cause) verbe à la fin (\all{wollen}). \all{damit} (but) verbe à la fin (\all{haben}). \cle{Respect strict de la place du verbe.} \all{um ... zu} évité car sujet différent (\all{Wir lüften} / \all{die Luft bleibt}).
    \item \textbf{Lettre :} En-tête standard (lieu, \all{den} + date). \all{Lieber Lukas} (masc. accusatif). \all{Weil}-phrase en position 1 (topicalisation) $\to$ verbe principal \cle{starten} en position 2. \all{Ich engagiere mich} (pronominal vrai réflexif), \all{kümmere mich um} (préposition figée + acc), \all{treffen uns} (réciproque). \all{Es wäre toll, wenn} (Konjunktiv II \all{wäre} + \all{wenn} pour politesse/souhait).
    \item \textbf{Version :} Adaptation culturelle : \all{Biomüll} $\to$ \emph{biodéchets} (terme technique français). \all{während du dich einseifst} $\to$ \emph{pendant que vous vous savonnez} (registre \emph{vous} pour affichage public). \all{Viren keine Chance haben} $\to$ \emph{les virus n'aient aucune chance} (subjonctif français pour but/négation).
    \item \textbf{Thème :} Registre formel \all{Sie} imposé. \all{Bitten Sie die Schülerinnen und Schüler, ... zu werfen / zu sparen / zu lüften} (infinitif avec \all{zu} après verbe de volonté/ordre). \all{Ihr Engagement} (nominalisation de \emph{votre engagement}) au lieu de \all{Sie engagieren sich} (plus administratif). \all{Mit freundlichen Grüßen} (formule standard lettres formelles).
\end{enumerate}
\end{exoresolu}

\courssub{Bilan de l'activité}
\begin{aretenir}[Ce que vous devez retenir de cette activité d'intégration]
\begin{itemize}
    \item La \textbf{production écrite} en allemand (tract, lettre) suit des \textbf{conventions de genre} précises : structure, registre, formules types. Le respect de ces codes assure la lisibilité et l'efficacité communicationnelle.
    \item L'\textbf{impératif} est l'outil grammatical central du \emph{conseil} et de l'\emph{injonction} : maîtriser ses 4 formes (\all{du, ihr, Sie, wir}) et la place de la particule / du pronom réfléchi est indispensable.
    \item Les \textbf{verbes pronominaux} ne sont pas de simples « verbes réfléchis » : ils expriment la réciprocité, l'auto-causation, ou s'utilisent avec des prépositions figées (\all{sich kümmern um}). Le cas du pronom (accusatif / datif) dépend de la structure.
    \item Les \textbf{subordonnées} (\all{weil}, \all{damit}, \all{um ... zu}) permettent de \textbf{justifier} (cause) ou de \textbf{finaliser} (but) une action. La \cle{règle du verbe en position finale} est non négociable.
    \item La \textbf{traduction} (thème/version) n'est pas un mot-à-mot : elle exige une \textbf{reformulation} dans la langue cible (ordre des mots, choix du registre, faux-amis, structures idiomatiques). La relecture croisée est l'étape clé pour progresser.
    \item Le \textbf{travail collaboratif} (rédaction, correction par les pairs, affichage) développe l'autonomie, l'esprit critique et la fierté linguistique. Vos tracts bilingues sont de vrais actes de communication dans l'espace scolaire !
\end{itemize}
\end{aretenir}

\begin{experience}[Prolongement possible : Projet eTwinning / Jumelage]
Si le lycée partenaire de Stuttgart répond favorablement, organisez une \textbf{visioconférence} (30 min) où chaque classe présente sa campagne, ses résultats (photos, poids des déchets collectés, sondage « air frais »...) et échange en allemand sur les différences / similitudes (ex. : \all{Pfandsystem} en Allemagne vs gestion des déchets à Yaoundé). Cela valide la compétence \emph{Expression orale en interaction} et ancre l'apprentissage dans une réalité interculturelle vivante.
\end{experience}

\newpage

\coursec{Méthodes et exercices résolus}

\courssub{Méthodes et exercices résolus}

\begin{methode}[Analyser et rédiger un tract (Flyer) pour la santé et l'environnement]
\textbf{Objectif :} Produire un texte court, incitatif et visuel (niveau A2) pour sensibiliser à un problème écologique ou sanitaire.
\begin{enumerate}
    \item \textbf{Analyser la consigne :} Identifier le \cle{public cible} (jeunes, familles, passants), le \cle{thème} (déchets, eau, nutrition, sport) et le \cle{format} (titre accrocheur, corps du texte, slogan final, coordonnées).
    \item \textbf{Structurer le tract :} 
    \begin{itemize}
        \item \textbf{Titre :} Court, impératif ou question (ex. \all{Schützt unser Wasser!} / \all{Wollen Sie gesund leben?}).
        \item \textbf{Arguments (3--4 points) :} Phrases courtes, verbes d'action (impératif), vocabulaire précis (\all{Müll trennen}, \all{Leitungswasser trinken}, \all{Fahrrad fahren}).
        \item \textbf{Slogan / Appel final :} Mémorisable, positif (ex. \all{Mach mit!}, \all{Gemeinsam für die Zukunft!}).
        \item \textbf{Infos pratiques :} Lieu, date, site web, QR-code.
    \end{itemize}
    \item \textbf{Choisir la grammaire :} \cle{Impératif} (2e pers. sg/pl, forme polie), \cle{verbes pronominaux} (\all{sich bewegen}, \all{sich ernähren}), \cle{subordonnées causales} (\all{weil}), \cle{conditionnelles} (\all{wenn}).
    \item \textbf{Rédiger en allemand :} Respecter la \cle{majuscule aux noms}, l'\cle{ordre des mots} (verbe en 2e position en principal, à la fin en subordonnée), l'\cle{orthographe} (Umlaute, ß).
    \item \textbf{Relire (checklist) :} Majuscules ? Verbes conjugués ? Cas (Akk/Dat après prépositions) ? Vocabulaire du chapitre ? Ponctuation allemande („ “) ?
\end{enumerate}
\end{methode}

\begin{exoresolu}[Rédiger un tract : « Une école propre et verte »]
\textbf{Énoncé :} Votre lycée organise une semaine « École verte ». Rédigez un tract en allemand (environ 80--100 mots) pour inciter les élèves à : trier les déchets, boire l'eau du robinet, venir à vélo ou à pied, planter un arbre. Utilisez l'impératif (forme \textit{Sie} ou \textit{ihr}), au moins une subordonnée avec \textit{weil} et un verbe pronominal.
\tcblower
\textbf{Solution détaillée :}
\begin{enumerate}
    \item \textbf{Planification :} Titre $\rightarrow$ 4 actions concrètes $\rightarrow$ Slogan $\rightarrow$ Contact.
    \item \textbf{Choix de la forme d'adresse :} \cle{Sie-Form} (polie, standard pour affichage public) ou \cle{ihr-Form} (familière, entre élèves). Ici : \textbf{Sie-Form} (plus formel pour un tract officiel).
    \item \textbf{Rédaction (brouillon mental) :}
    \begin{itemize}
        \item Titre: \all{Grüne Schule – Wir machen mit!}
        \item Point 1 (Déchets): \all{Trennen Sie Ihren Müll! Papier, Plastik und Bio gehören in verschiedene Tonnen.} (Impératif Sie, vocabulaire : \all{der Müll}, \all{die Tonne}, \all{trennen}).
        \item Point 2 (Eau): \all{Trinken Sie Leitungswasser! Das spart Plastikflaschen und Geld, weil unser Wasser sauber ist.} (Subordonnée \all{weil} $\rightarrow$ verbe à la fin \all{ist}).
        \item Point 3 (Transport): \all{Kommen Sie mit dem Fahrrad oder zu Fuß zur Schule! So schützen Sie das Klima und bewegen Sie sich gesund.} (Verbe pronominal \all{sich bewegen} $\rightarrow$ \all{bewegen Sie sich}).
        \item Point 4 (Arbre): \all{Pflanzen Sie einen Baum im Schulhof!}
        \item Slogan: \all{Gemeinsam für eine saubere Zukunft!}
        \item Contact: \all{Info: www.lycee-vert.cm}
    \end{itemize}
    \item \textbf{Texte final (allemand) :}
\end{enumerate}
\begin{textebox}[Tract : Grüne Schule]
\all{Grüne Schule – Wir machen mit!

Trennen Sie Ihren Müll! Papier, Plastik und Bio gehören in verschiedene Tonnen. Trinken Sie Leitungswasser! Das spart Plastikflaschen und Geld, weil unser Wasser sauber ist. Kommen Sie mit dem Fahrrad oder zu Fuß zur Schule! So schützen Sie das Klima und bewegen Sie sich gesund. Pflanzen Sie einen Baum im Schulhof!

Gemeinsam für eine saubere Zukunft!

Info: www.lycee-vert.cm}
\end{textebox}
\textbf{Traduction française :} « École verte – Nous participons ! Triez vos déchets ! Papier, plastique et bio vont dans des poubelles différentes. Buvez l'eau du robinet ! Cela économise des bouteilles en plastique et de l'argent, car notre eau est propre. Venez à l'école à vélo ou à pied ! Ainsi vous protégez le climat et vous bougez sainement. Plantez un arbre dans la cour de récréation ! Ensemble pour un avenir propre ! Info : www.lycee-vert.cm »
\textbf{Justifications grammaticales :} \all{Trennen/Trinken/Kommen/Pflanzen} = Impératif \textit{Sie} (infinitif + \textit{Sie}). \all{weil unser Wasser sauber \textbf{ist}} = subordonnée causale, verbe à la fin. \all{bewegen Sie \textbf{sich}} = impératif pronominal (accusatif). Majuscules sur tous les noms (\all{Müll, Tonnen, Leitungswasser, Fahrrad, Klima, Baum, Schulhof, Zukunft}). Guillemets non nécessaires ici (texte affiché).
\end{exoresolu}

\begin{methode}[Formuler des conseils avec l'impératif et les verbes pronominaux]
\textbf{Objectif :} Maîtriser les trois formes de l'impératif et l'emploi des verbes pronominaux (réfléchis / réciproques) pour donner des conseils de santé (\cle{Rat geben}).
\begin{enumerate}
    \item \textbf{Identifier le destinataire :} 
    \begin{itemize}
        \item \cle{du} (ami, famille) $\rightarrow$ radical du présent sans \textit{-e} (souvent) : \all{komm!}, \all{trink!}, \all{lies!} (verbes forts : voyelle modifiée \all{gib!}, \all{sprich!}, \all{iss!}).
        \item \cle{ihr} (groupe d'amis, classe) $\rightarrow$ forme du présent \textit{ihr} sans pronom : \all{kommt!}, \all{trinkt!}, \all{liest!}.
        \item \cle{Sie} (formel, public, adultes) $\rightarrow$ infinitif + \textit{Sie} : \all{kommen Sie!}, \all{trinken Sie!}, \all{lesen Sie!}.
    \end{itemize}
    \item \textbf{Verbes pronominaux (ex. \all{sich gesund ernähren}, \all{sich bewegen}, \all{sich entspannen}, \all{sich waschen}) :}
    \begin{itemize}
        \item Le pronom réfléchi \textbf{suit le verbe} à l'impératif.
        \item \textit{du} : \all{Ernähr dich gesund!} \quad \textit{ihr} : \all{Ernährt euch gesund!} \quad \textit{Sie} : \all{Ernähren Sie sich gesund!}
    \end{itemize}
    \item \textbf{Négation :} \all{Nicht} placé \textbf{après} le verbe (et après le pronom réfléchi le cas échéant) : \all{Rauchen Sie nicht!}, \all{Setz dich nicht in die Sonne!}.
    \item \textbf{Structures types pour conseiller :} \all{Man sollte ...} (on devrait), \all{Es ist wichtig, ... zu + Infinitiv}, \all{Ich rate Ihnen / dir, ... zu + Infinitiv}.
\end{enumerate}
\end{methode}

\begin{exoresolu}[Transformer des conseils : formes d'adresse et pronominaux]
\textbf{Énoncé :} Mettez les phrases suivantes à la forme d'adresse indiquée entre parenthèses. Traduisez en français.
\begin{enumerate}
    \item \all{Iss mehr Gemüse!} (Sie-Form)
    \item \all{Wascht euch die Hände vor dem Essen!} (du-Form)
    \item \all{Man sollte jeden Tag Sport treiben.} (Impératif ihr-Form, verbe pronominal \all{sich bewegen})
    \item \all{Trinken Sie keinen Alkohol!} (du-Form, pronominal \all{sich betrinken} -- sens figuré/argotique, ici on utilise \all{trinken} simple ou \all{verzichten auf} + Akk). \textbf{Correction pédagogique :} Utilisons \all{Verzichten Sie auf Alkohol!} $\rightarrow$ du-Form.
\end{enumerate}
\tcblower
\textbf{Solution détaillée :}
\begin{enumerate}
    \item \all{Iss mehr Gemüse!} $\rightarrow$ \textbf{\all{Essen Sie mehr Gemüse!}} (Infinitif \all{essen} + \all{Sie}). \textit{Traduction :} « Mangez plus de légumes ! » (Forme polie). \textbf{Attention :} Verbe fort \all{essen} (du isst), impératif \textit{du} = \all{iss} (pas \all{isse}).
    \item \all{Wascht euch die Hände ...} $\rightarrow$ \textbf{\all{Wasch dir die Hände vor dem Essen!}} (Radical \all{wasch} + pronom \all{dir} datif \textit{dir} car \all{die Hände} est l'objet direct / Akkusativ). \textit{Traduction :} « Lave-toi les mains avant de manger ! » \textbf{Grammaire :} \all{sich die Hände waschen} (Akkusativ \all{die Hände}, Datif \all{sich/dir}).
    \item \all{Man sollte ... Sport treiben} $\rightarrow$ \textbf{\all{Bewegt euch jeden Tag!}} (Impératif \textit{ihr} de \all{sich bewegen} : \all{bewegt} + \all{euch}). \textit{Traduction :} « Bougez-vous chaque jour ! » \textbf{Alternative :} \all{Treibt jeden Tag Sport!} (sans pronominal).
    \item \all{Verzichten Sie auf Alkohol!} $\rightarrow$ \textbf{\all{Verzichte auf Alkohol!}} (Impératif \textit{du} de \all{verzichten} : radical \all{verzicht} + \textit{e} $\rightarrow$ \all{verzichte} ; préposition \all{auf} + Akkusativ). \textit{Traduction :} « Renonce à l'alcool ! » \textbf{Note :} \all{sich betrinken} = « se saouler », pas « boire de l'alcool » au sens conseil santé.
\end{enumerate}
\textbf{Récapitulatif des formes impératives utilisées :}
\begin{center}
\begin{tabularx}{\linewidth}{|X|X|X|X|}
\hline
\rowcolor{popblueL} \textbf{Infinitif} & \textbf{du} & \textbf{ihr} & \textbf{Sie} \\ \hline
\all{essen} & \all{iss!} & \all{esst!} & \all{essen Sie!} \\ \hline
\all{sich waschen} (Dat) & \all{wasch dir ...!} & \all{wascht euch ...!} & \all{waschen Sie sich ...!} \\ \hline
\all{sich bewegen} (Akk) & \all{beweg dich!} & \all{bewegt euch!} & \all{bewegen Sie sich!} \\ \hline
\all{verzichten auf} + Akk & \all{verzichte auf ...!} & \all{verzichtet auf ...!} & \all{verzichten Sie auf ...!} \\ \hline
\end{tabularx}
\end{center}
\end{exoresolu}

\begin{methode}[Utiliser les subordonnées (weil, dass, wenn) pour argumenter]
\textbf{Objectif :} Construire des phrases complexes correctes (ordre des mots : verbe à la fin) pour justifier un conseil ou exprimer une conséquence (niveau A2/B1).
\begin{enumerate}
    \item \textbf{Conjonction \cle{weil} (parce que) :} Introduit la cause. \textbf{Verbe à la fin.}
    \all{Ich fahre mit dem Rad, \textbf{weil es gut für die Umwelt \textbf{ist}}.}
    \item \textbf{Conjonction \cle{dass} (que) :} Introduit un contenu (souvent après \all{ich denke}, \all{ich finde}, \all{es ist wichtig}). \textbf{Verbe à la fin.}
    \all{Ich finde, \textbf{dass Mülltrennung \textbf{wichtig} \textbf{ist}}.}
    \item \textbf{Conjonction \cle{wenn} (si / quand) :} Condition ou répétition. \textbf{Verbe à la fin.}
    \all{\textbf{Wenn} du Wasser sparst, \textbf{hilfst} du der Umwelt.} (Virgule avant la principale, verbe en 2e pos. \all{hilfst}).
    \item \textbf{Conjonction \cle{obwohl} (bien que) :} Concession. \textbf{Verbe à la fin.}
    \all{Er raucht, \textbf{obwohl er krank \textbf{ist}}.}
    \item \textbf{Piège francophone :} Ne pas mettre le verbe en 2e position dans la subordonnée (\all{*weil es ist gut} $\rightarrow$ \textbf{faux}). Ne pas oublier la virgule avant la conjonction.
\end{enumerate}
\end{methode}

\begin{exoresolu}[Compléter et traduire : subordonnées et ordre des mots]
\textbf{Énoncé :} Complétez les phrases par la conjonction adaptée (\all{weil}, \all{dass}, \all{wenn}, \all{obwohl}) et mettez le verbe entre parenthèses à la forme correcte (Présent). Traduisez.
\begin{enumerate}
    \item \all{Der Arzt sagt, \_\_\_\_\_\_ Rauchen \_\_\_\_\_\_ (sein) sehr gefährlich.}
    \item \all{\_\_\_\_\_\_\_ du Obst isst, \_\_\_\_\_\_ (haben) du mehr Vitamine.}
    \item \all{Wir trinken Leitungswasser, \_\_\_\_\_\_ es \_\_\_\_\_\_ (schmecken) gut und \_\_\_\_\_\_ (sparen) Geld.}
    \item \all{Viele Jugendliche bewegen sich zu wenig, \_\_\_\_\_\_ sie \_\_\_\_\_\_ (wissen), dass Sport gesund \_\_\_\_\_\_ (sein).}
\end{enumerate}
\tcblower
\textbf{Solution détaillée :}
\begin{enumerate}
    \item \all{Der Arzt sagt, \textbf{dass} Rauchen \textbf{ist} sehr gefährlich.} 
    \textit{Analyse :} \all{sagen} + subordonnée objet $\rightarrow$ \all{dass}. Verbe \all{sein} à la fin. \textit{Traduction :} « Le médecin dit que fumer est très dangereux. » (\all{Rauchen} = nom verbal, neutre, sujet).
    \item \all{\textbf{Wenn} du Obst isst, \textbf{hast} du mehr Vitamine.} 
    \textit{Analyse :} Condition $\rightarrow$ \all{wenn} + verbe à la fin (\all{isst}). Principale : verbe en 2e position (\all{hast}). \textbf{Attention :} \all{haben} (du hast) $\rightarrow$ \all{hast} (pas \all{hasst}). \textit{Traduction :} « Si tu manges des fruits, tu as plus de vitamines. »
    \item \all{Wir trinken Leitungswasser, \textbf{weil} es \textbf{gut schmeckt} und \textbf{Geld spart}.} 
    \textit{Analyse :} Cause $\rightarrow$ \all{weil}. Deux verbes coordonnés dans la subordonnée : \all{schmeckt} (es) et \all{spart} (es, sujet sous-entendu). Verbes en position finale (ordre : complément \all{gut} avant verbe \all{schmeckt}). \textit{Traduction :} « Nous buvons l'eau du robinet, car elle a bon goût et économise de l'argent. »
    \item \all{Viele Jugendliche bewegen sich zu wenig, \textbf{obwohl} sie \textbf{wissen}, \textbf{dass} Sport gesund \textbf{ist}.} 
    \textit{Analyse :} Concession $\rightarrow$ \all{obwohl} + verbe à la fin (\all{wissen}). Subordonnée imbriquée après \all{wissen} $\rightarrow$ \all{dass} + verbe à la fin (\all{ist}). \textit{Traduction :} « Beaucoup de jeunes bougent trop peu, bien qu'ils sachent que le sport est sain. » \textbf{Point clé :} \all{sich bewegen} (pronominal) $\rightarrow$ \all{bewegen sich}. \all{wissen} (savoir fait) vs \all{kennen} (connaître).
\end{enumerate}
\end{exoresolu}

\begin{methode}[Traduire un texte de l'allemand vers le français (Version)]
\textbf{Objectif :} Restituer le sens exact d'un texte allemand (niveau A2) en français naturel, en gérant les faux amis, la structure nominale allemande et les temps (Perfekt/Präteritum).
\begin{enumerate}
    \item \textbf{Lecture globale :} Comprendre le thème (environnement/santé), le ton (informatif, conseil), la structure (paragraphes).
    \item \textbf{Analyse phrase par phrase :}
    \begin{itemize}
        \item Repérer le \cle{verbe} (souvent en 2e pos ou à la fin) et son \cle{sujet} (majuscule !).
        \item Identifier les \cle{cas} (Nom/Acc/Dat/Gen) pour savoir « qui fait quoi à qui ».
        \item Détecter les \cle{verbes à particule séparable} (\all{aufstehen, mitbringen, vorhaben}) $\rightarrow$ particule en fin de phrase.
        \item Gérer le \cle{Perfekt} (\all{hat gemacht} = a fait / a fait) vs \cle{Präteritum} (\all{machte} = faisait / fit) : au choix selon le contexte narratif ou rapporté.
    \end{itemize}
    \item \textbf{Transposition lexicale (Faux amis fréquents) :}
    \begin{center}
    \begin{tabularx}{\linewidth}{|X|X|X|}
    \hline
    \rowcolor{popblueL} \textbf{Allemand} & \textbf{Piège (Français)} & \textbf{Vrai sens} \\ \hline
    \all{aktuell} & actuel & \textbf{actuel / en ce moment / récent} \\ \hline
    \all{eventuell} & éventuellement & \textbf{éventuel / possible / au cas où} \\ \hline
    \all{konsequent} & conséquent & \textbf{conséquent / logique / ferme} \\ \hline
    \all{sensibel} & sensible & \textbf{sensible / délicat / réactif} \\ \hline
    \all{das Gift} & le gift (cadeau) & \textbf{le poison} \\ \hline
    \all{das Rezept} & la recette & \textbf{l'ordonnance (médecin) / la recette (cuisine)} \\ \hline
    \end{tabularx}
    \end{center}
    \item \textbf{Rédaction française :} Éviter le calque syntaxique (ex. \all{Es gibt} $\neq$ « Il donne » mais « Il y a »). Privilégier le style naturel (passif, impersonnel, verbes d'état).
    \item \textbf{Relecture :} Cohérence, temps, accords, registres de langue.
\end{enumerate}
\end{methode}

\begin{exoresolu}[Version : « Gesund leben in der Stadt »]
\textbf{Énoncé :} Traduisez en français le texte suivant.
\begin{textebox}[Texte source]
\all{In großen Städten wie Berlin oder Yaoundé ist die Luft oft schlecht. Viele Autos produzieren Giftstoffe. Wer jeden Tag zur Arbeit fährt, atmet diese Schadstoffe ein. Das kann zu Atemwegserkrankungen führen. Deshalb raten Ärzte: „Bewegen Sie sich in Parks oder Wäldern! Dort ist die Luft besser.“ Auch die Ernährung spielt eine Rolle. Wer sich gesund ernährt, stärkt sein Immunsystem. Ein Rezept für ein langes Leben? Genug schlafen, wenig Stress, viel Wasser trinken und nicht rauchen.}
\end{textebox}
\tcblower
\textbf{Solution détaillée (analyse et traduction) :}
\begin{enumerate}
    \item \all{In großen Städten wie Berlin oder Yaoundé ist die Luft oft schlecht.}
    $\rightarrow$ « Dans les grandes villes comme Berlin ou Yaoundé, l'air est souvent mauvais. » (\all{Städte} pl. fém, \all{die Luft} fém sg, \all{schlecht} adj prédicatif).
    \item \all{Viele Autos produzieren Giftstoffe.}
    $\rightarrow$ « Beaucoup de voitures produisent des substances toxiques / des polluants. » (\all{Giftstoffe} = \all{Gift} poison + \all{Stoffe} substances. Faux ami \all{Gift} $\neq$ cadeau).
    \item \all{Wer jeden Tag zur Arbeit fährt, atmet diese Schadstoffe ein.}
    $\rightarrow$ « Celui qui / Quiconque va au travail chaque jour (conduit), inhale ces polluants. » (\all{Wer} = pronom relatif « celui qui », sujet. \all{fährt} verbe à la fin subordonnée \all{wer...}. \all{einatmen} séparable $\rightarrow$ \all{atmet ... ein}. \all{Schadstoffe} = substances nocives).
    \item \all{Das kann zu Atemwegserkrankungen führen.}
    $\rightarrow$ « Cela peut mener à / entraîner des maladies respiratoires. » (\all{führen zu} + Dat. \all{Atemwegserkrankungen} composé : \all{Atemwege} + \all{Erkrankungen}).
    \item \all{Deshalb raten Ärzte: „Bewegen Sie sich in Parks oder Wäldern!}
    $\rightarrow$ « C'est pourquoi les médecins conseillent : "Bougez-vous / Faites de l'exercice dans les parcs ou les forêts !" » (\all{raten} + Dat \textit{jemandem} + deux-points + discours direct. \all{sich bewegen} pronominal impératif \textit{Sie}).
    \item \all{Dort ist die Luft besser.“}
    $\rightarrow$ « Là-bas, l'air est meilleur. » (\all{besser} comparatif de \all{gut}).
    \item \all{Auch die Ernährung spielt eine Rolle.}
    $\rightarrow$ « L'alimentation aussi joue un rôle. / L'alimentation entre aussi en ligne de compte. »
    \item \all{Wer sich gesund ernährt, stärkt sein Immunsystem.}
    $\rightarrow$ « Qui se nourrit sainement / Celui qui mange sainement renforce son système immunitaire. » (\all{sich ernähren} pronominal, \all{stärken} renforce, \all{sein} possessif accordé à \all{Immunsystem} n. sg).
    \item \all{Ein Rezept für ein langes Leben?}
    $\rightarrow$ « Une recette / Une ordonnance pour une longue vie ? » (\all{Rezept} : ici sens figuré « recette », mais attention au sens médical « ordonnance »).
    \item \all{Genug schlafen, wenig Stress, viel Wasser trinken und nicht rauchen.}
    $\rightarrow$ « Bien dormir / Assez de sommeil, peu de stress, boire beaucoup d'eau et ne pas fumer. » (Énumération d'infinitifs nominaux ou impératifs implicites).
\end{enumerate}
\textbf{Traduction finale proposée :}
« Dans les grandes villes comme Berlin ou Yaoundé, l'air est souvent mauvais. De nombreuses voitures produisent des polluants. Quiconque se rend au travail chaque jour inhale ces substances nocives. Cela peut entraîner des maladies respiratoires. C'est pourquoi les médecins conseillent : "Faites de l'exercice dans les parcs ou les forêts ! Là-bas, l'air est meilleur." L'alimentation joue également un rôle. Qui se nourrit sainement renforce son système immunitaire. Une recette pour une longue vie ? Assez dormir, peu de stress, boire beaucoup d'eau et ne pas fumer. »
\end{exoresolu}

\begin{methode}[Traduire un texte du français vers l'allemand (Thème)]
\textbf{Objectif :} Produire un allemand correct (A2) à partir d'un énoncé français, en appliquant les règles de grammaire (cas, ordre des mots, conjugaison, majuscules) et le vocabulaire du chapitre.
\begin{enumerate}
    \item \textbf{Découpage en unités de sens :} Phrase par phrase, identifier Sujet -- Verbe -- Compléments.
    \item \textbf{Transfert grammatical (Checklist obligatoire) :}
    \begin{itemize}
        \item \textbf{Majuscules :} \textbf{Tous les noms} (\all{der Baum}, \all{die Gesundheit}, \all{das Wasser}).
        \item \textbf{Cas :} Sujet = Nominatif ; Objet direct = Accusatif (après \all{haben, sehen, trennen, pflanzen}) ; Objet indirect / après prépositions \all{mit, von, zu, bei} = Datif ; \all{wegen, während} = Génitif (niveau B1, éviter en A2).
        \item \textbf{Ordre des mots (Topologie) :}
        \begin{itemize}
            \item Phrase principale : Verbe conjugué en \textbf{2e position}. (Sujet ou Complément en 1ère).
            \item Subordonnée (\all{weil, dass, wenn}) : Verbe conjugué à la \textbf{fin}.
            \item Verbe à particule séparable : Particule à la \textbf{fin} (principale) ou collée au verbe (fin subordonnée).
            \item Infinitif (modaux, futur, \all{zu} + Inf) : À la \textbf{fin}.
        \end{itemize}
        \item \textbf{Verbes pronominaux :} Pronom accordé (Akk ou Dat) placé \textbf{après le verbe conjugué} (principale) ou \textbf{avant le verbe à la fin} (subordonnée).
        \item \textbf{Genre des noms :} Apprendre \textbf{toujours} avec l'article (\all{der/die/das}). Pluriel imprévisible $\rightarrow$ à mémoriser.
    \end{itemize}
    \item \textbf{Relecture croisée :} Relire l'allemand en se demandant : « Est-ce que \textit{der/die/das} est juste ? Le verbe est-il à la bonne place ? L'Umlaut/ß est-il là ? »
\end{enumerate}
\end{methode}

\begin{exoresolu}[Thème : Conseils pour un ami (Lettre / Message)]
\textbf{Énoncé :} Traduisez en allemand. Utilisez la forme \textit{du}. Contexte : Tu écris à ton ami Paul qui est stressé et mange mal.
\begin{enumerate}
    \item « Cher Paul, » 
    \item « Tu travailles trop et tu dors trop peu. »
    \item « Tu devrais manger plus de légumes et boire de l'eau, \textbf{parce que} c'est bon pour la santé. »
    \item « \textbf{Si} tu fais du sport chaque jour, tu te sens mieux. »
    \item « Ne fume pas et ne bois pas d'alcool ! »
    \item « Écris-moi bientôt ! Ton ami, Marc »
\end{enumerate}
\tcblower
\textbf{Solution détaillée (phrase par phrase) :}
\begin{enumerate}
    \item \all{Lieber Paul,} (Vocatif : forme = Nominatif masculin singulier \all{der} $\rightarrow$ \all{lieber}).
    \item \all{Du arbeitest zu viel und schläfst zu wenig.} 
    \textbf{Grammaire :} \all{arbeiten} (reg.) $\rightarrow$ \all{arbeitest}. \all{schlafen} (fort : a $\rightarrow$ ä) $\rightarrow$ \all{schläfst}. Coordination \all{und} : même structure (verbe 2e pos). \all{zu viel / zu wenig} (adverbes).
    \item \all{Du solltest mehr Gemüse essen und Wasser trinken, \textbf{weil} das gut für die Gesundheit \textbf{ist}.}
    \textbf{Grammaire :} \all{sollen} (modal) Konjunktiv II \all{solltest} (conseil) + Infinitif \all{essen/trinken} à la fin. Subordonnée \all{weil} $\rightarrow$ verbe \all{ist} à la fin. \all{das} (sujet neutre) = « cela / le fait de... ». \all{für} + Akk \all{die Gesundheit}.
    \item \all{\textbf{Wenn} du jeden Tag Sport \textbf{machst}, \textbf{fühlst} du dich besser.}
    \textbf{Grammaire :} \all{wenn} + verbe \all{machst} à la fin (subordonnée). Principale : verbe \all{fühlst} en 2e pos. \all{sich fühlen} pronominal $\rightarrow$ \all{fühlst du dich} (Akk). \all{besser} comparatif.
    \item \all{\textbf{Rauche} nicht und \textbf{trink} keinen Alkohol!}
    \textbf{Grammaire :} Impératif \textit{du}. \all{rauchen} (reg.) $\rightarrow$ \all{rauch} + \textit{e} (euphonie/usage courant) $\rightarrow$ \all{rauche}. \all{trinken} (fort e $\rightarrow$ i) $\rightarrow$ \all{trink} (pas \all{trinke} !). \all{keinen Alkohol} (Akk masc sg \all{kein}).
    \item \all{Schreib mir bald! \textbf{Dein} Marc}
    \textbf{Grammaire :} \all{schreiben} (fort ei $\rightarrow$ ie) $\rightarrow$ \all{schreib} (imp. du). \all{mir} (Dat). \all{Dein} (possessif nominatif masc, car \all{Marc} sujet implicite de la signature).
\end{enumerate}
\textbf{Texte final allemand :}
\begin{textebox}[Lettre à Paul]
\all{Lieber Paul,

du arbeitest zu viel und schläfst zu wenig. Du solltest mehr Gemüse essen und Wasser trinken, weil das gut für die Gesundheit ist. Wenn du jeden Tag Sport machst, fühlst du dich besser. Rauche nicht und trink keinen Alkohol! Schreib mir bald!

Dein Marc}
\end{textebox}
\textbf{Points de vigilance (Bac) :} \all{schläfst} (Umlaut), \all{fühlst du dich} (ordre + pronom), \all{trink} (impératif du fort), \all{keinen} (Akk), \all{weil ... ist} (verbe final), \all{wenn ..., fühlst} (structure V2).
\end{exoresolu}

\begin{methode}[Relire et corriger sa production : La check-list « Zéro Faute »]
\textbf{Objectif :} S'auto-évaluer avant de rendre une copie (Bac ou contrôle) pour gagner les points de « Langue correcte ».
\begin{enumerate}
    \item \textbf{Majuscules (Großschreibung) :} \cle{Tous les noms} ont une majuscule. \cle{Premier mot} de la phrase. \cle{Pronom formel \all{Sie} / \all{Ihnen} / \all{Ihr}}. \textbf{Pas de majuscule} sur \all{ich}, \all{du}, \all{ihr} (sauf début phrase), ni sur les adjectifs (\all{deutsch}, \all{französisch}).
    \item \textbf{Ordre des mots (Wortstellung) :} 
    \begin{itemize}
        \item Principal : \textbf{Verbe conjugué = Position 2}. (Sujet ou Complément en Pos 1).
        \item Subordonnée (\all{weil, dass, wenn, obwohl}) : \textbf{Verbe conjugué = Position finale}.
        \item Séparables : Particule en fin (principal) / collée (subordonnée).
    \end{itemize}
    \item \textbf{Cas (Kasus) \& Articles :} \cle{Nom} (Sujet), \cle{Akk} (Objet direct / après \all{durch, für, gegen, ohne, um}), \cle{Dat} (Objet indirect / après \all{aus, bei, mit, nach, seit, von, zu}). Vérifier \all{der/den/dem}, \all{die/die/der}, \all{das/das/dem}, \all{ein/einen/einem}.
    \item \textbf{Orthographe spécifique :} \cle{Umlaute} (ä, ö, ü) changent le sens (\all{schon} vs \all{schön}, \all{waschen} vs \all{wäscht}). \cle{ß (Eszett)} après voyelle longue/diphtongue (\all{Straße, groß, heiß}) vs \cle{ss} après voyelle courte (\all{dass, Wasser, muss}). \cle{Guillemets} „ “.
    \item \textbf{Faux amis \& Calques :} \all{bekommen} = recevoir (pas devenir). \all{aktuell} = actuel. \all{eventuell} = éventuel. \all{konsequent} = cohérent. \all{Das Gift} = poison. \all{Das Rezept} = ordonnance/recette. \all{Ich bin heiß} = j'ai chaud (sexuel) vs \all{Mir ist heiß} = j'ai chaud (température). \all{Ich langweile mich} = je m'ennuie (pas \all{Ich bin langweilig} = je suis ennuyeux).
\end{enumerate}
\end{methode}

\begin{exoresolu}[Correction d'un texte élève : « Mon week-end sain »]
\textbf{Énoncé :} Le texte ci-dessous contient \textbf{10 erreurs} majeures (grammaire, orthographe, cas, ordre des mots, vocabulaire). Identifiez-les, corrigez-les et expliquez la règle.
\begin{textebox}[Version élève à corriger]
\all{Am wochenende ich habe gut erholt. Ich bin in den park gegangen und habe viel sport getrieben. Das wetter war schön, weil die sonne geschienen hat. Ich habe gesund gegessen: ein salat und ein apfel. Abends ich habe nicht alkohol getrunken, sondern wasser. Das war gut für meine gesundheit. Ich hoffe, das du auch ein schönes wochenende hast.}
\end{textebox}
\tcblower
\textbf{Solution détaillée : Correction annotée}
\begin{center}
\begin{tabularx}{\linewidth}{|X|X|X|}
\hline
\rowcolor{popblueL} \textbf{Phrase élève (erreur soulignée)} & \textbf{Correction} & \textbf{Explication / Règle} \\ \hline
\all{Am \textbf{wochenende} ich habe gut \textbf{erholt}.} & \all{Am \textbf{Wochenende} \textbf{habe} ich mich gut \textbf{erholt}.} & 1. \textbf{Majuscule} : \all{Wochenende} (Nom). 2. \textbf{Ordre V2} : Complément \all{Am Wochenende} en Pos 1 $\rightarrow$ Verbe \all{habe} en Pos 2. 3. \textbf{Pronominal} : \all{sich erholen} $\rightarrow$ \all{habe mich erholt} (Akk). 4. \textbf{Participe} : \all{erholt} (régulier, correct). \\ \hline
\all{Ich bin in \textbf{den park} gegangen} & \all{Ich bin in \textbf{den Park} gegangen} & \textbf{Majuscule} : \all{Park} (Nom). \all{in} + Akk (direction) $\rightarrow$ \all{den} correct. \\ \hline
\all{und habe viel \textbf{sport getrieben}.} & \all{und \textbf{habe viel Sport getrieben}.} & \textbf{Majuscule} : \all{Sport} (Nom). \textbf{Ordre} : Objet \all{viel Sport} avant Participe \all{getrieben}. \\ \hline
\all{\textbf{Das wetter} war schön,} & \all{\textbf{Das Wetter} war schön,} & \textbf{Majuscule} : \all{Wetter}. \\ \hline
\all{\textbf{weil die sonne geschienen hat}.} & \all{\textbf{weil die Sonne} \textbf{geschienen hat}.} & \textbf{Majuscule} : \all{Sonne}. \textbf{Ordre subordonnée} : Participe \all{geschienen} puis auxiliaire \all{hat} à la \textbf{fin}. \\ \hline
\all{Ich habe gesund gegessen: \textbf{ein salat} und \textbf{ein apfel}.} & \all{Ich habe gesund gegessen: \textbf{einen Salat} und \textbf{einen Apfel}.} & \textbf{Majuscules} : \all{Salat, Apfel}. \textbf{Cas Accusatif} : \all{einen} (masc sg) pour les deux. \\ \hline
\all{\textbf{Abends ich habe nicht alkohol getrunken},} & \all{\textbf{Abends habe ich keinen Alkohol getrunken},} & \textbf{Ordre V2} : \all{Abends} (Pos 1) $\rightarrow$ \all{habe} (Pos 2) $\rightarrow$ \all{ich} (Pos 3). \textbf{Négation} : \all{keinen} (Akk masc) au lieu de \all{nicht alkohol}. \textbf{Majuscule} : \all{Alkohol}. \\ \hline
\all{sondern \textbf{wasser}.} & \all{sondern \textbf{Wasser}.} & \textbf{Majuscule} : \all{Wasser}. \\ \hline
\all{Das war gut für \textbf{meine gesundheit}.} & \all{Das war gut für \textbf{meine Gesundheit}.} & \textbf{Majuscule} : \all{Gesundheit}. \all{für} + Akk $\rightarrow$ \all{meine} (fém sg) correct. \\ \hline
\all{Ich hoffe, \textbf{das} du auch \textbf{ein schönes wochenende hast}.} & \all{Ich hoffe, \textbf{dass} du auch \textbf{ein schönes Wochenende hast}.} & \textbf{Conjonction} : \all{dass} (double s) pas \all{das} (pronom/démonstratif). \textbf{Majuscule} : \all{Wochenende}. \textbf{Ordre subordonnée} : \all{hast} à la fin (correct). \\ \hline
\end{tabularx}
\end{center}
\textbf{Texte corrigé final :}
\begin{textebox}[Version corrigée]
\all{Am Wochenende habe ich mich gut erholt. Ich bin in den Park gegangen und habe viel Sport getrieben. Das Wetter war schön, weil die Sonne geschienen hat. Ich habe gesund gegessen: einen Salat und einen Apfel. Abends habe ich keinen Alkohol getrunken, sondern Wasser. Das war gut für meine Gesundheit. Ich hoffe, dass du auch ein schönes Wochenende hast.}
\end{textebox}
\textbf{Bilan :} 10 erreurs corrigées (5 majuscules, 2 ordres des mots V2/subordonnée, 1 pronominal manquant, 1 cas/néation, 1 conjonction \all{dass}).
\end{exoresolu}

\begin{attention}[Les 5 erreurs qui coûtent des points au Bac]
\begin{enumerate}
    \item \textbf{Majuscules oubliées sur les noms :} \all{der apfel, die gesundheit, das wasser, im park, am wochenende} $\rightarrow$ \textbf{Zéro tolérance.} Tout nom = Majuscule. (\all{Apfel, Gesundheit, Wasser, Park, Wochenende}).
    \item \textbf{Ordre des mots dans la subordonnée (\all{weil, dass, wenn}) :} \all{*weil ich bin müde} / \all{*dass er kommt nicht} $\rightarrow$ \textbf{Verbe TOUJOURS à la fin.} (\all{weil ich müde \textbf{bin}}, \all{dass er nicht \textbf{kommt}}).
    \item \textbf{Verbes à particule séparable oubliée en fin de phrase :} \all{*Ich aufstehe um 7 Uhr} / \all{*Er raucht nicht auf} $\rightarrow$ \all{Ich stehe um 7 Uhr \textbf{auf}} / \all{Er hört \textbf{auf} zu rauchen}.
    \item \textbf{Faux amis \& Calques lexicaux :} \all{*Ich bekomme ein Geschenk} (pour « je deviens ») $\rightarrow$ \all{Ich \textbf{werde}...}. \all{*Das ist aktuell} (pour « actuel » adjectif) $\rightarrow$ \all{Das ist \textbf{aktuell}} (ok) mais \all{*eventuell} pour « éventuellement » (ok) attention \all{konsequent} $\neq$ conséquent (conséquent = \all{logisch/folgerichtig}). \all{*Ein Gift} pour « un cadeau » $\rightarrow$ \all{Ein \textbf{Geschenk}}.
    \item \textbf{Cas (Akk/Dat) après prépositions / Verbes :} \all{*Ich helfe \textbf{meinen} Freund} (Akk) $\rightarrow$ \all{Ich helfe \textbf{meinem} Freund} (Dat). \all{*Ich warte auf \textbf{meine} Mutter} (Dat) $\rightarrow$ \all{Ich warte auf \textbf{meine} Mutter} (Akk). \all{*Mit \textbf{mein} Vater} $\rightarrow$ \all{Mit \textbf{meinem} Vater}.
\end{enumerate}
\textbf{Astuce :} Appliquez la \cle{Méthode 6 (Check-list)} systématiquement en fin d'épreuve. 5 minutes de relecture = 5 à 10 points sauvés.
\end{attention}

\newpage

\coursec{Exercices}

\courssub{Je vérifie mes connaissances}

\begin{exercice}[QCM : environnement et santé]
Pour chaque question, entoure la bonne réponse.

1. \all{der Müll} signifie : a) la maladie \quad b) les déchets \quad c) le médecin

2. \all{die Gesundheit} signifie : a) la santé \quad b) l'hygiène \quad c) la nature

3. \all{die Verschmutzung} signifie : a) la protection \quad b) la pollution \quad c) le recyclage

4. Le pluriel de \all{der Baum} est : a) \all{die Baums} \quad b) \all{die Bäume} \quad c) \all{die Baumen}

5. \all{sich erholen} signifie : a) se reposer, récupérer \quad b) se répéter \quad c) se révolter

6. Le contraire de \all{gesund} est : a) \all{müde} \quad b) \all{krank} \quad c) \all{stark}
\end{exercice}

\begin{exercice}[Vrai ou faux ? Justifie ta réponse]
Indique si les affirmations sont vraies (V) ou fausses (F) et corrige les fausses.

1. À l'impératif 2\textsuperscript{e} personne du singulier, le verbe \all{nehmen} devient \all{Nimm!}

2. \all{die Medizin} est un faux ami : il signifie « le médecin ».

3. Dans une subordonnée introduite par \all{weil}, le verbe conjugué va à la fin.

4. Le verbe \all{sich waschen} se conjugue : \all{Ich wasche mir.}

5. Tous les noms allemands prennent une majuscule, même au milieu de la phrase.

6. \all{Recyceln Sie den Müll!} est un impératif à la forme de politesse.
\end{exercice}

\begin{exercice}[Vocabulaire : complète les mots]
Complète chaque mot allemand avec l'article et la voyelle manquante, puis traduis-le en français.

1. d\_ \_ Umw\_lt (l'environnement)

2. d\_ \_ Kr\_nkheit (la maladie)

3. d\_s W\_sser (l'eau)

4. d\_ \_ Luft (l'air)

5. d\_r \_rzt (le médecin)

6. d\_ \_ Bew\_gung (le mouvement, l'exercice physique)
\end{exercice}

\begin{exercice}[Conjugaison : l'impératif à compléter]
Complète les phrases avec l'impératif correct du verbe entre parenthèses (forme \all{du}, \all{ihr} ou \all{Sie} selon le contexte).

1. \all{\ldots\ mehr Wasser!} (trinken, du)

2. \all{\ldots\ den Müll in den Eimer!} (werfen, ihr)

3. \all{\ldots\ Sie bitte das Fenster!} (öffnen, Sie)

4. \all{\ldots\ dich nicht!} (sich erkälten, du)

5. \all{\ldots\ früh ins Bett!} (gehen, ihr)

6. \all{\ldots\ Sie die Umwelt!} (schützen, Sie)
\end{exercice}

\courssub{Je m'entraîne}

\begin{exercice}[Thème : phrases à traduire]
Traduis les phrases suivantes en allemand. Fais attention à l'impératif, aux pronominaux et à la place du verbe.

1. Repose-toi bien, parce que tu es malade.

2. Buvez beaucoup d'eau et mangez des fruits !

3. Je me lave les mains avant le repas.

4. Protégez l'environnement, parce que la nature est importante !

5. Nous allons chez le médecin, car mon frère est malade.
\end{exercice}

\begin{exercice}[Transformation : de l'infinitif à l'impératif]
Transforme chaque conseil en impératif à la forme indiquée.

1. \all{weniger Fleisch essen} (du) : \all{\ldots}

2. \all{jeden Tag Sport machen} (ihr) : \all{\ldots}

3. \all{das Handy nachts ausschalten} (Sie) : \all{\ldots}

4. \all{sich gesund ernähren} (du) : \all{\ldots}

5. \all{die Flaschen recyceln} (ihr) : \all{\ldots}

6. \all{früh aufstehen} (du) : \all{\ldots}
\end{exercice}

\begin{exercice}[Texte à trous : la journée de Paul]
Complète le texte avec les mots de la liste suivante (chaque mot ne sert qu'une fois) : \all{Umwelt, krank, weil, sich, Arzt, Wasser, gesund, Müll}.

\all{Paul fühlt \ldots\ nicht gut. Er ist \ldots\ und muss zum \ldots\ gehen. Der Arzt sagt: „Trinken Sie viel \ldots\ und bleiben Sie im Bett, \ldots\ Sie Fieber haben!“ Am nächsten Tag geht es Paul besser. Er will \ldots\ ernähren und die \ldots\ schützen: er trennt jetzt den \ldots\ zu Hause.}
\end{exercice}

\begin{exercice}[Appariements : mots et définitions]
Relie chaque mot allemand (1--6) à sa définition française (a--f).

\begin{center}
\begin{tabularx}{\linewidth}{|X|X|}
\hline
\rowcolor{popblueL} \textbf{Deutsch} & \textbf{Français} \\
\hline
1. \all{die Ernährung} & a) le fait de jeter les déchets par catégorie \\
\hline
2. \all{die Mülltrennung} & b) l'alimentation \\
\hline
3. \all{das Fieber} & c) un papier distribué pour informer ou sensibiliser \\
\hline
4. \all{das Flugblatt} & d) la température élevée du corps malade \\
\hline
5. \all{der Spaziergang} & e) le fait de réutiliser les matériaux \\
\hline
6. \all{das Recycling} & f) une promenade \\
\hline
\end{tabularx}
\end{center}
\end{exercice}

\courssub{Je me prépare au Bac}

\begin{exercice}[Compréhension de texte et questions de langue]
Lis le texte suivant, puis réponds aux questions en allemand par des phrases complètes.

\begin{textebox}[Gesund leben in Yaoundé]
In Yaoundé gibt es viele Probleme mit der Umwelt: Auf den Straßen liegt oft Müll, und die Luft ist manchmal schlecht, weil es viele Autos und Motorräder gibt. Die Schüler des Gymnasiums wollen das ändern. Jeden Freitag sammeln sie den Müll auf dem Schulhof und trennen Plastik, Papier und Glas. „Wir müssen unsere Stadt schützen, weil sie unser Zuhause ist“, sagt Amina, eine Schülerin der Klasse Première. Die Schüler essen auch gesünder: Sie kaufen Obst auf dem Markt und trinken viel Wasser. Der Direktor ist stolz auf seine Schüler und sagt: „Bewegt euch mehr, schlaft genug und bleibt gesund!“
\end{textebox}

\textbf{A. Compréhension} (réponds en allemand)

1. Quels sont les deux problèmes d'environnement mentionnés dans le texte ?

2. Que font les élèves chaque vendredi ?

3. Pourquoi Amina veut-elle protéger la ville ?

4. Que font les élèves pour manger plus sainement ?

5. Quels sont les trois conseils du directeur ?

\textbf{B. Questions de langue}

6. Releve dans le texte deux impératifs et précise la forme utilisée (\all{du}, \all{ihr} ou \all{Sie}).

7. Releve deux subordonnées avec \all{weil} et explique la place du verbe.

8. Trouve dans le texte un verbe pronominal et conjugue-le au présent à toutes les personnes.
\end{exercice}

\begin{exercice}[Thème et version type Bac]
\textbf{A. Thème.} Traduis les cinq phrases suivantes en allemand.

1. Lève-toi tôt et fais du sport tous les jours !

2. Elle se repose parce qu'elle est fatiguée.

3. Ne jetez pas les déchets dans la rue !

4. Nous devons protéger l'eau, car elle est précieuse.

5. Lavez-vous les mains avant de manger !

\textbf{B. Version.} Traduis le texte suivant en français.

\begin{textebox}[Mülltrennung an einer Schule in Deutschland]
An der Schule in Freiburg gibt es drei verschiedene Eimer: einen für Papier, einen für Plastik und einen für Restmüll. Die Schüler lernen schon in der fünften Klasse, wie man den Müll trennt. „Das ist wichtig, weil wir die Natur schützen müssen“, erklärt die Lehrerin Frau Berger. Jede Woche bringt eine Klasse das Papier zum Container. Die Schule bekommt dafür ein bisschen Geld und kauft neue Bücher für die Bibliothek. Die Schüler sind stolz, weil sie der Umwelt helfen.
\end{textebox}
\end{exercice}

\begin{exercice}[Production écrite type Bac]
Choisis \textbf{un} des deux sujets suivants.

\textbf{Sujet 1 : Lettre à un ami.} Ton ami Thomas est souvent malade. Écris-lui une lettre en allemand : donne-lui \textbf{trois conseils} pour rester en bonne santé (alimentation, sommeil, sport). Utilise au moins \textbf{deux impératifs}, \textbf{un verbe pronominal} et \textbf{une subordonnée avec} \all{weil}. (80 à 100 mots)

Modèle de début : \all{Lieber Thomas, wie geht es dir? Ich habe gehört, dass du oft krank bist. Ich gebe dir drei Ratschläge \ldots}

\textbf{Sujet 2 : Tract.} Avec ta classe, tu participes à une campagne de sensibilisation. Rédige en allemand un tract (\all{das Flugblatt}) contre la pollution dans ta ville. Ton tract doit contenir : un \textbf{titre accrocheur}, au moins \textbf{quatre impératifs}, \textbf{deux subordonnées} (avec \all{weil} ou \all{dass}) et une \textbf{phrase de conclusion} motivante. (60 à 80 mots)

Modèle de titre : \all{„Saubere Stadt -- gesunde Menschen!“}
\end{exercice}

\begin{methode}[Réussir la production écrite]
\begin{enumerate}
\item \textbf{Analyse le sujet} : type de texte (lettre ou tract), destinataire (\all{du} pour un ami, \all{Sie} ou \all{ihr} pour le public), contraintes obligatoires (impératifs, pronominal, subordonnée).
\item \textbf{Planifie} : pour la lettre, trois conseils = trois paragraphes courts ; pour le tract, titre + 4 conseils + conclusion.
\item \textbf{Rédige} avec des phrases simples et sûres : sujet -- verbe en 2\textsuperscript{e} position -- compléments.
\item \textbf{Vérifie les contraintes} : souligne tes impératifs, ton pronominal, ta subordonnée avec \all{weil} (verbe à la fin !).
\item \textbf{Relis} : majuscules des noms, Umlaute, place du verbe, cas après les prépositions (\all{zum Arzt}, \all{vor dem Essen}).
\end{enumerate}
\end{methode}

\begin{exemplebox}[Corrigé possible du sujet 1 (lettre)]
\all{Lieber Thomas,}

\all{wie geht es dir? Ich habe gehört, dass du oft krank bist. Ich gebe dir drei Ratschläge. Erstens: Iss mehr Obst und Gemüse und trink viel Wasser, weil eine gesunde Ernährung wichtig ist. Zweitens: Schlaf genug, mindestens acht Stunden pro Nacht. Drittens: Mach jeden Tag Sport oder geh zu Fuß zur Schule. Beweg dich mehr, dann fühlst du dich besser! Ich hoffe, dass du bald gesund bist.}

\all{Viele Grüße}

\all{dein Paul}

Traduction : « Cher Thomas, comment vas-tu ? J'ai entendu que tu es souvent malade. Je te donne trois conseils. Premièrement : mange plus de fruits et de légumes et bois beaucoup d'eau, parce qu'une alimentation saine est importante. Deuxièmement : dors assez, au moins huit heures par nuit. Troisièmement : fais du sport tous les jours ou va à l'école à pied. Bouge plus, alors tu te sentiras mieux ! J'espère que tu seras bientôt en bonne santé. Amitiés, ton Paul. »
\end{exemplebox}

\begin{exemplebox}[Corrigé possible du sujet 2 (tract)]
\all{„Saubere Stadt -- gesunde Menschen!“}

\all{Unsere Stadt ist schmutzig, weil viele Leute den Müll auf die Straße werfen. Das müssen wir ändern! Werft den Müll in den Eimer! Trennt Plastik, Papier und Glas! Recycelt eure Flaschen! Spart Wasser, weil es kostbar ist! Wir wissen, dass eine saubere Umwelt uns gesund hält. Helft mit -- gemeinsam sind wir stark!}

Traduction : « Ville propre -- gens en bonne santé ! Notre ville est sale parce que beaucoup de gens jettent les déchets dans la rue. Nous devons changer cela ! Jetez les déchets à la poubelle ! Triez le plastique, le papier et le verre ! Recyclez vos bouteilles ! Économisez l'eau, parce qu'elle est précieuse ! Nous savons qu'un environnement propre nous garde en bonne santé. Participez -- ensemble, nous sommes forts ! »
\end{exemplebox}

\begin{exercice}[Reprise d'erreurs : corrige les fautes]
Chaque phrase contient une faute typique des francophones (place du verbe, majuscule, faux ami, impératif incorrect, cas). Corrige-la et explique la règle.

1. \all{Ich gehe in die schule um acht Uhr.}

2. \all{Weil ich krank bin, ich gehe zum Arzt.}

3. \all{Nehmt dein Medizin!} (pour dire « Prends tes médicaments ! »)

4. \all{Trinkst du mehr Wasser!} (conseil à un ami)

5. \all{Ich wasche mir die Hande.}

6. \all{Sie muss das Umwelt schützen.}

7. \all{Gestern ich habe Sport gemacht.}

8. \all{Er besucht die Arzt, weil er Fieber hat.}
\end{exercice}

\begin{attention}[Les 5 pièges de ce chapitre]
\begin{itemize}
\item \textbf{Faux amis} : \all{die Medizin} = le médicament (le médecin = \all{der Arzt}) ; \all{die Luft} = l'air (pas « l'aile ») ; \all{das Gymnasium} = le lycée (pas le gymnase).
\item \textbf{Impératif} : pas de pronom sujet aux formes \all{du} et \all{ihr} ; verbes à changement de voyelle : \all{nehmen} $\rightarrow$ \all{Nimm!}, \all{essen} $\rightarrow$ \all{Iss!}, \all{lesen} $\rightarrow$ \all{Lies!}
\item \textbf{Pronominaux} : \all{mich/dich} à l'accusatif, mais \all{mir/dir} au datif quand il y a un complément : \all{Ich wasche mich.} / \all{Ich wasche mir die Hände.}
\item \textbf{Subordonnées} : après \all{weil, dass, wenn}, le verbe conjugué va à la fin ; si la subordonnée est en tête, inversion dans la principale.
\item \textbf{Majuscules et pluriels} : majuscule à tous les noms ; attention aux pluriels à Umlaut : \all{der Baum} $\rightarrow$ \all{die Bäume}, \all{die Hand} $\rightarrow$ \all{die Hände}, \all{die Stadt} $\rightarrow$ \all{die Städte}.
\end{itemize}
\end{attention}

\newpage

\coursec{Corrigés des exercices}

\courssub{Je vérifie mes connaissances}

\begin{corrige}[Exercice 1 — QCM : environnement et santé]
1. b) les déchets. \quad 2. a) la santé. \quad 3. b) la pollution. \quad 4. b) \all{die Bäume} (pluriel avec Umlaut, sans -s). \quad 5. a) se reposer, récupérer. \quad 6. b) \all{krank} (malade).

Rappel : \all{der Müll} = les déchets ; \all{die Gesundheit} = la santé ; \all{die Verschmutzung} = la pollution ; \all{der Baum} $\rightarrow$ \all{die Bäume} ; \all{sich erholen} = se reposer, récupérer ; \all{gesund} $\leftrightarrow$ \all{krank}.
\end{corrige}

\begin{corrige}[Exercice 2 — Vrai ou faux ?]
1. \textbf{Vrai.} \all{nehmen} est un verbe à changement de voyelle (e $\rightarrow$ i) : \all{Nimm!} (radical seul, sans -e final).

2. \textbf{Faux.} \all{die Medizin} signifie « le médicament, la médecine » ; « le médecin » se dit \all{der Arzt}. C'est bien un faux ami, mais l'affirmation inverse le sens.

3. \textbf{Vrai.} Après \all{weil}, le verbe conjugué se place en dernière position : \all{..., weil ich krank bin.}

4. \textbf{Faux.} On dit \all{Ich wasche mich.} (accusatif). \all{mir} (datif) n'est possible qu'avec un complément de partie du corps : \all{Ich wasche mir die Hände.}

5. \textbf{Vrai.} En allemand, tout nom commun prend une majuscule : \all{die Schule, das Wasser, der Arzt}.

6. \textbf{Vrai.} \all{Recyceln Sie!} est l'impératif de politesse : verbe à l'infinitif + \all{Sie} juste après.
\end{corrige}

\begin{corrige}[Exercice 3 — Vocabulaire : complète les mots]
1. \all{die Umwelt} : l'environnement. \quad 2. \all{die Krankheit} : la maladie. \quad 3. \all{das Wasser} : l'eau. \quad 4. \all{die Luft} : l'air. \quad 5. \all{der Arzt} : le médecin. \quad 6. \all{die Bewegung} : le mouvement, l'exercice physique.

Astuce mémorisation : les noms en \all{-ung} (\all{die Bewegung}), en \all{-heit} (\all{die Krankheit}) et en \all{-tion} (\all{die Information}) sont toujours féminins (\all{die}).
\end{corrige}

\begin{corrige}[Exercice 4 — Conjugaison : l'impératif à compléter]
1. \all{Trink mehr Wasser!} (radical seul, forme \all{du})

2. \all{Werft den Müll in den Eimer!} (radical + -t, forme \all{ihr})

3. \all{Öffnen Sie bitte das Fenster!} (infinitif + \all{Sie}, forme de politesse)

4. \all{Erkälte dich nicht!} (le pronom réfléchi \all{dich} suit immédiatement l'impératif)

5. \all{Geht früh ins Bett!} (forme \all{ihr})

6. \all{Schützen Sie die Umwelt!} (forme de politesse)
\end{corrige}

\courssub{Je m'entraîne}

\begin{corrige}[Exercice 5 — Thème : phrases à traduire]
1. \all{Erhol dich gut, weil du krank bist.} (ou \all{Ruh dich gut aus, weil du krank bist.}) — pronominal à l'impératif : \all{dich} après le verbe ; verbe de la subordonnée \all{weil} à la fin.

2. \all{Trinken Sie viel Wasser und essen Sie Obst!} (ou, à la 2\textsuperscript{e} personne du pluriel : \all{Trinkt viel Wasser und esst Obst!})

3. \all{Ich wasche mir die Hände vor dem Essen.} — datif \all{mir} + article défini, car le possesseur est clair ; \all{vor} + datif : \all{vor dem Essen}.

4. \all{Schützen Sie die Umwelt, weil die Natur wichtig ist!}

5. \all{Wir gehen zum Arzt, denn mein Bruder ist krank.} — \all{denn} = coordination : le verbe reste en 2\textsuperscript{e} position ; variante avec \all{weil} : \all{..., weil mein Bruder krank ist.}
\end{corrige}

\begin{corrige}[Exercice 6 — Transformation : de l'infinitif à l'impératif]
1. \all{Iss weniger Fleisch!} (verbe à changement de voyelle e $\rightarrow$ i : \all{essen} $\rightarrow$ \all{iss})

2. \all{Macht jeden Tag Sport!}

3. \all{Schalten Sie das Handy nachts aus!} (particule séparable \all{aus} renvoyée à la fin)

4. \all{Ernähre dich gesund!} (pronom \all{dich} après l'impératif)

5. \all{Recycelt die Flaschen!}

6. \all{Steh früh auf!} (particule \all{auf} à la fin)
\end{corrige}

\begin{corrige}[Exercice 7 — Texte à trous : la journée de Paul]
\all{Paul fühlt \textbf{sich} nicht gut. Er ist \textbf{krank} und muss zum \textbf{Arzt} gehen. Der Arzt sagt: „Trinken Sie viel \textbf{Wasser} und bleiben Sie im Bett, \textbf{weil} Sie Fieber haben!“ Am nächsten Tag geht es Paul besser. Er will sich \textbf{gesund} ernähren und die \textbf{Umwelt} schützen: er trennt jetzt den \textbf{Müll} zu Hause.}

Remarque : le mot \all{sich} sert deux fois dans le texte (\all{sich fühlen} et \all{sich ernähren}) ; la liste contient donc huit mots pour huit trous. Traduction : « Paul ne se sent pas bien. Il est malade et doit aller chez le médecin. Le médecin dit : „Buvez beaucoup d'eau et restez au lit, parce que vous avez de la fièvre !“ Le lendemain, Paul va mieux. Il veut se nourrir sainement et protéger l'environnement : il trie maintenant les déchets à la maison. »
\end{corrige}

\begin{corrige}[Exercice 8 — Appariements : mots et définitions]
1 -- b ; 2 -- a ; 3 -- d ; 4 -- c ; 5 -- f ; 6 -- e.

\begin{center}
\begin{tabularx}{\linewidth}{|X|X|}
\hline
\rowcolor{popblueL} \textbf{Deutsch} & \textbf{Français} \\
\hline
\all{die Ernährung} & l'alimentation \\
\hline
\all{die Mülltrennung} & le tri des déchets \\
\hline
\all{das Fieber} & la fièvre \\
\hline
\all{das Flugblatt, die Flugblätter} & le tract, le prospectus \\
\hline
\all{der Spaziergang, die Spaziergänge} & la promenade \\
\hline
\all{das Recycling} & le recyclage \\
\hline
\end{tabularx}
\end{center}
\end{corrige}

\courssub{Je me prépare au Bac}

\begin{corrige}[Exercice 9 — Compréhension de texte et questions de langue]
\textbf{A. Compréhension}

1. \all{Auf den Straßen liegt oft Müll, und die Luft ist manchmal schlecht.} (« Il y a souvent des déchets dans les rues et l'air est parfois mauvais. »)

2. \all{Jeden Freitag sammeln sie den Müll auf dem Schulhof und trennen Plastik, Papier und Glas.} (« Chaque vendredi, ils ramassent les déchets dans la cour et trient le plastique, le papier et le verre. »)

3. \all{Sie will die Stadt schützen, weil sie ihr Zuhause ist.} (« Elle veut protéger la ville parce qu'elle est sa maison. »)

4. \all{Sie kaufen Obst auf dem Markt und trinken viel Wasser.} (« Ils achètent des fruits au marché et boivent beaucoup d'eau. »)

5. \all{Er sagt: „Bewegt euch mehr, schlaft genug und bleibt gesund!“} (« Bougez plus, dormez assez et restez en bonne santé ! »)

\textbf{B. Questions de langue}

6. \all{Bewegt euch!}, \all{schlaft!}, \all{bleibt!} : impératifs à la forme \all{ihr} (radical + -t). On trouve aussi dans le texte \all{Trinken Sie!} et \all{bleiben Sie!} : forme de politesse \all{Sie} (infinitif + \all{Sie}).

7. \all{..., weil es viele Autos und Motorräder gibt.} et \all{..., weil sie unser Zuhause ist.} : dans une subordonnée introduite par \all{weil}, le verbe conjugué (\all{gibt}, \all{ist}) se place à la fin de la proposition.

8. \all{sich bewegen} (dans \all{Bewegt euch!}) : \all{ich bewege mich, du bewegst dich, er/sie/es bewegt sich, wir bewegen uns, ihr bewegt euch, sie/Sie bewegen sich.}

Barème indicatif (10 points) : A. 5 $\times$ 1 pt (phrase complète, information exacte) ; B. 6 : 1 pt, 7 : 2 pts, 8 : 2 pts.
\end{corrige}

\begin{corrige}[Exercice 10 — Thème et version type Bac]
\textbf{A. Thème}

1. \all{Steh früh auf und mach jeden Tag Sport!} (ou \all{Stehen Sie früh auf und machen Sie jeden Tag Sport!})

2. \all{Sie ruht sich aus, weil sie müde ist.} (ou \all{Sie erholt sich, weil sie müde ist.})

3. \all{Werfen Sie den Müll nicht auf die Straße!} (ou \all{Werft den Müll nicht auf die Straße!})

4. \all{Wir müssen das Wasser schützen, denn es ist kostbar.} (ou \all{..., weil es kostbar ist.})

5. \all{Waschen Sie sich die Hände vor dem Essen!} (ou \all{Wascht euch die Hände vor dem Essen!})

Points de vigilance : particules séparables à la fin (\all{aufstehen} $\rightarrow$ \all{Steh ... auf!}) ; \all{weil} renvoie le verbe à la fin ; \all{vor} + datif (\all{vor dem Essen}).

\textbf{B. Version}

« Dans l'école de Fribourg, il y a trois poubelles différentes : une pour le papier, une pour le plastique et une pour les déchets résiduels. Les élèves apprennent dès la classe de cinquième comment on trie les déchets. „C'est important, parce que nous devons protéger la nature“, explique la professeure, M\textsuperscript{me} Berger. Chaque semaine, une classe apporte le papier au conteneur. L'école reçoit en échange un peu d'argent et achète de nouveaux livres pour la bibliothèque. Les élèves sont fiers, parce qu'ils aident l'environnement. »

Points de vigilance : \all{der Restmüll} = les déchets résiduels (non recyclables) ; \all{der Umwelt helfen} : \all{helfen} se construit avec le datif ; \all{die fünfte Klasse} = la classe de cinquième (système allemand) ; \all{der Container} = le conteneur.

Barème indicatif (10 points) : thème 5 $\times$ 1 pt (impératif correct, place du verbe, cas) ; version 5 pts (fidélité du sens 3 pts, qualité du français 2 pts).
\end{corrige}

\begin{corrige}[Exercice 11 — Production écrite type Bac]
\textbf{Corrigé possible du sujet 1 (lettre) :}

\all{Lieber Thomas,}

\all{wie geht es dir? Ich habe gehört, dass du oft krank bist. Ich gebe dir drei Ratschläge. Erstens: Iss mehr Obst und Gemüse und trink viel Wasser, weil eine gesunde Ernährung wichtig ist. Zweitens: Schlaf genug, mindestens acht Stunden pro Nacht. Drittens: Mach jeden Tag Sport oder geh zu Fuß zur Schule. Beweg dich mehr, dann fühlst du dich besser! Ich hoffe, dass du bald gesund bist.}

\all{Viele Grüße}

\all{dein Paul}

Traduction : « Cher Thomas, comment vas-tu ? J'ai entendu que tu es souvent malade. Je te donne trois conseils. Premièrement : mange plus de fruits et de légumes et bois beaucoup d'eau, parce qu'une alimentation saine est importante. Deuxièmement : dors assez, au moins huit heures par nuit. Troisièmement : fais du sport tous les jours ou va à l'école à pied. Bouge plus, alors tu te sentiras mieux ! J'espère que tu seras bientôt en bonne santé. Amitiés, ton Paul. »

Contraintes respectées : impératifs \all{Iss}, \all{trink}, \all{Schlaf}, \all{Mach}, \all{Beweg} ; pronominal \all{sich bewegen} / \all{sich fühlen} ; subordonnées avec \all{weil} et \all{dass} (verbe à la fin).

\textbf{Corrigé possible du sujet 2 (tract) :}

\all{„Saubere Stadt -- gesunde Menschen!“}

\all{Unsere Stadt ist schmutzig, weil viele Leute den Müll auf die Straße werfen. Das müssen wir ändern! Werft den Müll in den Eimer! Trennt Plastik, Papier und Glas! Recycelt eure Flaschen! Spart Wasser, weil es kostbar ist! Wir wissen, dass eine saubere Umwelt uns gesund hält. Helft mit -- gemeinsam sind wir stark!}

Traduction : « Ville propre -- gens en bonne santé ! Notre ville est sale parce que beaucoup de gens jettent les déchets dans la rue. Nous devons changer cela ! Jetez les déchets à la poubelle ! Triez le plastique, le papier et le verre ! Recyclez vos bouteilles ! Économisez l'eau, parce qu'elle est précieuse ! Nous savons qu'un environnement propre nous garde en bonne santé. Participez -- ensemble, nous sommes forts ! »

Contraintes respectées : titre accrocheur ; quatre impératifs \all{ihr} (\all{Werft, Trennt, Recycelt, Spart}) + \all{Helft mit!} ; subordonnées \all{weil} et \all{dass} ; conclusion motivante.

Barème indicatif (10 points) : respect du type de texte et des contraintes 4 pts ; correction grammaticale (impératifs, place du verbe, cas) 3 pts ; vocabulaire du chapitre 2 pts ; orthographe et majuscules 1 pt.
\end{corrige}

\begin{corrige}[Exercice 12 — Reprise d'erreurs : corrige les fautes]
1. \all{Ich gehe um acht Uhr in die Schule.} Règle : majuscule à tous les noms (\all{die Schule}).

2. \all{Weil ich krank bin, gehe ich zum Arzt.} Règle : quand la subordonnée \all{weil} ouvre la phrase, elle compte comme la 1\textsuperscript{re} position ; le verbe de la principale vient juste après la virgule (pas de double sujet).

3. \all{Nimm deine Medikamente!} Règle : conseil à une personne = forme \all{du} (\all{nimm}, changement de voyelle e $\rightarrow$ i) ; \all{die Medizin} est un faux ami (« la médecine, le médicament »), « les médicaments » = \all{die Medikamente} ; possessif accordé au pluriel : \all{deine}.

4. \all{Trink mehr Wasser!} Règle : l'impératif \all{du} se forme avec le radical seul, sans -st ni pronom : \all{trinken} $\rightarrow$ \all{Trink!}

5. \all{Ich wasche mir die Hände.} Règle : pluriel de \all{die Hand} = \all{die Hände} (Umlaut obligatoire).

6. \all{Sie muss die Umwelt schützen.} Règle : \all{die Umwelt} est féminin ; à l'accusatif, l'article reste \all{die}.

7. \all{Gestern habe ich Sport gemacht.} Règle : le verbe conjugué est toujours en 2\textsuperscript{e} position ; si un complément de temps ouvre la phrase, le sujet passe après le verbe (inversion).

8. \all{Er besucht den Arzt, weil er Fieber hat.} Règle : \all{der Arzt} est masculin ; à l'accusatif : \all{den Arzt}.
\end{corrige}

\newpage

\coursec{Fiche bilan}

\begin{popfigure}
\begin{center}\tcbox[colback=popblue,colframe=popblue,coltext=white,arc=9pt,boxsep=4pt,left=6pt,right=6pt]{\bfseries\small Umwelt und Gesundheit (Bilan)}\end{center}
\vspace{-2pt}
\begin{tcbraster}[raster columns=3,raster equal height=rows,raster column skip=6pt,raster row skip=6pt,enhanced,blankest]
\begin{tcolorbox}[enhanced,colback=white,colframe=popgreen,boxrule=0.9pt,arc=3pt,title={Lexique Umwelt},fonttitle=\bfseries\scriptsize,coltitle=white,colbacktitle=popgreen,left=4pt,right=4pt,top=2pt,bottom=2pt,boxsep=1pt,toptitle=1.5pt,bottomtitle=1.5pt,halign=left]
\begin{itemize}[nosep,leftmargin=*,itemsep=1pt,font=\scriptsize]
\item die Umwelt der Müll
\item recyceln schützen
\end{itemize}
\end{tcolorbox}
\begin{tcolorbox}[enhanced,colback=white,colframe=poporange,boxrule=0.9pt,arc=3pt,title={Lexique Gesundheit},fonttitle=\bfseries\scriptsize,coltitle=white,colbacktitle=poporange,left=4pt,right=4pt,top=2pt,bottom=2pt,boxsep=1pt,toptitle=1.5pt,bottomtitle=1.5pt,halign=left]
\begin{itemize}[nosep,leftmargin=*,itemsep=1pt,font=\scriptsize]
\item die Gesundheit der Arzt
\item gesund leben sich bewegen
\end{itemize}
\end{tcolorbox}
\begin{tcolorbox}[enhanced,colback=white,colframe=poppurple,boxrule=0.9pt,arc=3pt,title={Impératif conseiller},fonttitle=\bfseries\scriptsize,coltitle=white,colbacktitle=poppurple,left=4pt,right=4pt,top=2pt,bottom=2pt,boxsep=1pt,toptitle=1.5pt,bottomtitle=1.5pt,halign=left]
\begin{itemize}[nosep,leftmargin=*,itemsep=1pt,font=\scriptsize]
\item Iss! Geh! Trennt den Müll!
\end{itemize}
\end{tcolorbox}
\begin{tcolorbox}[enhanced,colback=white,colframe=poppink,boxrule=0.9pt,arc=3pt,title={Verbes pronominaux},fonttitle=\bfseries\scriptsize,coltitle=white,colbacktitle=poppink,left=4pt,right=4pt,top=2pt,bottom=2pt,boxsep=1pt,toptitle=1.5pt,bottomtitle=1.5pt,halign=left]
\begin{itemize}[nosep,leftmargin=*,itemsep=1pt,font=\scriptsize]
\item sich freuen sich waschen
\end{itemize}
\end{tcolorbox}
\begin{tcolorbox}[enhanced,colback=white,colframe=popgold,boxrule=0.9pt,arc=3pt,title={Subordonnées verbe final},fonttitle=\bfseries\scriptsize,coltitle=white,colbacktitle=popgold,left=4pt,right=4pt,top=2pt,bottom=2pt,boxsep=1pt,toptitle=1.5pt,bottomtitle=1.5pt,halign=left]
\begin{itemize}[nosep,leftmargin=*,itemsep=1pt,font=\scriptsize]
\item weil, dass wenn, damit
\end{itemize}
\end{tcolorbox}
\begin{tcolorbox}[enhanced,colback=white,colframe=popblue!70,boxrule=0.9pt,arc=3pt,title={Traduction thème / version},fonttitle=\bfseries\scriptsize,coltitle=white,colbacktitle=popblue!70,left=4pt,right=4pt,top=2pt,bottom=2pt,boxsep=1pt,toptitle=1.5pt,bottomtitle=1.5pt,halign=left]
\begin{itemize}[nosep,leftmargin=*,itemsep=1pt,font=\scriptsize]
\item ordre des mots cas et genres
\end{itemize}
\end{tcolorbox}
\begin{tcolorbox}[enhanced,colback=white,colframe=popdark,boxrule=0.9pt,arc=3pt,title={Production tract, lettre},fonttitle=\bfseries\scriptsize,coltitle=white,colbacktitle=popdark,left=4pt,right=4pt,top=2pt,bottom=2pt,boxsep=1pt,toptitle=1.5pt,bottomtitle=1.5pt,halign=left]
\begin{itemize}[nosep,leftmargin=*,itemsep=1pt,font=\scriptsize]
\item conseils lettre à un ami
\end{itemize}
\end{tcolorbox}
\end{tcbraster}

\legende{Carte mentale de la leçon : environnement, santé, grammaire et production.}
\end{popfigure}

\begin{bilan}
\textbf{1. Lexique clé à connaître par cœur}
\begin{itemize}
  \item \all{die Umwelt} (l'environnement), \all{der Umweltschutz} (la protection de l'environnement), \all{die Natur, -en} (la nature), \all{der Müll} (les déchets), \all{der Müll trennen} (trier les déchets), \all{recyceln} (recycler), \all{sparen} (économiser), \all{die Luft} (l'air), \all{das Wasser} (l'eau), \all{der Baum, -äume} (l'arbre), \all{verschmutzen} (polluer).
  \item \all{die Gesundheit} (la santé), \all{gesund / krank} (en bonne santé / malade), \all{der Arzt, -ärzte} (le médecin), \all{die Krankheit, -en} (la maladie), \all{der Sport} (le sport), \all{sich bewegen} (bouger, faire de l'exercice), \all{sich gesund ernähren} (se nourrir sainement), \all{der Schlaf} (le sommeil), \all{aufgeben / aufhören mit} (arrêter de), \all{das Rauchen} (le tabagisme).
\end{itemize}

\textbf{2. Grammaire à connaître par cœur}
\begin{center}
\begin{tabularx}{\linewidth}{|X|X|}
\hline
\rowcolor{popblueL} \textbf{Règle} & \textbf{Exemple} \\
\hline
Impératif (du) : radical du verbe, souvent sans \all{-e} ; verbes forts en e $\rightarrow$ i/ie changent la voyelle & \all{Trink Wasser!} ; \all{Iss mehr Obst!} ; \all{Lies den Text!} \\
\hline
Impératif (ihr / Sie) : forme du présent sans pronom & \all{Trennt den Müll!} ; \all{Schützen Sie die Natur!} \\
\hline
Verbes pronominaux : pronom réfléchi au cas exigé (\all{mich/dich/sich} acc., \all{mir/dir/sich} dat.) & \all{Ich wasche mich.} / \all{Ich wasche mir die Hände.} \\
\hline
Subordonnée : le verbe conjugué va à la fin & \all{Ich mache Sport, weil ich gesund bleiben will.} \\
\hline
\all{damit} + subordonnée = pour que / afin de & \all{Wir recyceln, damit die Umwelt sauber bleibt.} \\
\hline
Verbe à particule séparable : particule en fin de phrase principale & \all{Er hört mit dem Rauchen auf.} \\
\hline
\end{tabularx}
\end{center}

\textbf{3. Méthodes express}
\begin{itemize}
  \item \emph{Tract} : titre accrocheur + impératifs courts + 3 ou 4 conseils + slogan final. Ex. : \all{„Rette unsere Umwelt! Trenne deinen Müll!“}
  \item \emph{Lettre à un ami} : \all{Lieber Paul, / Liebe Anna,} $\rightarrow$ corps avec \all{weil}, \all{dass}, \all{damit} $\rightarrow$ \all{Viele Grüße / Dein(e) ...}
  \item \emph{Thème (fr. $\rightarrow$ all.)} : repérer le temps, placer le verbe conjugué en 2\textsuperscript{e} position, vérifier genre, cas et majuscules.
  \item \emph{Version (all. $\rightarrow$ fr.)} : repérer d'abord le verbe conjugué, puis le sujet ; traduire la subordonnée après avoir trouvé le verbe final.
  \item \emph{Relecture} : verbe en 2\textsuperscript{e} position ? majuscule aux noms ? particule séparable en fin ? pronom réfléchi correct ?
\end{itemize}
\end{bilan}

\begin{aretenir}[Checklist avant le Bac]
\begin{itemize}
  \item $\square$ Je connais le vocabulaire de l'environnement (\all{die Umwelt, der Müll, recyceln, schützen}) avec les articles.
  \item $\square$ Je connais le vocabulaire de la santé (\all{die Gesundheit, der Arzt, sich bewegen, gesund}) avec les articles.
  \item $\square$ Je forme l'impératif aux trois personnes : \all{Komm! Kommt! Kommen Sie!}
  \item $\square$ Je n'oublie pas le changement de voyelle à l'impératif : \all{Iss! Lies! Nimm!}
  \item $\square$ Je place le pronom réfléchi au bon cas : \all{Ich freue mich.} / \all{Ich wasche mir die Hände.}
  \item $\square$ Je mets le verbe conjugué à la fin de la subordonnée : \all{..., weil ich krank bin.}
  \item $\square$ Je sépare la particule en fin de phrase : \all{Er steht früh auf.}
  \item $\square$ Je mets une majuscule à tous les noms allemands.
  \item $\square$ Je structure un tract : titre, impératifs, conseils, slogan.
  \item $\square$ Je rédige une lettre avec formules d'appel et de politesse correctes.
  \item $\square$ En traduction, je vérifie l'ordre des mots, les cas et les temps.
  \item $\square$ Je relis ma production pour corriger orthographe, genres et accords.
\end{itemize}
\end{aretenir}

\findecours

\end{document}
